% Circle Geometry — Pure Mathematics with Mechanics Upper Sixth — Cours signé OnBuch+
% Official MINESEC syllabus. Compiler avec : tectonic PMM01-circle-geometry.tex
\newcommand{\DOCMATIERE}{Pure Mathematics with Mechanics}
\newcommand{\DOCNIVEAU}{Upper Sixth}
\newcommand{\DOCMODULE}{Upper Sixth — Pure mathematics with mechanics}
\newcommand{\DOCLECON}{1}
\newcommand{\DOCTITRE}{Circle Geometry}
% OnBuch+ — préambule « cours signés » ENRICHI (compilé avec tectonic / XeLaTeX)
% Système visuel « Pop » d'OnBuch+ : fond crème (jamais blanc), bordures encre
% épaisses, ombres « dures » décalées, titres Archivo Black en capitales.
% Version MATHÉMATIQUES : graphes (pgfplots/tikz), boîtes pédagogiques
% (définition, propriété, méthode, attention, le savais-tu, exercice résolu,
% exercice, TP, bilan), page de garde et signature OnBuch+ ancrée en bas.
%
% Macros attendues dans main.tex AVANT \input{preamble} :
%   \DOCMATIERE  \DOCNIVEAU  \DOCMODULE  \DOCLECON (numéro)  \DOCTITRE
\documentclass[11pt]{article}

\usepackage{fontspec}
\usepackage[english]{babel}
\usepackage[a4paper,margin=2cm,top=3.4cm,bottom=2.8cm,headheight=1.4cm,footskip=1.3cm]{geometry}
\usepackage{amsmath,amssymb,amsfonts}
\usepackage{mathtools} % \xmapsto, \coloneqq... souvent utilisés par les modèles
\usepackage{cancel} % simplifications barrées (bilans de forces)
% \abs{z} (module, valeur absolue) et \norm{u} (norme), utilisés par les modèles
\providecommand{\abs}{}\let\abs\relax\DeclarePairedDelimiter{\abs}{\lvert}{\rvert}
\providecommand{\norm}{}\let\norm\relax\DeclarePairedDelimiter{\norm}{\lVert}{\rVert}
\usepackage[table]{xcolor} % [table] : fournit aussi \rowcolors (alternance de lignes)
\usepackage{pagecolor}
\usepackage{tikz}
\usetikzlibrary{arrows.meta,positioning,calc,shapes.geometric,shapes.misc,shapes.arrows,decorations.pathmorphing,decorations.pathreplacing,decorations.markings,patterns,shadows,mindmap,trees,fit,backgrounds,matrix,intersections,angles,quotes}
\usepackage{pgfplots}
\usepackage[version=4]{mhchem} % équations nucléaires \ce{^{235}_{92}U}
\usepackage[european,siunitx]{circuitikz} % schémas électriques
\pgfplotsset{compat=1.18}
\usepgfplotslibrary{fillbetween,groupplots} % aires entre courbes (calcul intégral)
\usepackage{enumitem}
\usepackage{fancyhdr}
% [most] charge toutes les bibliothèques tcolorbox (listings, minted, raster,
% documentation...) et gonfle inutilement la mémoire TeX sur un document long
% (~300 boîtes) : on ne charge que ce qui est réellement utilisé.
\usepackage[skins,breakable]{tcolorbox}
\usepackage{graphicx}
\usepackage{colortbl}
\usepackage{booktabs}
\usepackage{tabularx}
\usepackage{diagbox} % case d'en-tête diagonale des tableaux à double entrée
% Colonnes extensibles centrée / gauche / droite (C, L, R), souvent utilisées par les modèles
\newcolumntype{C}{>{\centering\arraybackslash}X}
\newcolumntype{L}{>{\raggedright\arraybackslash}X}
\newcolumntype{R}{>{\raggedleft\arraybackslash}X}
\usepackage{array}
\usepackage{multicol}
\usepackage{siunitx}
% parse-numbers=false : accepte aussi \SI{2,5 \times 10^{-3}}{...}
\sisetup{locale=UK,output-decimal-marker={.},group-minimum-digits=4,parse-numbers=false}
\DeclareSIUnit\ohm{\ensuremath{\symup{\Omega}}} % U+2126 absent de la police
\DeclareSIUnit\division{div} % réglages d'oscilloscope (ms/div, V/div)
\DeclareSIUnit\tr{tr} % tours (vitesse de rotation, ex. \SI{1500}{\tr\per\minute})
\usepackage{qrcode}
% \degree : commande fréquemment utilisée par les modèles pour un angle ou une
% température en dehors de \SI{}{} (non fournie par les paquets chargés).
\providecommand{\degree}{^{\circ}}
% \qed (fin de démonstration), utilisé par les modèles sans amsthm
\providecommand{\qed}{\hfill$\square$}
% \barre : double barre des valeurs interdites dans un tableau de variations (tkz-tab)
\providecommand{\barre}{\ensuremath{\|}}
% environnement proof (amsthm non chargé) utilisé par les modèles
\ifdefined\proof\else\newenvironment{proof}[1][Proof]{\par\noindent\textit{#1.}\ }{\qed\par}\fi
\usepackage{lastpage}
\usepackage{hyperref}
\hypersetup{hidelinks}

% ============================================================
% Palette « Pop » OnBuch+ — mêmes hex que la classe P (plus_ui.dart)
% ============================================================
\definecolor{poporange}{HTML}{F59321}
\definecolor{popdark}{HTML}{E07A0C}
\definecolor{popcream}{HTML}{FFF6E8}
\definecolor{popink}{HTML}{1C1714}
\definecolor{poppurple}{HTML}{7A5AE0}
\definecolor{popgreen}{HTML}{2E9E6A}
\definecolor{poppink}{HTML}{E0457B}
\definecolor{popblue}{HTML}{2F7DE1}
\definecolor{popgold}{HTML}{C98A12}
\definecolor{popmuted}{HTML}{8A7F76}
\definecolor{popline}{HTML}{EADFD2}
\definecolor{popgray}{HTML}{8A7F76}
\colorlet{popgrey}{popgray}
% Alias tolérants (le modèle nomme parfois les couleurs différemment)
\colorlet{popwhite}{white}
\colorlet{popblack}{popink}
\colorlet{popred}{poppink}
\colorlet{popyellow}{popgold}
% Teintes claires pour fonds de tableaux / zones
\colorlet{popblueL}{popblue!10!white}
\colorlet{poporangeL}{poporange!14!white}
\colorlet{popgreenL}{popgreen!12!white}
\colorlet{poppurpleL}{poppurple!12!white}
\colorlet{poppinkL}{poppink!12!white}
\colorlet{popgoldL}{popgold!14!white}

\pagecolor{popcream}

% ============================================================
% Typographie
% ============================================================
\defaultfontfeatures{Ligatures=TeX}
\setmainfont{PlusJakartaSans-Medium.ttf}[Path=../../../fonts/,BoldFont=PlusJakartaSans-Bold.ttf]
% Maths en sans-serif (Fira Math) avec chiffres et lettres droites de Plus
% Jakarta, pour que les formules se fondent dans le texte.
\usepackage[math-style=ISO,bold-style=ISO]{unicode-math}
\setmathfont{FiraMath-Regular.otf}[Path=../../../fonts/]
\setmathfont{PlusJakartaSans-Medium.ttf}[Path=../../../fonts/,range={up/num,up/{Latin,latin},\mathup}]
% (icomma retiré : décimales avec point en anglais)
% Caractères mathématiques Unicode absents des polices de texte (Plus Jakarta,
% Archivo Black) : rendus par leur équivalent mathématique, en texte comme en
% formule — sinon ils s'affichent en carré vide (ex. « Arithmétique dans ℤ »).
\usepackage{newunicodechar}
\newunicodechar{ℝ}{\ensuremath{\mathbb{R}}}
\newunicodechar{ℤ}{\ensuremath{\mathbb{Z}}}
\newunicodechar{ℂ}{\ensuremath{\mathbb{C}}}
\newunicodechar{ℕ}{\ensuremath{\mathbb{N}}}
\newunicodechar{ℚ}{\ensuremath{\mathbb{Q}}}
\newunicodechar{𝔻}{\ensuremath{\mathbb{D}}}
\newunicodechar{α}{\ensuremath{\alpha}}
\newunicodechar{β}{\ensuremath{\beta}}
\newunicodechar{γ}{\ensuremath{\gamma}}
\newunicodechar{δ}{\ensuremath{\delta}}
\newunicodechar{ε}{\ensuremath{\varepsilon}}
\newunicodechar{θ}{\ensuremath{\theta}}
\newunicodechar{λ}{\ensuremath{\lambda}}
\newunicodechar{μ}{\ensuremath{\mu}}
\newunicodechar{σ}{\ensuremath{\sigma}}
\newunicodechar{τ}{\ensuremath{\tau}}
\newunicodechar{φ}{\ensuremath{\varphi}}
\newunicodechar{ψ}{\ensuremath{\psi}}
\newunicodechar{ω}{\ensuremath{\omega}}
\newunicodechar{Σ}{\ensuremath{\Sigma}}
% majuscules produites par \MakeUppercase dans les titres (α-aminés -> Α)
\newunicodechar{Α}{\ensuremath{\alpha}}
\newunicodechar{Β}{\ensuremath{\beta}}
\newunicodechar{Γ}{\ensuremath{\Gamma}}
\newunicodechar{Δ}{\ensuremath{\Delta}}
\newunicodechar{Ε}{\ensuremath{\varepsilon}}
\newunicodechar{Θ}{\ensuremath{\Theta}}
\newunicodechar{Λ}{\ensuremath{\Lambda}}
\newunicodechar{Μ}{\ensuremath{\mu}}
\newunicodechar{Τ}{\ensuremath{\tau}}
\newunicodechar{Φ}{\ensuremath{\Phi}}
\newunicodechar{Ψ}{\ensuremath{\Psi}}
\newunicodechar{Ω}{\ensuremath{\Omega}}
\newunicodechar{↦}{\ensuremath{\mapsto}}
\newunicodechar{∈}{\ensuremath{\in}}
\newunicodechar{∉}{\ensuremath{\notin}}
\newunicodechar{∧}{\ensuremath{\wedge}}
\newunicodechar{⇔}{\ensuremath{\Leftrightarrow}}
\newunicodechar{⟺}{\ensuremath{\Longleftrightarrow}}
\newunicodechar{⇒}{\ensuremath{\Rightarrow}}
\newunicodechar{⟹}{\ensuremath{\Longrightarrow}}
\newunicodechar{∘}{\ensuremath{\circ}}
\newunicodechar{∪}{\ensuremath{\cup}}
\newunicodechar{∩}{\ensuremath{\cap}}
\newunicodechar{≡}{\ensuremath{\equiv}}
\newunicodechar{⊂}{\ensuremath{\subset}}
\newunicodechar{⊥}{\ensuremath{\perp}}
\newunicodechar{∃}{\ensuremath{\exists}}
\newunicodechar{∀}{\ensuremath{\forall}}
\newunicodechar{∛}{\ensuremath{\sqrt[3]{\,}}}
\newunicodechar{∜}{\ensuremath{\sqrt[4]{\,}}}
\newunicodechar{⃗}{\ensuremath{{}^{\to}}}
\newunicodechar{ⁿ}{\ifmmode^{n}\else\textsuperscript{\ensuremath{n}}\fi}
\newunicodechar{ˣ}{\ifmmode^{x}\else\textsuperscript{\ensuremath{x}}\fi}
\newunicodechar{ᵇ}{\ifmmode^{b}\else\textsuperscript{\ensuremath{b}}\fi}
\newunicodechar{ʳ}{\ifmmode^{r}\else\textsuperscript{\ensuremath{r}}\fi}
\newunicodechar{ᵘ}{\ifmmode^{u}\else\textsuperscript{\ensuremath{u}}\fi}
\newunicodechar{ʸ}{\ifmmode^{y}\else\textsuperscript{\ensuremath{y}}\fi}
\newunicodechar{ᵃ}{\ifmmode^{a}\else\textsuperscript{\ensuremath{a}}\fi}
\newunicodechar{ᵉ}{\ifmmode^{e}\else\textsuperscript{\ensuremath{e}}\fi}
\newunicodechar{ᵏ}{\ifmmode^{k}\else\textsuperscript{\ensuremath{k}}\fi}
\newunicodechar{ᵖ}{\ifmmode^{p}\else\textsuperscript{\ensuremath{p}}\fi}
\newunicodechar{ᴺ}{\ifmmode^{N}\else\textsuperscript{\ensuremath{N}}\fi}
\newunicodechar{ᶜ}{\ifmmode^{c}\else\textsuperscript{\ensuremath{c}}\fi}
\newunicodechar{ʰ}{\ifmmode^{h}\else\textsuperscript{\ensuremath{h}}\fi}
\newunicodechar{ᵠ}{\ifmmode^{q}\else\textsuperscript{\ensuremath{q}}\fi}
\newunicodechar{ₙ}{\ifmmode_{n}\else\textsubscript{\ensuremath{n}}\fi}
\newunicodechar{ₐ}{\ifmmode_{a}\else\textsubscript{\ensuremath{a}}\fi}
\newunicodechar{ᵢ}{\ifmmode_{i}\else\textsubscript{\ensuremath{i}}\fi}
\newunicodechar{ᵣ}{\ifmmode_{r}\else\textsubscript{\ensuremath{r}}\fi}
\newunicodechar{ₖ}{\ifmmode_{k}\else\textsubscript{\ensuremath{k}}\fi}
\newunicodechar{ᵦ}{\ifmmode_{\beta}\else\textsubscript{\ensuremath{\beta}}\fi}
\newunicodechar{ₓ}{\ifmmode_{x}\else\textsubscript{\ensuremath{x}}\fi}
\newfontfamily\popdisplay{ArchivoBlack-Regular.ttf}[Path=../../../fonts/]
\newfontfamily\poplogo{SpaceGrotesk-Bold.ttf}[Path=../../../fonts/]
\newfontfamily\popsemi{PlusJakartaSans-SemiBold.ttf}[Path=../../../fonts/]
\newfontfamily\popxbold{PlusJakartaSans-ExtraBold.ttf}[Path=../../../fonts/]
\newfontfamily\popxboldls{PlusJakartaSans-ExtraBold.ttf}[Path=../../../fonts/,LetterSpace=3.0]

\setlist[enumerate]{leftmargin=1.7em,itemsep=5pt,topsep=4pt}
\setlist[itemize]{leftmargin=1.7em,itemsep=4pt,topsep=4pt,label={\color{poporange}\textbullet}}
\setlength{\parindent}{0pt}
\setlength{\parskip}{5pt}
\renewcommand{\arraystretch}{1.25}

% pgfplots : style Pop par défaut
\pgfplotsset{
  every axis/.append style={axis line style={popink,line width=1pt},tick style={popink},
    grid style={popline,line width=0.5pt},label style={font=\small},tick label style={font=\footnotesize}},
  popcurve/.style={poporange,line width=1.8pt,no markers,smooth},
}
% Même style disponible dans un \draw TikZ simple (hors pgfplots)
\tikzset{popcurve/.style={poporange,line width=1.8pt}}
\tikzset{popgrid/.style={popline,very thin}}
\colorlet{popcurve}{poporange} % le nom du style est parfois utilisé comme couleur

% ============================================================
% Pill / squiggle
% ============================================================
\newcommand{\poppill}[3]{%
  \tikz[baseline=-0.55ex]\node[draw=popink,line width=1.2pt,rounded corners=2.6pt,%
    fill=#2,inner xsep=6.5pt,inner ysep=3pt]{%
    \popxboldls\fontsize{7}{7}\selectfont\color{#3}\MakeUppercase{#1}};%
}
\newcommand{\popsquiggle}[1]{%
  \tikz[baseline=-0.4ex]{
    \draw[#1,line width=2.6pt,line cap=round]
      plot [smooth] coordinates {(0,0.09) (0.34,0.01) (0.68,0.21) (1.02,0.01) (1.36,0.10)};
  }%
}
% Mot-clé surligné (marqueur orange sous le texte)
\newcommand{\cle}[1]{{\popxbold\color{popdark}#1}}

% ============================================================
% En-tête / pied
% ============================================================
\pagestyle{fancy}
\fancyhf{}
\renewcommand{\headrulewidth}{0pt}
\renewcommand{\footrulewidth}{0pt}
\lhead{\raisebox{-0.15ex}{\poplogo\fontsize{13}{13}\selectfont\color{popink}OnBuch\color{poporange}+}}
\rhead{\poppill{\DOCMATIERE\ \textbullet\ \DOCNIVEAU\ \textbullet\ Chapter \DOCLECON}{white}{popink}}
\lfoot{\popxbold\fontsize{7.6}{7.6}\selectfont\color{popmuted}\MakeUppercase{Signed course OnBuch+}\ \textbullet\ {\popsemi\DOCTITRE}}
\rfoot{\tikz[baseline=-0.6ex]\node[rounded corners=8.5pt,fill=popink,minimum height=17pt,minimum width=17pt,inner xsep=5pt,inner sep=0]{\popxbold\fontsize{8}{8}\selectfont\color{popcream}\thepage\,/\,\pageref*{LastPage}};}

% ============================================================
% Titres
% ============================================================
\newcommand{\chaptitre}[2]{%
  \begin{tcolorbox}[enhanced,colback=poporange,colframe=popink,boxrule=2.6pt,arc=11pt,
      drop shadow=popink,left=15pt,right=15pt,top=12pt,bottom=11pt,before skip=3pt,after skip=18pt]
    \poppill{#2}{popink}{popcream}\par\vspace{9pt}
    {\raggedright\hyphenpenalty=10000\exhyphenpenalty=10000\popdisplay\fontsize{22}{24}\selectfont\color{popcream}\MakeUppercase{#1}\par}
  \end{tcolorbox}
}
% Section numérotée : pastille orange avec le numéro + titre Archivo
\newcounter{popsec}
\newcommand{\coursec}[1]{%
  \stepcounter{popsec}\setcounter{popsub}{0}%
  \par\vspace{16pt}\noindent
  \begin{minipage}{\linewidth}
  \tikz[baseline=-0.7ex]\node[draw=popink,line width=1.6pt,fill=poporange,rounded corners=4pt,minimum size=22pt,inner sep=0]{\popdisplay\fontsize{12}{12}\selectfont\color{popcream}\thepopsec};%
  \hspace{8pt}{\popdisplay\fontsize{15}{17}\selectfont\color{popink}\MakeUppercase{#1}}\par
  \vspace{4pt}\noindent{\color{poporange}\rule{36pt}{3.6pt}}
  \end{minipage}\par\nopagebreak\vspace{8pt}
}
\newcounter{popsub}
\newcommand{\courssub}[1]{%
  \stepcounter{popsub}%
  \par\vspace{9pt}\noindent{\popxbold\fontsize{11.5}{13}\selectfont\color{popink}{\color{poporange}\thepopsec.\thepopsub}\hspace{6pt}#1}\par\nopagebreak\vspace{4pt}
}

% ============================================================
% Boîtes pédagogiques — même recette : fond blanc, bordure épaisse colorée,
% ombre dure encre, badge plein. Titre optionnel [#1].
% ============================================================
\tcbset{popbase/.style={breakable,enhanced,colback=white,boxrule=2.3pt,arc=9pt,
  shadow={3pt}{-3pt}{0pt}{popink},left=14pt,right=14pt,top=11pt,bottom=11pt,before skip=11pt,after skip=15pt,
  fontupper=\popsemi\fontsize{10.6}{14.5}\selectfont\color{popink},
  fontlower=\popsemi\fontsize{10.6}{14.5}\selectfont\color{popink}}}
% Coche dessinée (le glyphe ✓ n'existe pas dans les polices du document)
\newcommand{\popcheck}[1]{\tikz[baseline=-0.5ex]\draw[#1,line width=1.6pt,line cap=round,line join=round](-0.13,0.0)--(-0.04,-0.09)--(0.13,0.1);}
\providecommand{\coche}{\popcheck{popgreen}}
\providecommand{\resultat}[1]{\par\smallskip\noindent\fcolorbox{popgreen}{popgreenL}{\parbox{\dimexpr\linewidth-2\fboxsep-2\fboxrule\relax}{\textbf{Result.} #1}}\par\smallskip}
\newcommand{\popboxhead}[3]{\poppill{#1}{#2}{white}\ifx\relax#3\relax\else\hspace{6pt}{\popxbold\fontsize{10.6}{12}\selectfont\color{popink}#3}\fi\par\vspace{7pt}}

\newtcolorbox{aretenir}[1][]{popbase,colframe=popgreen,before upper={\popboxhead{Key points}{popgreen}{#1}}}
\newtcolorbox{exemplebox}[1][]{popbase,colframe=poporange,before upper={\popboxhead{Exemple}{poporange}{#1}}}
\newtcolorbox{definition}[1][]{popbase,colframe=popblue,before upper={\popboxhead{Definition}{popblue}{#1}}}
\newtcolorbox{propriete}[1][]{popbase,colframe=poppurple,before upper={\popboxhead{Property}{poppurple}{#1}}}
\newtcolorbox{methode}[1][]{popbase,colframe=popgold,before upper={\popboxhead{Method}{popgold}{#1}}}
\newtcolorbox{attention}[1][]{popbase,colframe=poppink,before upper={\popboxhead{Warning}{poppink}{#1}}}
\newtcolorbox{experience}[1][]{popbase,colframe=popblue,colback=popblueL,before upper={\popboxhead{Experiment}{popblue}{#1}}}
% « Le savais-tu ? » : carte violette pleine (contexte, culture, Cameroun)
\newtcolorbox{savaistu}[1][]{popbase,colback=poppurple,colframe=popink,
  fontupper=\popsemi\fontsize{10.4}{14.2}\selectfont\color{white},
  before upper={\poppill{Did you know?}{white}{poppurple}\ifx\relax#1\relax\else\hspace{6pt}{\popxbold\color{white}#1}\fi\par\vspace{7pt}}}
% Exercice résolu : énoncé en haut, \tcblower puis la solution sur fond crème
\newcounter{popexo}\newcounter{popexr}
\newtcolorbox{exoresolu}[1][]{popbase,colframe=poporange,colbacklower=popcream,
  segmentation style={popink,line width=1.2pt,dashed},
  before upper={\stepcounter{popexr}\popboxhead{Worked example \thepopexr}{poporange}{#1}},
  before lower={\poppill{Solution}{popink}{popcream}\par\vspace{6pt}}}
% Exercice (non résolu en place) — la correction est dans la partie Corrigés
\newtcolorbox{exercice}[1][]{popbase,colframe=popink,
  before upper={\stepcounter{popexo}\popboxhead{Exercise \thepopexo}{popink}{#1}}}
% Corrigé
\newtcolorbox{corrige}[1][]{popbase,colframe=popgreen,colback=popgreenL,
  before upper={\popboxhead{Solution}{popgreen}{#1}}}
% Objectifs / prérequis / bilan
\newtcolorbox{objectifs}{popbase,colframe=popink,colback=poporangeL,
  before upper={\popboxhead{Chapter objectives}{popink}{}}}
\newtcolorbox{prerequis}{popbase,colframe=popink,colback=popblueL,
  before upper={\popboxhead{Prerequisites}{popblue}{}}}
\newtcolorbox{bilan}{popbase,colframe=popink,colback=white,boxrule=2.8pt,
  before upper={\popboxhead{Fiche bilan}{poporange}{L'essentiel à connaître pour l'examen}}}
% Figure encadrée avec légende
\newtcolorbox{popfigure}[1][]{enhanced,breakable=false,colback=white,colframe=popline,boxrule=1.4pt,arc=8pt,
  left=10pt,right=10pt,top=10pt,bottom=8pt,before skip=10pt,after skip=12pt,halign=center,
  lower separated=false,halign lower=center,
  fontlower=\popsemi\fontsize{9}{11.5}\selectfont\color{popmuted},
  #1}
\newcommand{\legende}[1]{\par\vspace{6pt}{\popsemi\fontsize{9}{11.5}\selectfont\color{popmuted}#1}}

% ============================================================
% Signature OnBuch+
% ============================================================
\newcommand{\poplogoinline}[1]{{\poplogo\fontsize{#1}{#1}\selectfont\color{popink}OnBuch\color{poporange}+}}
\newcommand{\signatureonbuch}{%
  \begin{tcolorbox}[enhanced,colback=popink,colframe=popink,boxrule=2.2pt,arc=10pt,
      drop shadow=poporange,left=14pt,right=14pt,top=10pt,bottom=10pt,before skip=0pt,after skip=0pt]
    \tikz[baseline=-0.7ex]\node[circle,fill=popgreen,draw=popcream,line width=1.3pt,minimum size=22pt,inner sep=0]{\popcheck{white}};%
    \hspace{9pt}\begin{minipage}[c]{0.62\linewidth}
      {\popxbold\fontsize{12}{13}\selectfont\color{popcream}Signed\ \textbullet\ {\poplogo OnBuch\color{poporange}+}}\par\vspace{2pt}
      {\raggedright\popsemi\fontsize{8}{10.5}\selectfont\color{popline}\MakeUppercase{OnBuch+ teaching team}\\\MakeUppercase{Follows the official MINESEC syllabus}\par}
    \end{minipage}\hfill
    {\poplogo\fontsize{22}{22}\selectfont\color{popcream}OnBuch\color{poporange}+}
  \end{tcolorbox}%
}

% ============================================================
% Page de garde
% #1 titre · #2 matière · #3 niveau · #4 module · #5 n° leçon · #6 durée ·
% #7 description / objectifs (texte ou itemize)
% ============================================================
\newcommand{\pagegarde}[7]{%
  \thispagestyle{empty}
  \newgeometry{a4paper,margin=1.7cm,top=1.6cm,bottom=1.4cm}%
  \begin{tikzpicture}[remember picture,overlay]
    \fill[poporange!18] ([xshift=-3.2cm,yshift=-3.6cm]current page.north east) circle (4.6cm);
    \fill[poppurple!14] ([xshift=2.4cm,yshift=7.6cm]current page.south west) circle (3.8cm);
  \end{tikzpicture}%
  {\poplogo\fontsize{30}{30}\selectfont\color{popink}OnBuch\color{poporange}+}\hfill
  \raisebox{0.6ex}{\poppill{Signed course \textbullet\ Premium}{popink}{popcream}}\par
  \vspace{1pt}\popsquiggle{poppurple}\par
  \vspace{8pt}
  \begin{tcolorbox}[enhanced,colback=poporange,colframe=popink,boxrule=2.8pt,arc=13pt,
      drop shadow=popink,left=18pt,right=18pt,top=12pt,bottom=13pt,after skip=10pt]
    \poppill{#2 \textbullet\ #3}{popink}{popcream}\hspace{5pt}\poppill{#4}{white}{popink}\par\vspace{8pt}
    {\popdisplay\fontsize{14}{14}\selectfont\color{popink}CHAPTER #5}\par\vspace{4pt}
    {\raggedright\hyphenpenalty=10000\exhyphenpenalty=10000\popdisplay\fontsize{25}{27}\selectfont\color{popcream}\MakeUppercase{#1}\par}
    \vspace{8pt}
    {\popxbold\fontsize{9.5}{9.5}\selectfont\color{popink}\MakeUppercase{Suggested time: #6}}
  \end{tcolorbox}
  \begin{tcolorbox}[enhanced,colback=white,colframe=popink,boxrule=2.2pt,arc=10pt,
      drop shadow=popink,left=16pt,right=16pt,top=10pt,bottom=10pt,after skip=8pt]
    \poppill{In this chapter}{poppurple}{white}\par\vspace{5pt}
    \popsemi\fontsize{9.3}{12.6}\selectfont\color{popink}#7
  \end{tcolorbox}
  \noindent\begin{minipage}[t]{0.487\linewidth}
    \begin{tcolorbox}[enhanced,colback=popgreen,colframe=popink,boxrule=2.2pt,arc=10pt,
        drop shadow=popink,left=11pt,right=11pt,top=10pt,bottom=10pt,height=3.1cm,valign=center]
      \tcbox[colback=white,colframe=popink,boxrule=1.3pt,arc=5pt,boxsep=0pt,nobeforeafter,
        left=3pt,right=3pt,top=3pt,bottom=3pt,valign=center]{\qrcode[height=1.9cm]{https://play.google.com/store/apps/details?id=com.luvvix.onbuch&hl=en}}%
      \hspace{8pt}\begin{minipage}[c]{0.5\linewidth}
        {\popdisplay\fontsize{11}{12.5}\selectfont\color{white}\MakeUppercase{Download OnBuch}}\par\vspace{4pt}
        {\raggedright\popsemi\fontsize{8.4}{11}\selectfont\color{white}Free \textbullet\ Google Play\\AI tutor, past papers, results\par}
      \end{minipage}
    \end{tcolorbox}
  \end{minipage}\hfill
  \begin{minipage}[t]{0.487\linewidth}
    \begin{tcolorbox}[enhanced,colback=poppurple,colframe=popink,boxrule=2.2pt,arc=10pt,
        drop shadow=popink,left=11pt,right=11pt,top=10pt,bottom=10pt,height=3.1cm,valign=center]
      \tikz[baseline=-0.7ex]\node[circle,fill=white,draw=popink,line width=1.3pt,minimum size=1.6cm,inner sep=0]{\popdisplay\fontsize{24}{24}\selectfont\color{poppurple}+};%
      \hspace{8pt}\begin{minipage}[c]{0.6\linewidth}
        {\popdisplay\fontsize{11}{12.5}\selectfont\color{white}\MakeUppercase{Go OnBuch+}}\par\vspace{4pt}
        {\raggedright\popsemi\fontsize{8.4}{11}\selectfont\color{white}All signed courses, unlimited AI tutor, PDF sheets — OnBuch+ tab in the app\par}
      \end{minipage}
    \end{tcolorbox}
  \end{minipage}
  \vspace{8pt}
  \popcontenu
  \vfill
  \signatureonbuch
  \restoregeometry
  \newpage
}

% Rangée « Ce cours contient » (page de garde)
\newcommand{\poptile}[3]{% #1 couleur #2 gros texte #3 légende
  \begin{tcolorbox}[enhanced,colback=#1,colframe=popink,boxrule=2pt,arc=9pt,shadow={2.5pt}{-2.5pt}{0pt}{popink},
      left=6pt,right=6pt,top=9pt,bottom=9pt,halign=center,valign=center,height=1.75cm,width=\linewidth,nobeforeafter]
    {\popdisplay\fontsize{12.5}{13}\selectfont\color{white}\MakeUppercase{#2}}\par\vspace{3pt}
    {\centering\popsemi\fontsize{7.8}{9.5}\selectfont\color{white}#3\par}
  \end{tcolorbox}}
\newcommand{\popcontenu}{%
  \par\noindent{\popxboldls\fontsize{7.5}{7.5}\selectfont\color{popmuted}THIS SIGNED COURSE CONTAINS}\par\vspace{7pt}
  \noindent\begin{minipage}[t]{0.235\linewidth}\poptile{popblue}{Course}{complete, illustrated, step by step}\end{minipage}\hfill
  \begin{minipage}[t]{0.235\linewidth}\poptile{poporange}{Methods}{and worked examples}\end{minipage}\hfill
  \begin{minipage}[t]{0.235\linewidth}\poptile{poppink}{Exercises}{exam style + solutions}\end{minipage}\hfill
  \begin{minipage}[t]{0.235\linewidth}\poptile{popgreen}{Summary sheet}{the essentials to remember}\end{minipage}\par}

% Sommaire
\newtcolorbox{sommaire}{enhanced,colback=white,colframe=popink,boxrule=2.2pt,arc=10pt,
  drop shadow=popink,left=18pt,right=18pt,top=13pt,bottom=13pt,before skip=6pt,after skip=20pt,
  before upper={\poppill{Contents}{poppurple}{white}\par\vspace{9pt}\popsemi\fontsize{10.6}{14.5}\selectfont\color{popink}}}

% Signature de fin de cours — poussée tout en bas de la dernière page
\newcommand{\findecours}{%
  \par\vfill\nopagebreak
  \begin{center}\popsquiggle{poporange}\end{center}\vspace{4pt}
  \signatureonbuch
}
\AtBeginDocument{\def\llbracket{\mathopen{[\mkern-3.5mu[}}\def\rrbracket{\mathclose{]\mkern-3.5mu]}}} % intervalles d'entiers (glyphe ⟦ absent de la police)
% Glyphes absents de Fira Math : redéfinis à partir de glyphes disponibles
\AtBeginDocument{%
  \def\setminus{\mathbin{\backslash}}\let\smallsetminus\setminus
  \def\vdots{\mathord{\vbox{\baselineskip3.2pt\lineskiplimit0pt\kern4pt\hbox{.}\hbox{.}\hbox{.}}}}%
  \def\ddots{\mathinner{\mkern1mu\raise6.5pt\hbox{.}\mkern2mu\raise3.6pt\hbox{.}\mkern2mu\raise0.7pt\hbox{.}\mkern1mu}}%
  \def\triangle{\mathord{\scalebox{0.9}{$\symup{\Delta}$}}}%
  \def\nabla{\mathord{\raisebox{\depth}{\scalebox{1}[-1]{$\symup{\Delta}$}}}}%
  \def\checkmark{\popcheck{popdark}}%
  \def\star{\mathbin{\ast}}\def\bigstar{\mathord{\ast}}%
  \def\diamond{\mathbin{\scalebox{0.6}{$\Diamond$}}}%
  \def\bigcup{\mathop{\mathchoice{\raisebox{-0.3em}{\scalebox{1.7}{$\cup$}}}{\scalebox{1.25}{$\cup$}}{\cup}{\cup}}\displaylimits}%
  \def\bigcap{\mathop{\mathchoice{\raisebox{-0.3em}{\scalebox{1.7}{$\cap$}}}{\scalebox{1.25}{$\cap$}}{\cap}{\cap}}\displaylimits}%
  \def\lesseqqgtr{\mathrel{\lessgtr}}\def\bigoplus{\mathop{\oplus}\displaylimits}%
}
\providecommand{\ssi}{if and only if} % abréviation fréquente
\AtBeginDocument{\providecommand{\wideparen}{\overparen}} % arc de cercle

%% FIGFIX-BEGIN v4 — correctifs globaux des figures (docs/onbuch-generation/tools/figfix)
\usepackage{adjustbox}\usepackage{etoolbox}
% toute figure TikZ/pgfplots plus large que la ligne est réduite à la largeur disponible
\BeforeBeginEnvironment{tikzpicture}{\begin{adjustbox}{max width=\linewidth}}
\AfterEndEnvironment{tikzpicture}{\end{adjustbox}}
% légendes pgfplots lisibles par-dessus les courbes
\pgfplotsset{every axis/.append style={legend style={fill=white,fill opacity=0.92,text opacity=1,draw=black!15,inner sep=3pt}}}
% étiquettes d'arêtes (syntaxe edge["..."]) avec fond blanc
\tikzset{every edge quotes/.append style={fill=white,inner sep=1.2pt}}
%% FIGFIX-END
\begin{document}

\pagegarde{Circle Geometry}{Pure Mathematics with Mechanics}{Upper Sixth}{Upper Sixth — Pure mathematics with mechanics}{1}{22 periods}{This chapter develops the complete analytic study of the circle in the Cartesian plane: its various equations, tangents, intersections with lines and other circles, parametric form, radical axis and locus problems. It is one of the most frequently examined topics in the GCE Advanced Level Pure Mathematics with Mechanics paper.
\vspace{4pt}
\begin{multicols}{2}\small
\begin{itemize}
\item By the end of this chapter I can write the equation of a circle given its centre and radius, a diameter, or three points on it.
\item By the end of this chapter I can convert the general equation x² + y² + 2gx + 2fy + c = 0 to centre-radius form and interpret g, f, c.
\item By the end of this chapter I can find the equation of a tangent to a circle at a given point or from an external point.
\item By the end of this chapter I can state and apply the condition for a line to touch a circle.
\item By the end of this chapter I can use and convert parametric equations of a circle.
\item By the end of this chapter I can compute the length of a tangent from an external point.
\end{itemize}\end{multicols}}

\begin{sommaire}
\begin{multicols}{2}
\begin{enumerate}
\item Section 1: Equations of a Circle — Centre, Radius, Diameter
\item Section 2: The General Equation and Circles Through Three Points
\item Section 3: Tangents to a Circle
\item Section 4: Parametric Equations and Length of Tangent
\item Section 5: Intersection of a Circle and a Line; Two Circles
\item Section 6: Families of Circles Through Intersections
\item Section 7: Radical Axis and Locus Problems
\item Section 8: Circles Under Constraints — Touching Axes and Lines
\item Section 9: Synthesis, GCE Techniques and Revision
\item Integration activity
\item Methods and worked examples
\item Exercises
\item Solutions to the exercises
\item Summary sheet
\end{enumerate}
\end{multicols}
\end{sommaire}

\begin{objectifs}
\begin{itemize}
\item By the end of this chapter I can write the equation of a circle given its centre and radius, a diameter, or three points on it.
\item By the end of this chapter I can convert the general equation x² + y² + 2gx + 2fy + c = 0 to centre-radius form and interpret g, f, c.
\item By the end of this chapter I can find the equation of a tangent to a circle at a given point or from an external point.
\item By the end of this chapter I can state and apply the condition for a line to touch a circle.
\item By the end of this chapter I can use and convert parametric equations of a circle.
\item By the end of this chapter I can compute the length of a tangent from an external point.
\item By the end of this chapter I can find the intersections of a circle with a line or with another circle.
\item By the end of this chapter I can write the equation of a circle through the intersection of two circles, or of a line and a circle, using the family-of-circles method.
\item By the end of this chapter I can determine the radical axis of two circles and the conditions for two circles to touch.
\item By the end of this chapter I can solve locus problems and constrained circle problems (touching axes, touching a line, centre on a line).
\end{itemize}
\end{objectifs}

\begin{prerequis}
\begin{itemize}
\item Coordinate geometry of the straight line: gradients, equations of lines, perpendicularity condition (Lower Sixth).
\item Distance between two points and midpoint formula.
\item Pythagoras' theorem and basic circle theorems from O-Level geometry (angle in a semicircle, perpendicular from centre to chord).
\item Solving simultaneous equations, including one linear and one quadratic.
\item Completing the square for a quadratic expression.
\end{itemize}
\end{prerequis}

\newpage

\coursec{Section 1: Equations of a Circle — Centre, Radius, Diameter}

In Yaoundé, a telecommunications engineer must position a relay antenna so that its signal covers three villages located at known points on a map: the coverage zone is a circle, and the antenna must sit at its centre. In Douala, a roundabout is designed so that a straight road just grazes it — a tangent problem. When a water tank's circular rim must pass through three given pillars, or when two circular irrigation zones overlap, every question reduces to finding the equation of a circle and studying its position relative to lines and other circles. This chapter gives you the algebraic tools to solve all these problems with precision.

\courssub{The Circle as a Locus — Definition and the Origin-Centred Circle}

\begin{definition}[Circle as a Locus]
A \cle{circle} is the set of all points in a plane that are at a fixed distance, called the \cle{radius}, from a fixed point called the \cle{centre}.
\end{definition}

Let the centre be the origin $O(0,0)$ and the radius be $r > 0$. For any point $P(x,y)$ on the circle, the distance $OP$ equals $r$. By the distance formula:
\[
OP = \sqrt{(x-0)^2 + (y-0)^2} = \sqrt{x^2 + y^2}.
\]
Hence $\sqrt{x^2 + y^2} = r$, and squaring both sides (both are non-negative, so squaring preserves equality) gives the fundamental equation.

\begin{propriete}[Circle Centred at the Origin]
The equation of a circle with centre $O(0,0)$ and radius $r$ is
\[
\boxed{x^2 + y^2 = r^2}.
\]
\end{propriete}

\begin{popfigure}
\begin{tikzpicture}[scale=1.2]
\draw[popline, thick] (0,0) circle (2.5);
\draw[->, popdark] (-3.2,0) -- (3.2,0) node[right] {$x$};
\draw[->, popdark] (0,-3.2) -- (0,3.2) node[above] {$y$};
\coordinate (O) at (0,0);
\coordinate (P) at (45:2.5);
\coordinate (Q) at (2.5*cos 45,0);
\draw[popblue, thick] (O) -- (P) node[midway, above left] {$r$};
\draw[poporange, dashed] (P) -- (Q) node[midway, right] {$y$};
\draw[popgreen, thick] (O) -- (Q) node[midway, below] {$x$};
\fill[popink] (O) circle (2pt) node[below left] {$O(0,0)$};
\fill[popblue] (P) circle (2pt) node[above right] {$P(x,y)$};
\draw (0.5,0) arc (0:45:0.5) node[midway, right, xshift=2pt] {$\theta$};
\end{tikzpicture}
\legende{Circle centred at the origin: right triangle with legs $|x|$, $|y|$ and hypotenuse $r$, so $x^2+y^2=r^2$ by Pythagoras.}
\end{popfigure}

\begin{exemplebox}[Reading the radius from the equation]
Find the radius of the circle $x^2 + y^2 = 50$.
\end{exemplebox}
\[
r^2 = 50 \quad\Rightarrow\quad r = \sqrt{50} = 5\sqrt{2}.
\]

\begin{exemplebox}[Writing the equation given the radius]
Write the equation of the circle centred at the origin with radius $\sqrt{13}$.
\end{exemplebox}
\[
x^2 + y^2 = (\sqrt{13})^2 = 13.
\]

\courssub{Translation of the Centre — Equation $(x-a)^2+(y-b)^2=r^2$}

When the centre is $C(a,b)$, the distance $CP$ for $P(x,y)$ is
\[
CP = \sqrt{(x-a)^2 + (y-b)^2}.
\]
Setting $CP = r$ and squaring yields the standard form.

\begin{propriete}[Standard Equation of a Circle]
The circle with centre $C(a,b)$ and radius $r$ has equation
\[
\boxed{(x-a)^2 + (y-b)^2 = r^2}.
\]
\end{propriete}

\begin{popfigure}
\begin{tikzpicture}[scale=1.1]
\draw[popline, thick] (2,1) circle (2.2);
\draw[->, popdark] (-1,0) -- (5,0) node[right] {$x$};
\draw[->, popdark] (0,-1.5) -- (0,4) node[above] {$y$};
\coordinate (C) at (2,1);
\coordinate (P) at (2+2.2*cos 50, 1+2.2*sin 50);
\coordinate (Q) at (2+2.2*cos 50,1);
\draw[popblue, thick] (C) -- (P) node[midway, above left] {$r$};
\draw[poporange, dashed] (P) -- (Q) node[midway, right] {$y-b$};
\draw[popgreen, thick] (C) -- (Q) node[midway, below] {$x-a$};
\fill[popink] (C) circle (2pt) node[below left] {$C(a,b)$};
\fill[popblue] (P) circle (2pt) node[above right] {$P(x,y)$};
\draw (2.4,1) arc (0:50:0.4) node[midway, right, xshift=2pt] {$\theta$};
\end{tikzpicture}
\legende{Circle with centre $C(a,b)$: the right triangle has legs $x-a$ and $y-b$, so $(x-a)^2+(y-b)^2=r^2$.}
\end{popfigure}

\begin{aretenir}[Reading Centre and Radius]
From $(x-a)^2 + (y-b)^2 = r^2$:
\begin{itemize}
\item Centre: $C(a,b)$.
\item Radius: $r = \sqrt{r^2}$ (positive square root).
\end{itemize}
\textbf{Watch the signs:} $(x-3)^2$ means $a=3$; $(x+2)^2 = (x-(-2))^2$ means $a=-2$.
\end{aretenir}

\begin{attention}[Sign Trap]
In $(x-a)^2+(y-b)^2=r^2$ the centre is $(a,b)$, \textbf{not} $(-a,-b)$.
Example: $(x+4)^2+(y-5)^2=9$ has centre $(-4,5)$, radius $3$.
\end{attention}

\begin{exoresolu}[Centre and radius from the equation]
Find the centre and radius of the circle $(x-2)^2 + (y+3)^2 = 25$.
\tcblower
Comparing with $(x-a)^2+(y-b)^2=r^2$:
$a=2$, $b=-3$, $r^2=25 \Rightarrow r=5$.
Centre $C(2,-3)$, radius $5$.
\end{exoresolu}

\begin{exoresolu}[Writing the equation given centre and radius]
Find the equation of the circle with centre $C(-1,4)$ and radius $6$.
\tcblower
$a=-1$, $b=4$, $r=6 \Rightarrow r^2=36$.
\[
(x-(-1))^2 + (y-4)^2 = 36 \quad\Rightarrow\quad (x+1)^2 + (y-4)^2 = 36.
\]
\end{exoresolu}

\courssub{Circle Given Centre and a Point on the Circle (Lesson 3)}

If the centre $C(a,b)$ is known and a point $P(x_1,y_1)$ lies on the circle, the radius is the distance $CP$:
\[
r = CP = \sqrt{(x_1-a)^2 + (y_1-b)^2}.
\]
Then the equation is $(x-a)^2+(y-b)^2 = r^2$.

\begin{methode}[Circle from Centre and a Point]
\begin{enumerate}
\item Compute $r^2 = (x_1-a)^2 + (y_1-b)^2$.
\item Write $(x-a)^2 + (y-b)^2 = r^2$.
\end{enumerate}
\end{methode}

\begin{exoresolu}[Centre and a point on the circle]
Find the equation of the circle with centre $C(3,-2)$ passing through $P(6,2)$.
\tcblower
$r^2 = (6-3)^2 + (2-(-2))^2 = 3^2 + 4^2 = 9+16 = 25$.
Equation: $(x-3)^2 + (y+2)^2 = 25$.
\end{exoresolu}

\courssub{Circle on a Given Diameter (Lesson 5)}

Let $A(x_1,y_1)$ and $B(x_2,y_2)$ be the endpoints of a diameter.

\begin{propriete}[Centre and Radius from Diameter]
\begin{itemize}
\item Centre $M$ is the midpoint of $AB$:
\[
M\left(\frac{x_1+x_2}{2},\ \frac{y_1+y_2}{2}\right).
\]
\item Radius $r = \dfrac{1}{2}AB = \dfrac{1}{2}\sqrt{(x_2-x_1)^2 + (y_2-y_1)^2}$.
\end{itemize}
\end{propriete}

\begin{propriete}[Alternative Form — Angle in a Semicircle]
For any point $P(x,y)$ on the circle, $\angle APB = 90^\circ$ (angle in a semicircle), so $\overrightarrow{PA} \cdot \overrightarrow{PB} = 0$:
\[
(x-x_1)(x-x_2) + (y-y_1)(y-y_2) = 0.
\]
This is the \cle{diameter form} of the circle equation.
\end{propriete}

\begin{popfigure}
\begin{tikzpicture}[scale=1.1]
\draw[popline, thick] (1,1) circle (2);
\coordinate (A) at (-1,1);
\coordinate (B) at (3,1);
\coordinate (M) at (1,1);
\coordinate (P) at (1+2*cos 110, 1+2*sin 110);
\draw[popdark, thick] (A) -- (B);
\draw[popblue, thick] (A) -- (P) -- (B);
\draw[poporange, dashed] (M) -- (P);
\fill[popink] (A) circle (2pt) node[left] {$A(x_1,y_1)$};
\fill[popink] (B) circle (2pt) node[right] {$B(x_2,y_2)$};
\fill[poppurple] (M) circle (2pt) node[below] {$M$};
\fill[popgreen] (P) circle (2pt) node[above left] {$P(x,y)$};
\draw[popdark] (P) ++(20:0.25) arc (20:110:0.25);
\end{tikzpicture}
\legende{Circle on diameter $AB$: $\angle APB = 90^\circ$ gives $(x-x_1)(x-x_2)+(y-y_1)(y-y_2)=0$.}
\end{popfigure}

\begin{methode}[Circle on Diameter $AB$]
\begin{enumerate}
\item \textbf{Midpoint method:} Find $M\left(\frac{x_1+x_2}{2},\frac{y_1+y_2}{2}\right)$, compute $r = \frac{1}{2}AB$, write $(x-x_M)^2+(y-y_M)^2=r^2$.
\item \textbf{Diameter form (often faster):} Write directly $(x-x_1)(x-x_2)+(y-y_1)(y-y_2)=0$ and expand if needed.
\end{enumerate}
\end{methode}

\begin{exoresolu}[Circle on a diameter]
Find the equation of the circle on $A(2,-1)$ and $B(6,3)$ as diameter.
\tcblower
\textbf{Method 1 (Midpoint):}
$M\left(\frac{2+6}{2},\frac{-1+3}{2}\right) = (4,1)$.
$AB = \sqrt{(6-2)^2+(3-(-1))^2} = \sqrt{16+16} = \sqrt{32} = 4\sqrt{2}$.
$r = 2\sqrt{2}$, $r^2 = 8$.
Equation: $(x-4)^2 + (y-1)^2 = 8$.

\textbf{Method 2 (Diameter form):}
$(x-2)(x-6) + (y+1)(y-3) = 0$.
$x^2-8x+12 + y^2-2y-3 = 0 \;\Rightarrow\; x^2+y^2-8x-2y+9=0$.
Complete squares: $(x-4)^2-16 + (y-1)^2-1 + 9 = 0 \;\Rightarrow\; (x-4)^2+(y-1)^2=8$. \checkmark
\end{exoresolu}

\courssub{Position of a Point Relative to a Circle}

\begin{propriete}[Point Inside / On / Outside]
Let circle $C$ have centre $O(a,b)$ and radius $r$. For a point $P(x_0,y_0)$, compute $d = OP = \sqrt{(x_0-a)^2+(y_0-b)^2}$.
\begin{itemize}
\item $d < r$ \quad$\Rightarrow$\quad $P$ is \cle{inside} the circle.
\item $d = r$ \quad$\Rightarrow$\quad $P$ lies \cle{on} the circle.
\item $d > r$ \quad$\Rightarrow$\quad $P$ is \cle{outside} the circle.
\end{itemize}
Equivalently, compare $(x_0-a)^2+(y_0-b)^2$ with $r^2$ (avoids the square root).
\end{propriete}

\begin{exoresolu}[Position of a point]
Determine whether $P(5,2)$ lies inside, on, or outside the circle $(x-1)^2+(y+2)^2=20$.
\tcblower
Centre $C(1,-2)$, $r^2=20$.
$CP^2 = (5-1)^2 + (2-(-2))^2 = 4^2 + 4^2 = 16+16 = 32$.
Since $32 > 20$, $CP^2 > r^2$, so $P$ is \textbf{outside} the circle.
\end{exoresolu}

\begin{exercice}[Practice 1.1]
\begin{enumerate}
\item Write the equation of the circle centred at the origin with radius $\sqrt{17}$.
\item Find the centre and radius of $(x+3)^2+(y-4)^2=49$.
\item Find the equation of the circle with centre $(-2,5)$ passing through $(1,9)$.
\item Find the equation of the circle on $A(-3,2)$ and $B(5,-4)$ as diameter (use both methods).
\item Determine the position of $Q(0,0)$ relative to the circle $(x-2)^2+(y+3)^2=10$.
\end{enumerate}
\end{exercice}

\begin{corrige}[Practice 1.1]
\begin{enumerate}
\item $x^2+y^2=17$.
\item Centre $(-3,4)$, radius $7$.
\item $r^2=(1+2)^2+(9-5)^2=9+16=25$; $(x+2)^2+(y-5)^2=25$.
\item Midpoint $M(1,-1)$, $AB=\sqrt{8^2+(-6)^2}=10$, $r=5$; $(x-1)^2+(y+1)^2=25$. Diameter form: $(x+3)(x-5)+(y-2)(y+4)=0 \Rightarrow x^2-2x-15+y^2+2y-8=0 \Rightarrow (x-1)^2+(y+1)^2=25$.
\item $C(2,-3)$, $r^2=10$. $CQ^2=(0-2)^2+(0+3)^2=4+9=13 > 10$; $Q$ is outside.
\end{enumerate}
\end{corrige}

\begin{savaistu}[Historical Note]
The Greek mathematician Apollonius of Perga (c. 262–190 BC) systematically studied conic sections, including the circle. The coordinate approach we use today was developed by René Descartes (1596–1650) and Pierre de Fermat (1607–1665), linking geometry and algebra — the foundation of analytic geometry taught in Cameroonian lycées since the 1960s.
\end{savaistu}

\coursec{Section 2: The General Equation and Circles Through Three Points}

In Section 1 we built the equation of a circle from its geometric definition: the set of points at a fixed distance $r$ from a fixed centre $(a,b)$. This gave us the \cle{centre-radius form}
\[
(x-a)^2+(y-b)^2=r^2.
\]
This form is ideal when you \emph{know} the centre and radius. But in many GCE problems the equation arrives already expanded, looking like $x^2+y^2-6x+4y-12=0$, and nothing about it immediately reveals the centre or radius. Worse, sometimes you are not given the centre at all --- only three points the circle must pass through. This section gives you the two essential tools: how to \emph{recognise and decode} an expanded circle equation (the \cle{general form}), and how to \emph{construct} a circle from three non-collinear points.

\courssub{From Centre-Radius Form to the General Equation}

Start from the centre-radius form and simply expand the brackets:
\begin{align*}
(x-a)^2+(y-b)^2 &= r^2\\
x^2-2ax+a^2+y^2-2by+b^2 &= r^2\\
x^2+y^2-2ax-2by+(a^2+b^2-r^2) &= 0.
\end{align*}
Now rename the constants: put $g=-a$, $f=-b$ and $c=a^2+b^2-r^2$. The equation becomes the famous \cle{general equation of a circle}.

\begin{definition}[General equation of a circle]
The \cle{general equation} of a circle is
\[
x^2+y^2+2gx+2fy+c=0,
\]
where $g$, $f$, $c$ are real constants. Its \cle{centre} is $(-g,-f)$ and its \cle{radius} is
\[
r=\sqrt{g^2+f^2-c}.
\]
\end{definition}

Notice the bookkeeping: the centre coordinates are \emph{minus half} the coefficients of $x$ and $y$ (because the coefficients are $2g$ and $2f$), and the constant term $c$ absorbs $a^2+b^2-r^2$.

\begin{methode}[Completing the square: general form $\to$ centre-radius form]
Given $x^2+y^2+2gx+2fy+c=0$:
\begin{enumerate}
\item Group the $x$-terms and the $y$-terms: $(x^2+2gx)+(y^2+2fy)=-c$.
\item Complete each square: add $g^2$ and $f^2$ to \textbf{both} sides:
\[
(x+g)^2-g^2+(y+f)^2-f^2=-c.
\]
\item Rearrange: $(x+g)^2+(y+f)^2=g^2+f^2-c$.
\item Read off: centre $(-g,-f)$, radius $\sqrt{g^2+f^2-c}$.
\end{enumerate}
\end{methode}

\begin{exemplebox}[Decoding a general equation]
Find the centre and radius of $x^2+y^2-6x+4y-12=0$.

\emph{By formula:} $2g=-6\Rightarrow g=-3$; $2f=4\Rightarrow f=2$; $c=-12$.
Centre $=(-g,-f)=(3,-2)$, radius $=\sqrt{9+4+12}=\sqrt{25}=5$.

\emph{By completing the square:} $(x^2-6x)+(y^2+4y)=12$, so $(x-3)^2-9+(y+2)^2-4=12$, i.e. $(x-3)^2+(y+2)^2=25$. Same answer --- and this is exactly the circle of Section 1.
\end{exemplebox}

\begin{popfigure}
\begin{center}
\begin{tikzpicture}
\begin{axis}[
  width=9cm, axis equal image,
  axis lines=middle, xlabel=$x$, ylabel=$y$,
  xmin=-3, xmax=9, ymin=-8, ymax=4,
  xtick={-2,0,2,4,6,8}, ytick={-6,-4,-2,0,2},
  grid=both, grid style={popline!40},
]
\addplot[popblue, thick, domain=0:360, samples=200] ({3+5*cos(x)},{-2+5*sin(x)});
\addplot[only marks, mark=*, mark size=2pt, poporange] coordinates {(3,-2)};
\node[popdark, anchor=south west] at (axis cs:3.1,-1.9) {$C(3,-2)$};
\draw[poporange, thick, ->] (axis cs:3,-2) -- (axis cs:8,-2) node[midway, above, popdark] {$r=5$};
\end{axis}
\end{tikzpicture}
\end{center}
\legende{The circle $x^2+y^2-6x+4y-12=0$: centre $(-g,-f)=(3,-2)$, radius $\sqrt{g^2+f^2-c}=5$.}
\end{popfigure}

\courssub{Does the General Equation Always Give a Circle?}

\begin{attention}[When is it really a circle?]
The radius formula involves $\sqrt{g^2+f^2-c}$, so we must examine the sign of $g^2+f^2-c$:
\begin{itemize}
\item $g^2+f^2-c>0$: a genuine circle of positive radius;
\item $g^2+f^2-c=0$: a \cle{point circle} --- the locus is the single point $(-g,-f)$;
\item $g^2+f^2-c<0$: \textbf{no real locus at all} (a sum of two squares cannot be negative).
\end{itemize}
For example $x^2+y^2+2x+4y+9=0$ gives $g^2+f^2-c=1+4-9=-4<0$: no real point satisfies it.
\end{attention}

\begin{attention}[Not every quadratic is a circle]
An equation represents a circle only if
\begin{enumerate}
\item the coefficients of $x^2$ and $y^2$ are \textbf{equal} (and non-zero), and
\item there is \textbf{no $xy$ term}.
\end{enumerate}
Thus $2x^2+2y^2+3x-4y+1=0$ \emph{is} a circle (divide by $2$ first!), but $x^2+2y^2-4x=0$ is an ellipse and $x^2+y^2+3xy=1$ is not a circle either. Always divide through to make the coefficients of $x^2$ and $y^2$ equal to $1$ before reading off $g$, $f$, $c$.
\end{attention}

\begin{aretenir}[Recognising and decoding a circle equation]
\begin{itemize}
\item A circle equation has \textbf{equal coefficients of $x^2$ and $y^2$} and \textbf{no $xy$ term}.
\item After normalising to $x^2+y^2+2gx+2fy+c=0$: centre $(-g,-f)$, radius $\sqrt{g^2+f^2-c}$.
\item Real circle $\iff g^2+f^2-c>0$; equality gives a single point; negative gives nothing.
\item If in doubt, complete the square --- it never lies.
\end{itemize}
\end{aretenir}

\courssub{Equation of a Circle Through Three Points}

Here is a beautiful fact you already know from O-Level geometry: \emph{through any three non-collinear points there passes exactly one circle} --- the circumcircle of the triangle they form, whose centre is the intersection of the perpendicular bisectors of the sides. Algebraically, the general equation has exactly \textbf{three unknowns} $g$, $f$, $c$, so three points give three linear equations: the system has a unique solution. Geometry and algebra agree perfectly.

\begin{methode}[Circle through three given points --- algebraic method]
To find the circle through $P_1(x_1,y_1)$, $P_2(x_2,y_2)$, $P_3(x_3,y_3)$ (non-collinear):
\begin{enumerate}
\item Write the general equation $x^2+y^2+2gx+2fy+c=0$.
\item Substitute each point in turn to obtain three linear equations in $g$, $f$, $c$.
\item Solve the simultaneous system (elimination is usually fastest).
\item State the equation; if required, give centre $(-g,-f)$ and radius $\sqrt{g^2+f^2-c}$.
\item \textbf{Verify} by substituting one or all of the points back.
\end{enumerate}
\end{methode}

\begin{exoresolu}[GCE-style: circle through three points]
Find the equation of the circle passing through $A(1,2)$, $B(5,4)$ and $C(-3,2)$. Give its centre and radius, and verify that all three points lie on it.
\tcblower
\textbf{Solution.} Let the circle be $x^2+y^2+2gx+2fy+c=0$.

Substituting $A(1,2)$: $\;1+4+2g+4f+c=0 \Rightarrow 2g+4f+c=-5$ \quad (i)

Substituting $B(5,4)$: $\;25+16+10g+8f+c=0 \Rightarrow 10g+8f+c=-41$ \quad (ii)

Substituting $C(-3,2)$: $\;9+4-6g+4f+c=0 \Rightarrow -6g+4f+c=-13$ \quad (iii)

(ii)$-$(i): $8g+4f=-36 \Rightarrow 2g+f=-9$ \quad (iv)

(i)$-$(iii): $8g=8 \Rightarrow g=1$.

From (iv): $f=-9-2(1)=-11$. From (i): $c=-5-2(1)-4(-11)=-5-2+44=37$.

\[
\[\boxed{x^2+y^2+2x-22y+37=0}\]
\]
Centre: $(-g,-f)=(-1,11)$. Radius: $\sqrt{g^2+f^2-c}=\sqrt{1+121-37}=\sqrt{85}$.

\textbf{Verification} with $A(1,2)$: $1+4+2-44+37=0$ \checkmark; with $B(5,4)$: $25+16+10-88+37=0$ \checkmark; with $C(-3,2)$: $9+4-6-44+37=0$ \checkmark.
\end{exoresolu}

\courssub{The Geometric Viewpoint: Perpendicular Bisectors}

Why does the algebraic method work? Because the centre of the circle is equidistant from all three points, and the locus of points equidistant from two given points is precisely the \cle{perpendicular bisector} of the segment joining them. Hence the centre lies on the perpendicular bisector of $AB$ \emph{and} on that of $BC$: their intersection is the centre. This is exactly the circumcircle construction from O-Level, and it gives an alternative method: find two perpendicular bisectors, intersect them to get the centre, then compute the radius as the distance to any vertex.

\begin{popfigure}
\begin{center}
\begin{tikzpicture}[scale=0.85]
\coordinate (A) at (1,2);
\coordinate (B) at (5,4);
\coordinate (C) at (-3,2);
\coordinate (O) at (-1,11);
\draw[popblue, thick] (O) circle ({sqrt(max(0,85))});
\draw[popdark, thick] (A)--(B)--(C)--cycle;
\coordinate (M1) at ($(A)!0.5!(B)$);
\coordinate (M2) at ($(A)!0.5!(C)$);
\draw[poporange, thick, dashed] ($(M1)!-1.2!(O)$) -- ($(M1)!1.6!(O)$);
\draw[popgreen, thick, dashed] ($(M2)!-0.4!(O)$) -- ($(M2)!1.9!(O)$);
\fill[popink] (A) circle (2pt) node[below right, popdark] {$A$};
\fill[popink] (B) circle (2pt) node[right, popdark] {$B$};
\fill[popink] (C) circle (2pt) node[below left, popdark] {$C$};
\fill[poporange] (O) circle (2pt) node[above, popdark] {centre};
\fill[popmuted] (M1) circle (1.5pt);
\fill[popmuted] (M2) circle (1.5pt);
\end{tikzpicture}
\end{center}
\legende{The perpendicular bisectors of two chords of the triangle meet at the centre of the circumcircle.}
\end{popfigure}

\begin{exemplebox}[Using perpendicular bisectors]
Find the centre of the circle through $A(1,2)$, $B(5,4)$, $C(-3,2)$ geometrically.

Midpoint of $AC$ is $(-1,2)$; $AC$ is horizontal, so its perpendicular bisector is the vertical line $x=-1$. Midpoint of $AB$ is $(3,3)$; gradient of $AB$ is $\tfrac{4-2}{5-1}=\tfrac12$, so the perpendicular bisector has gradient $-2$: $y-3=-2(x-3)$. Setting $x=-1$: $y-3=-2(-4)=8$, so $y=11$. Centre $(-1,11)$ --- exactly as found algebraically. Radius $=\sqrt{(1-(-1))^2+(2-11)^2}=\sqrt{4+81}=\sqrt{85}$. \checkmark
\end{exemplebox}

\begin{attention}[Collinear points]
If the three points are \textbf{collinear}, no circle passes through them --- the linear system will be inconsistent and the perpendicular bisectors will be parallel. Before starting, a quick gradient check can save you time in the exam.
\end{attention}

\begin{savaistu}[Circles everywhere in Cameroon]
The circumcircle idea is behind a practical technique used by carpenters and masons: to find the centre of a circular arc (say, of a round water tank or an arch in a building in Yaound\'e), mark any three points on the arc and construct two perpendicular bisectors with a straightedge --- their intersection is the centre. The same mathematics locates the epicentre of a signal from three receiving stations.
\end{savaistu}

\begin{exercice}[Practice]
\begin{enumerate}
\item Find the centre and radius of each circle, or state that there is no real locus:
\begin{itemize}
\item $x^2+y^2+8x-10y+5=0$
\item $x^2+y^2-4x+6y+13=0$
\item $3x^2+3y^2-12x+6y-1=0$
\end{itemize}
\item Find the equation of the circle through $(0,0)$, $(4,0)$ and $(0,6)$, and state its centre and radius.
\item Find the equation of the circle through $(2,1)$, $(6,5)$ and $(4,-3)$. Verify your answer.
\end{enumerate}
\end{exercice}

\begin{corrige}[Exercise 1]
\textbf{1.} (a) $g=4$, $f=-5$, $c=5$: centre $(-4,5)$, $r=\sqrt{16+25-5}=\sqrt{36}=6$.
(b) $g=-2$, $f=3$, $c=13$: $g^2+f^2-c=4+9-13=0$: point circle at $(2,-3)$.
(c) Divide by $3$: $x^2+y^2-4x+2y-\tfrac13=0$: centre $(2,-1)$, $r=\sqrt{4+1+\tfrac13}=\sqrt{\tfrac{16}{3}}=\tfrac{4}{\sqrt3}$.

\textbf{2.} Substituting $(0,0)$ gives $c=0$. $(4,0)$: $16+8g=0\Rightarrow g=-2$. $(0,6)$: $36+12f=0\Rightarrow f=-3$. Equation: $x^2+y^2-4x-6y=0$; centre $(2,3)$, radius $\sqrt{4+9}=\sqrt{13}$.

\textbf{3.} General form with the three points:
$(2,1)$: $4+1+4g+2f+c=0 \Rightarrow 4g+2f+c=-5$ \quad (i)
$(6,5)$: $36+25+12g+10f+c=0 \Rightarrow 12g+10f+c=-61$ \quad (ii)
$(4,-3)$: $16+9+8g-6f+c=0 \Rightarrow 8g-6f+c=-25$ \quad (iii)

(ii)$-$(i): $8g+8f=-56 \Rightarrow g+f=-7$ \quad (iv)
(iii)$-$(i): $4g-8f=-20 \Rightarrow g-2f=-5$ \quad (v)

(iv)$-$(v): $3f=-2 \Rightarrow f=-\tfrac23$.
From (iv): $g=-7-f=-7+\tfrac23=-\tfrac{19}{3}$.
From (i): $c=-5-4g-2f=-5-4(-\tfrac{19}{3})-2(-\tfrac23)=-5+\tfrac{76}{3}+\tfrac43=-5+\tfrac{80}{3}=\tfrac{65}{3}$.

Equation: $x^2+y^2+2(-\tfrac{19}{3})x+2(-\tfrac23)y+\tfrac{65}{3}=0$, i.e.
\[
\[\boxed{3x^2+3y^2-38x-4y+65=0}\]
\]
Centre: $(\tfrac{19}{3},\tfrac23)$. Radius: $\sqrt{g^2+f^2-c}=\sqrt{\tfrac{361}{9}+\tfrac49-\tfrac{65}{3}}=\sqrt{\tfrac{365-195}{9}}=\sqrt{\tfrac{170}{9}}=\tfrac{\sqrt{170}}{3}$.

\textbf{Verification}:
$(2,1)$: $12+3-76-4+65=0$ \checkmark
$(6,5)$: $108+75-228-20+65=0$ \checkmark
$(4,-3)$: $48+27-152+12+65=0$ \checkmark
\end{corrige}

\begin{aretenir}[Key points of Section 2]
\begin{itemize}
\item General equation: $x^2+y^2+2gx+2fy+c=0$; centre $(-g,-f)$; radius $\sqrt{g^2+f^2-c}$.
\item Completing the square converts general form to centre-radius form --- master it.
\item Real circle $\iff g^2+f^2-c>0$.
\item Three non-collinear points determine a unique circle: substitute into the general form and solve three linear equations, \emph{or} intersect two perpendicular bisectors.
\item Always verify by substituting the given points back into your final equation.
\end{itemize}
\end{aretenir}

\coursec{Section 3: Tangents to a Circle}

In Section 2 we learnt to recognise the general equation $x^{2}+y^{2}+2gx+2fy+c=0$ as a circle and to find the unique circle through three given points. We now ask a natural geometric question: given a circle and a straight line, when does the line \emph{just touch} the circle, and how do we write the equation of such a line? This is the theory of \cle{tangents}, one of the most frequently examined topics in the GCE Advanced Level paper.

\courssub{What is a tangent?}

\begin{definition}[Tangent to a circle]
A \cle{tangent} to a circle is a straight line which meets the circle at \emph{exactly one point}. This point is called the \cle{point of contact}. A line which cuts the circle in two distinct points is called a \emph{secant}; a line which does not meet the circle at all is an \emph{exterior} line.
\end{definition}

Imagine a secant line rotating about one of its intersection points: as the second intersection point slides along the circle and merges with the first, the secant becomes a tangent. The tangent is therefore the \emph{limiting position} of a secant — this is exactly the idea behind differentiation, and it explains why the tangent gives the direction of the curve at the point of contact.

\begin{propriete}[Radius--tangent perpendicularity]
The tangent at any point $T$ of a circle is perpendicular to the radius through the point of contact. That is, if $C$ is the centre, then $CT \perp$ tangent at $T$.
\end{propriete}

\textbf{Proof (recalled from O-Level).} Let $\ell$ be the tangent at $T$ and suppose $CT$ is \emph{not} perpendicular to $\ell$. Drop the perpendicular from $C$ to $\ell$, meeting $\ell$ at $Q \neq T$. Then $CQ < CT$ (perpendicular is the shortest distance from a point to a line), so $CQ < r$, which means $Q$ lies \emph{inside} the circle. But then the line $\ell$ passes through an interior point and must cut the circle twice — contradicting the fact that $\ell$ is a tangent. Hence $CT \perp \ell$. $\blacksquare$

\begin{popfigure}
\begin{center}
\begin{tikzpicture}[scale=1.1]
\coordinate (C) at (0,0);
\coordinate (T) at (1.5,2);
\draw[popblue, thick] (C) circle (2.5);
\draw[poporange, very thick] (C) -- (T);
\draw[popdark, thick] ($(T)!1.6cm!90:(C)$) -- ($(T)!1.6cm!-90:(C)$);
\draw[popgreen] ($(T)!0.28cm!(C)$) -- ($(T)!0.28cm!(C)!0.28cm!90:(T)$) -- ($(T)!0.28cm!90:(C)$);
\fill[popdark] (C) circle (1.5pt) node[below left] {$C$};
\fill[popdark] (T) circle (1.5pt) node[above right] {$T$};
\node[poporange] at (0.55,1.15) {$r$};
\node[popdark] at (3.2,1.1) {tangent};
\end{tikzpicture}
\end{center}
\legende{The tangent at $T$ is perpendicular to the radius $CT$.}
\end{popfigure}

\courssub{Tangent to $x^{2}+y^{2}=r^{2}$ at a point on the circle}

\begin{propriete}[Tangent formula — centre at the origin]
The equation of the tangent to the circle $x^{2}+y^{2}=r^{2}$ at the point $P(x_{1},y_{1})$ \textbf{on the circle} is
\[ xx_{1}+yy_{1}=r^{2}. \]
The proof is examinable and is given below.
\end{propriete}

\textbf{Derivation (gradient method).} The centre is $O(0,0)$, so the radius $OP$ has gradient $\dfrac{y_{1}}{x_{1}}$ (assuming $x_{1}\neq 0$). By the perpendicularity theorem, the tangent has gradient $-\dfrac{x_{1}}{y_{1}}$. Using point--slope form through $P(x_{1},y_{1})$:
\begin{align*}
y-y_{1} &= -\frac{x_{1}}{y_{1}}(x-x_{1})\\
y_{1}y - y_{1}^{2} &= -x_{1}x + x_{1}^{2}\\
xx_{1}+yy_{1} &= x_{1}^{2}+y_{1}^{2}.
\end{align*}
Since $P$ lies on the circle, $x_{1}^{2}+y_{1}^{2}=r^{2}$, giving $xx_{1}+yy_{1}=r^{2}$. (The cases $x_{1}=0$ or $y_{1}=0$ give horizontal/vertical tangents, which the formula also covers.) $\blacksquare$

\begin{attention}[Check the point lies on the circle!]
The formula $xx_{1}+yy_{1}=r^{2}$ is valid \textbf{only when $(x_{1},y_{1})$ lies on the circle}. If the point is outside or inside, the formula gives a different line (the \emph{polar}), not a tangent. Always verify $x_{1}^{2}+y_{1}^{2}=r^{2}$ before applying it.
\end{attention}

\begin{exemplebox}[Tangent at a given point]
Find the tangent to $x^{2}+y^{2}=25$ at $P(3,4)$.

\emph{Solution.} Check: $3^{2}+4^{2}=25$ \checkmark, so $P$ is on the circle. The tangent is
\[ 3x+4y=25. \]
\end{exemplebox}

\courssub{Tangent to $(x-a)^{2}+(y-b)^{2}=r^{2}$}

\begin{propriete}[Tangent formula — general centre]
The tangent to $(x-a)^{2}+(y-b)^{2}=r^{2}$ at the point $(x_{1},y_{1})$ on the circle is
\[ (x-a)(x_{1}-a)+(y-b)(y_{1}-b)=r^{2}. \]
\end{propriete}

This follows by translating the origin to $(a,b)$, or by the memorable \cle{replace one factor} rule: in each squared bracket, keep one factor as $(x-a)$ and replace the other by $(x_{1}-a)$; similarly for $y$. The constant $r^{2}$ is unchanged.

\begin{exoresolu}[Tangent to a shifted circle]
Find the equation of the tangent to the circle $(x-2)^{2}+(y+1)^{2}=10$ at the point $(5,0)$.
\tcblower
\textbf{Solution.} Check the point: $(5-2)^{2}+(0+1)^{2}=9+1=10$ \checkmark.

Apply the replace-one-factor rule with $(a,b)=(2,-1)$, $(x_{1},y_{1})=(5,0)$:
\begin{align*}
(x-2)(5-2)+(y+1)(0+1) &= 10\\
3(x-2)+(y+1) &= 10\\
3x+y &= 15.
\end{align*}
\emph{Verification by gradient:} centre $C(2,-1)$, radius gradient $=\frac{0-(-1)}{5-2}=\frac{1}{3}$, so tangent gradient $=-3$, and $y-0=-3(x-5)$ gives $3x+y=15$. \checkmark
\end{exoresolu}

\courssub{Tangent to the general circle equation}

\begin{propriete}[Tangent to $x^{2}+y^{2}+2gx+2fy+c=0$ at $(x_{1},y_{1})$]
The equation of the tangent to the circle $x^{2}+y^{2}+2gx+2fy+c=0$ at the point $(x_{1},y_{1})$ on the circle is
\[ xx_{1}+yy_{1}+g(x+x_{1})+f(y+y_{1})+c=0. \]
\end{propriete}

\textbf{Derivation (replace rule).} Starting from the circle equation, apply the following substitutions (the \cle{replace rule} for the general form):
\[ x^{2}\to xx_{1},\quad y^{2}\to yy_{1},\quad 2x\to x+x_{1},\quad 2y\to y+y_{1},\quad c\to c. \]
This yields $xx_{1}+yy_{1}+g(x+x_{1})+f(y+y_{1})+c=0$. One can verify that this line passes through $(x_{1},y_{1})$ and is perpendicular to the radius joining the centre $(-g,-f)$ to $(x_{1},y_{1})$.

\begin{exemplebox}[Tangent to a circle in general form]
Find the tangent to $x^{2}+y^{2}-6x+4y-12=0$ at the point $(7,1)$.

\emph{Solution.} Check the point: $7^{2}+1^{2}-42+4-12=49+1-42+4-12=0$ \checkmark.
Here $g=-3$, $f=2$, $c=-12$. The tangent is
\begin{align*}
7x+1y-3(x+7)+2(y+1)-12 &= 0\\
7x+y-3x-21+2y+2-12 &= 0\\
4x+3y-31 &= 0.
\end{align*}
\end{exemplebox}

\courssub{Lesson 8: Condition for a line to touch a circle}

\begin{propriete}[Tangency condition for $y=mx+c$ and $x^{2}+y^{2}=r^{2}$]
The line $y=mx+c$ is a tangent to the circle $x^{2}+y^{2}=r^{2}$ if and only if
\[ c^{2}=r^{2}(1+m^{2}). \]
Both derivations below are examinable.
\end{propriete}

\textbf{Derivation 1 — equal roots (discriminant).} Substitute $y=mx+c$ into $x^{2}+y^{2}=r^{2}$:
\begin{align*}
x^{2}+(mx+c)^{2} &= r^{2}\\
(1+m^{2})x^{2}+2mcx+(c^{2}-r^{2}) &= 0.
\end{align*}
The line touches the circle exactly when this quadratic in $x$ has \emph{equal roots}, i.e.\ $\Delta=0$:
\begin{align*}
(2mc)^{2}-4(1+m^{2})(c^{2}-r^{2}) &= 0\\
m^{2}c^{2}-(c^{2}-r^{2}+m^{2}c^{2}-m^{2}r^{2}) &= 0\\
r^{2}(1+m^{2})-c^{2} &= 0 \quad\Longrightarrow\quad c^{2}=r^{2}(1+m^{2}). \quad\blacksquare
\end{align*}

\textbf{Derivation 2 — perpendicular distance.} The line touches the circle exactly when the perpendicular distance from the centre $(0,0)$ to the line $mx-y+c=0$ equals the radius:
\[ \frac{|c|}{\sqrt{m^{2}+1}}=r \quad\Longrightarrow\quad c^{2}=r^{2}(1+m^{2}). \quad\blacksquare \]

\begin{propriete}[Tangents with given slope to $x^{2}+y^{2}=r^{2}$]
The two tangents of gradient $m$ to the circle $x^{2}+y^{2}=r^{2}$ are
\[ y = mx \pm r\sqrt{1+m^{2}}. \]
\end{propriete}

\begin{propriete}[General tangency condition]
The line $lx+my+n=0$ touches the circle with centre $(a,b)$ and radius $r$ if and only if
\[ \frac{|la+mb+n|}{\sqrt{l^{2}+m^{2}}}=r. \]
\end{propriete}

\begin{methode}[Two ways to test or impose tangency]
\begin{enumerate}
\item \textbf{Discriminant method:} substitute the line into the circle equation; the resulting quadratic must satisfy $\Delta=0$.
\item \textbf{Distance method:} the perpendicular distance from the centre to the line must equal the radius.
\end{enumerate}
The distance method is usually faster; the discriminant method also locates the point of contact (the repeated root).
\end{methode}

\begin{popfigure}
\begin{center}
\begin{tikzpicture}
\begin{axis}[
  axis equal image,
  axis lines=middle,
  xlabel={$x$}, ylabel={$y$},
  xmin=-8, xmax=8, ymin=-8, ymax=10,
  width=9cm, height=8cm,
  xtick={-5,5}, ytick={-5,5},
  tick label style={font=\small},
]
\addplot[popblue, thick, domain=0:360, samples=100] ({5*cos(x)},{5*sin(x)});
\addplot[poporange, very thick, domain=-8:6.5] {0.75*x+6.25};
\addplot[popdark, dashed, domain=-8:8] {0.75*x-6.25};
\addplot[only marks, mark=*, popdark, mark size=1.5pt] coordinates {(-3,4)};
\node[popdark, anchor=south west] at (axis cs:-3,4) {\small $T$};
\node[poporange, anchor=west] at (axis cs:2.2,7.9) {\small $y=\frac{3}{4}x+\frac{25}{4}$};
\node[popmuted, anchor=north west] at (axis cs:-7.5,-5.6) {\small $y=\frac{3}{4}x-\frac{25}{4}$};
\end{axis}
\end{tikzpicture}
\end{center}
\legende{The circle $x^{2}+y^{2}=25$ with the two tangents of gradient $m=\frac{3}{4}$: $c=\pm 5\sqrt{1+\frac{9}{16}}=\pm\frac{25}{4}$, illustrating $c^{2}=25(1+m^{2})$.}
\end{popfigure}

\begin{exoresolu}[Applying the tangency condition]
Show that the line $3x-4y+20=0$ is a tangent to the circle $x^{2}+y^{2}=16$, and find the point of contact.
\tcblower
\textbf{Solution.} Distance from centre $(0,0)$ to the line:
\[ d=\frac{|3(0)-4(0)+20|}{\sqrt{3^{2}+(-4)^{2}}}=\frac{20}{5}=4=r. \]
Since $d=r$, the line touches the circle. \checkmark

For the point of contact, write $y=\frac{3x+20}{4}$ and substitute into $x^{2}+y^{2}=16$:
\begin{align*}
x^{2}+\frac{(3x+20)^{2}}{16} &= 16\\
16x^{2}+9x^{2}+120x+400 &= 256\\
25x^{2}+120x+144 &= 0\\
(5x+12)^{2} &= 0 \quad\Longrightarrow\quad x=-\frac{12}{5}.
\end{align*}
Then $y=\dfrac{3(-\frac{12}{5})+20}{4}=\dfrac{16}{5}$. The point of contact is $\left(-\dfrac{12}{5},\dfrac{16}{5}\right)$. Note the repeated root — the algebraic signature of tangency.
\end{exoresolu}

\courssub{Tangents from an external point}

From a point $P$ \emph{outside} a circle, exactly \textbf{two} tangents can be drawn to the circle. To find their equations, write a general line through $P$ with gradient $m$, and impose the tangency condition (distance from centre $=$ radius). The resulting equation in $m$ is quadratic, giving the two gradients.

\begin{popfigure}
\begin{center}
\begin{tikzpicture}[scale=1.0]
\coordinate (C) at (0,0);
\coordinate (P) at (5.5,0);
\coordinate (T1) at (1.636,1.909);
\coordinate (T2) at (1.636,-1.909);
\draw[popblue, thick] (C) circle (2.5);
\draw[poporange, thick] (P) -- (T1);
\draw[poporange, thick] (P) -- (T2);
\draw[popgreen, dashed] (C) -- (T1);
\draw[popgreen, dashed] (C) -- (T2);
\draw[popmuted, dashed] (C) -- (P);
\fill[popdark] (C) circle (1.5pt) node[below left] {$C$};
\fill[popdark] (P) circle (1.5pt) node[right] {$P$};
\fill[popdark] (T1) circle (1.5pt) node[above left] {$T_{1}$};
\fill[popdark] (T2) circle (1.5pt) node[below left] {$T_{2}$};
\end{tikzpicture}
\end{center}
\legende{From an external point $P$ there are exactly two tangents, touching the circle at $T_{1}$ and $T_{2}$; $CT_{1}\perp PT_{1}$ and $CT_{2}\perp PT_{2}$.}
\end{popfigure}

\begin{exoresolu}[Tangents from an external point]
Find the equations of the tangents from the point $P(5,5)$ to the circle $x^{2}+y^{2}=10$.
\tcblower
\textbf{Solution.} First check $P$ is external: $5^{2}+5^{2}=50>10$ \checkmark.

A line through $(5,5)$ with gradient $m$: $y-5=m(x-5)$, i.e.\ $mx-y+(5-5m)=0$.

Distance from the centre $(0,0)$ must equal $r=\sqrt{10}$:
\begin{align*}
\frac{|5-5m|}{\sqrt{m^{2}+1}} &= \sqrt{10}\\
25(1-m)^{2} &= 10(m^{2}+1)\\
25-50m+25m^{2} &= 10m^{2}+10\\
15m^{2}-50m+15 &= 0\\
3m^{2}-10m+3 &= 0\\
(3m-1)(m-3) &= 0 \quad\Longrightarrow\quad m=\frac{1}{3}\ \text{or}\ m=3.
\end{align*}
The two tangents are:
\[ y-5=\frac{1}{3}(x-5)\ \Longrightarrow\ x-3y+10=0, \qquad y-5=3(x-5)\ \Longrightarrow\ 3x-y-10=0. \]
\end{exoresolu}

\begin{attention}[Do not lose a tangent!]
From an external point there are \textbf{exactly two} tangents. If your equation in $m$ yields only one finite value, the second tangent is \emph{vertical} (infinite gradient): test the line $x=x_{P}$ separately. Also, squaring both sides of the distance equation can introduce no extra solutions here, but always keep both roots of the quadratic in $m$.
\end{attention}

\begin{aretenir}[Key results of this section]
\begin{itemize}
\item The tangent at $T$ is perpendicular to the radius $CT$.
\item Tangent to $x^{2}+y^{2}=r^{2}$ at $(x_{1},y_{1})$ on the circle: $xx_{1}+yy_{1}=r^{2}$.
\item Tangent to $(x-a)^{2}+(y-b)^{2}=r^{2}$ at $(x_{1},y_{1})$: $(x-a)(x_{1}-a)+(y-b)(y_{1}-b)=r^{2}$.
\item Tangent to $x^{2}+y^{2}+2gx+2fy+c=0$ at $(x_{1},y_{1})$: $xx_{1}+yy_{1}+g(x+x_{1})+f(y+y_{1})+c=0$.
\item $y=mx+c$ touches $x^{2}+y^{2}=r^{2} \iff c^{2}=r^{2}(1+m^{2})$; the tangents are $y=mx\pm r\sqrt{1+m^{2}}$.
\item General test: $\dfrac{|la+mb+n|}{\sqrt{l^{2}+m^{2}}}=r$, or discriminant $=0$ after substitution.
\item From an external point: exactly two tangents.
\end{itemize}
\end{aretenir}

\begin{exercice}[Tangents]
\begin{enumerate}
\item Find the tangent to $x^{2}+y^{2}=29$ at the point $(2,5)$.
\item Find the tangent to $(x-1)^{2}+(y-3)^{2}=25$ at the point $(4,7)$.
\item Find the tangent to $x^{2}+y^{2}-6x+4y-12=0$ at the point $(7,1)$.
\item Show that $y=2x+10$ touches $x^{2}+y^{2}=20$ and find the point of contact.
\item Find the values of $k$ for which $y=x+k$ is a tangent to $x^{2}+y^{2}=8$.
\item Find the equations of the tangents from $(0,5)$ to $x^{2}+y^{2}=9$.
\item Find the equations of the tangents of gradient $2$ to the circle $x^{2}+y^{2}-4x+6y-12=0$.
\end{enumerate}
\end{exercice}

\begin{corrige}[Exercise 1]
\begin{enumerate}
\item $2^{2}+5^{2}=29$ \checkmark; tangent: $2x+5y=29$.
\item $(4-1)^{2}+(7-3)^{2}=25$ \checkmark; $(x-1)(3)+(y-3)(4)=25$ gives $3x+4y=40$.
\item Circle: $g=-3$, $f=2$, $c=-12$. Point $(7,1)$ satisfies the equation. Tangent: $7x+1y-3(x+7)+2(y+1)-12=0 \Rightarrow 4x+3y-31=0$.
\item Distance from $(0,0)$: $\frac{|10|}{\sqrt{5}}=2\sqrt{5}=\sqrt{20}=r$ \checkmark. Substitution gives $x^{2}+(2x+10)^{2}=20$, i.e.\ $5x^{2}+40x+80=0$, $(x+4)^{2}=0$; contact point $(-4,2)$.
\item $c^{2}=r^{2}(1+m^{2})$: $k^{2}=8(2)=16$, so $k=\pm 4$.
\item Line $y=mx+5$: distance $\frac{5}{\sqrt{m^{2}+1}}=3$ gives $m^{2}=\frac{16}{9}$, $m=\pm\frac{4}{3}$. Tangents: $4x-3y+15=0$ and $4x+3y-15=0$.
\item Centre $(2,-3)$, radius $r=\sqrt{4+9+12}=5$. Tangents with gradient $2$: $y=2x+c$. Distance from centre: $\frac{|2(2)-(-3)+c|}{\sqrt{5}}=\frac{|7+c|}{\sqrt{5}}=5 \Rightarrow |c+7|=5\sqrt{5}$. So $c=-7\pm5\sqrt{5}$. Tangents: $y=2x-7+5\sqrt{5}$ and $y=2x-7-5\sqrt{5}$.
\end{enumerate}
\end{corrige}

\begin{savaistu}[Tangents in everyday life]
When a wheel rolls along a road in Douala, the road is \emph{tangent} to the wheel at the point of contact — and the radius to that point is vertical, perpendicular to the road. The same principle is used by engineers designing gear teeth and cam profiles: smooth contact between two curves always happens along a common tangent.
\end{savaistu}

\coursec{Section 4: Parametric Equations and Length of Tangent}

In Section 3 we described a circle by a single Cartesian equation and studied its tangents. That equation, however, mixes $x$ and $y$ together: it tells us \emph{which} points lie on the circle, but not \emph{where} a moving point on the circle is at a given moment. Imagine a child on the merry-go-round at the amusement park in Douala: to describe her position at time $t$, the most natural quantity is the \emph{angle} through which the ride has turned. This is exactly the idea of a \cle{parametric equation}: instead of relating $x$ and $y$ directly, we express both coordinates in terms of a single auxiliary variable, the \cle{parameter}. In this section we build the parametric equations of a circle, learn how to pass from the parametric form to the Cartesian form and back, derive the tangent at a point of parameter $\theta$, and finish with a beautiful and powerful result: the length of the tangent drawn to a circle from an external point.

\courssub{Parametric Equations of a Circle}

\textbf{Observation: describing a point on a circle by an angle.}\par
Take a circle of radius $r$ centred at the origin $O$, and let $P(x, y)$ be any point on it. Join $OP$ and let $\theta$ be the angle that $OP$ makes with the positive $x$-axis, measured anticlockwise. Dropping a perpendicular from $P$ to the $x$-axis creates a right-angled triangle with hypotenuse $OP = r$, so by the very definitions of cosine and sine,
\[
\cos\theta = \frac{x}{r} \qquad\text{and}\qquad \sin\theta = \frac{y}{r}.
\]
Hence $x = r\cos\theta$ and $y = r\sin\theta$. As $\theta$ varies from $0$ to $2\pi$, the point $P$ travels once round the whole circle. The angle $\theta$ is the parameter.

\begin{definition}[Parametric equations of a circle]
The \cle{parametric equations} of the circle with centre $C(a, b)$ and radius $r$ are
\[
x = a + r\cos\theta, \qquad y = b + r\sin\theta, \qquad 0 \le \theta < 2\pi,
\]
where $\theta$ is the angle, measured anticlockwise from the positive $x$-direction, that the radius $CP$ makes with the horizontal. In particular, for a circle centred at the origin,
\[
x = r\cos\theta, \qquad y = r\sin\theta.
\]
The point $(a + r\cos\theta,\ b + r\sin\theta)$ is often called ``the point $\theta$'' on the circle.
\end{definition}

The derivation for a general centre is a simple translation: if $P(x,y)$ lies on the circle of centre $C(a,b)$ and radius $r$, then the vector $\overrightarrow{CP}$ has length $r$ and makes angle $\theta$ with the positive $x$-direction, so $\overrightarrow{CP} = (r\cos\theta,\ r\sin\theta)$, giving $P = C + \overrightarrow{CP} = (a + r\cos\theta,\ b + r\sin\theta)$.

\begin{popfigure}
\begin{center}
\begin{tikzpicture}[scale=1.5]
\draw[popline,->] (-1.2,0) -- (4.6,0) node[below left] {$x$};
\draw[popline,->] (0,-1.2) -- (0,3.4) node[below left] {$y$};
\coordinate (C) at (2,1.3);
\draw[popblue,thick] (C) circle (1.3);
\fill[popdark] (C) circle (0.045) node[below left] {$C(a,b)$};
\coordinate (P) at ($(C)+(40:1.3)$);
\fill[poporange] (P) circle (0.05) node[above right, text=popdark] {$P(a+r\cos\theta,\ b+r\sin\theta)$};
\draw[popdark,thick] (C) -- (P) node[midway, above left] {$r$};
\draw[popmuted,dashed] (C) -- ($(C)+(1.9,0)$);
\draw[popgreen,thick] ($(C)+(0.55,0)$) arc (0:40:0.55);
\node[text=popgreen] at ($(C)+(20:0.85)$) {$\theta$};
\fill[popdark] (0,0) circle (0.04) node[below left] {$O$};
\end{tikzpicture}
\end{center}
\legende{The parameter $\theta$ locates the point $P$ on the circle of centre $C(a,b)$ and radius $r$.}
\end{popfigure}

\textbf{Geometric meaning of the parameter.}\par
The parameter $\theta$ is \emph{not} the angle at the point $P$; it is the angle at the \emph{centre}, between the positive $x$-direction and the radius $CP$. Particular values are worth memorising:
\begin{center}
\begin{tabular}{|c|c|c|c|c|c|}
\hline
\rowcolor{popblueL} $\theta$ & $0$ & $\dfrac{\pi}{2}$ & $\pi$ & $\dfrac{3\pi}{2}$ & $\dfrac{\pi}{4}$ \\ \hline
Point on $x^2+y^2=r^2$ & $(r,0)$ & $(0,r)$ & $(-r,0)$ & $(0,-r)$ & $\left(\dfrac{r}{\sqrt2},\dfrac{r}{\sqrt2}\right)$ \\ \hline
\end{tabular}
\end{center}

\begin{exemplebox}[Locating points for given parameter values]
The circle $x^2 + y^2 = 25$ has parametric equations $x = 5\cos\theta$, $y = 5\sin\theta$.
\begin{itemize}
\item For $\theta = \dfrac{\pi}{3}$: $x = 5\cos\dfrac{\pi}{3} = \dfrac{5}{2}$, $y = 5\sin\dfrac{\pi}{3} = \dfrac{5\sqrt3}{2}$. The point is $\left(\dfrac{5}{2}, \dfrac{5\sqrt3}{2}\right)$.
\item For $\theta = \pi$: the point is $(-5, 0)$.
\end{itemize}
Conversely, the point $(3, 4)$ lies on this circle since $3^2 + 4^2 = 25$; its parameter satisfies $\cos\theta = \dfrac{3}{5}$, $\sin\theta = \dfrac{4}{5}$, so $\theta = \tan^{-1}\dfrac{4}{3} \approx 0.927$ rad (first quadrant, since both coordinates are positive).
\end{exemplebox}

\begin{attention}[The parameter is measured from the centre]
A frequent error is to measure $\theta$ from the origin when the centre is $(a,b) \ne (0,0)$. The angle $\theta$ is always the angle of the \textbf{radius from the centre to the point}, measured from the positive $x$-direction. Also, do not confuse the parameter $\theta$ with the polar angle of $P$ about the origin: they coincide only when the circle is centred at $O$.
\end{attention}

\begin{exemplebox}[Finding the parameter of a given point on a translated circle]
Find the parameter $\theta$ of the point $P(5, 7)$ on the circle $(x-2)^2 + (y-3)^2 = 25$.
\tcblower
The centre is $C(2,3)$ and $r=5$. We have
\[
\cos\theta = \frac{5-2}{5} = \frac{3}{5}, \qquad \sin\theta = \frac{7-3}{5} = \frac{4}{5}.
\]
Since both are positive, $\theta$ is in the first quadrant: $\theta = \tan^{-1}\frac{4}{3} \approx 0.927$ rad.
\end{exemplebox}

\courssub{Conversion Between Parametric and Cartesian Forms}

Given parametric equations, how do we recover the Cartesian equation? The key is the Pythagorean identity $\cos^2\theta + \sin^2\theta = 1$, which destroys the parameter.

\begin{methode}[Eliminating the parameter $\theta$]
To convert the parametric equations of a circle to Cartesian form:
\begin{enumerate}
\item Isolate $\cos\theta$ and $\sin\theta$: from $x = a + r\cos\theta$ get $\cos\theta = \dfrac{x-a}{r}$; from $y = b + r\sin\theta$ get $\sin\theta = \dfrac{y-b}{r}$.
\item Substitute into the identity $\cos^2\theta + \sin^2\theta = 1$:
\[
\left(\frac{x-a}{r}\right)^2 + \left(\frac{y-b}{r}\right)^2 = 1.
\]
\item Multiply through by $r^2$: $(x-a)^2 + (y-b)^2 = r^2$.
\end{enumerate}
Conversely, given a Cartesian circle $(x-a)^2 + (y-b)^2 = r^2$, set $\dfrac{x-a}{r} = \cos\theta$ and $\dfrac{y-b}{r} = \sin\theta$ to obtain the parametric form.
\end{methode}

\begin{exoresolu}[From parametric to Cartesian and back]
\textbf{(a)} A curve is given by $x = 2 + 3\cos\theta$, $y = -1 + 3\sin\theta$. Show that it is a circle and state its centre and radius.
\textbf{(b)} Write down parametric equations for the circle $x^2 + y^2 - 4x + 6y - 12 = 0$.
\tcblower
\textbf{(a)} Isolate the trigonometric functions:
\[
\cos\theta = \frac{x-2}{3}, \qquad \sin\theta = \frac{y+1}{3}.
\]
Applying $\cos^2\theta + \sin^2\theta = 1$:
\[
\frac{(x-2)^2}{9} + \frac{(y+1)^2}{9} = 1
\quad\Longrightarrow\quad
(x-2)^2 + (y+1)^2 = 9.
\]
This is a circle of centre $(2, -1)$ and radius $3$.

\textbf{(b)} Complete the squares:
\[
x^2 - 4x + y^2 + 6y = 12
\;\Longrightarrow\;
(x-2)^2 - 4 + (y+3)^2 - 9 = 12
\;\Longrightarrow\;
(x-2)^2 + (y+3)^2 = 25.
\]
Hence the centre is $(2, -3)$ and $r = 5$, so parametric equations are
\[
x = 2 + 5\cos\theta, \qquad y = -3 + 5\sin\theta, \qquad 0 \le \theta < 2\pi.
\]
\end{exoresolu}

\begin{exemplebox}[Conversion with a non-standard parameter range]
The curve $x = 1 + 2\cos 2t$, $y = -2 + 2\sin 2t$ ($0 \le t < \pi$) is also a circle. Let $\theta = 2t$; then $x = 1 + 2\cos\theta$, $y = -2 + 2\sin\theta$ with $0 \le \theta < 2\pi$. Centre $(1,-2)$, radius $2$.
\end{exemplebox}

\courssub{Tangent at the Point with Parameter $\theta$}

Parametric form gives a very quick route to the tangent. Consider the circle $x^2 + y^2 = r^2$ and the point $P(r\cos\theta, r\sin\theta)$ on it. The radius $OP$ has gradient $\dfrac{r\sin\theta}{r\cos\theta} = \tan\theta$, so the tangent at $P$, being perpendicular to $OP$, has gradient $-\cot\theta = -\dfrac{\cos\theta}{\sin\theta}$. Its equation is
\[
y - r\sin\theta = -\frac{\cos\theta}{\sin\theta}\left(x - r\cos\theta\right).
\]
Multiplying by $\sin\theta$:
\[
y\sin\theta - r\sin^2\theta = -x\cos\theta + r\cos^2\theta
\quad\Longrightarrow\quad
x\cos\theta + y\sin\theta = r\left(\cos^2\theta + \sin^2\theta\right) = r.
\]

\begin{propriete}[Tangent at the point $\theta$]
The equation of the tangent to the circle $x^2 + y^2 = r^2$ at the point $(r\cos\theta, r\sin\theta)$ is
\[
x\cos\theta + y\sin\theta = r.
\]
More generally, the tangent to $(x-a)^2 + (y-b)^2 = r^2$ at the point $(a + r\cos\theta,\ b + r\sin\theta)$ is
\[
(x-a)\cos\theta + (y-b)\sin\theta = r.
\]
The derivation above (gradient of radius, perpendicular gradient, point--gradient form) is examinable at GCE Advanced Level.
\end{propriete}

\begin{exemplebox}[Tangent at a given point]
Find the tangent to $x^2 + y^2 = 25$ at the point $\theta = \dfrac{\pi}{3}$.

Here $r = 5$, $\cos\dfrac{\pi}{3} = \dfrac12$, $\sin\dfrac{\pi}{3} = \dfrac{\sqrt3}{2}$. The tangent is
\[
\frac{1}{2}x + \frac{\sqrt3}{2}y = 5
\quad\Longrightarrow\quad
x + \sqrt3\,y = 10.
\]
Check: the point of contact $\left(\dfrac52, \dfrac{5\sqrt3}{2}\right)$ satisfies the equation, and the distance from $O$ to the line is $\dfrac{10}{\sqrt{1+3}} = 5 = r$, as expected for a tangent.
\end{exemplebox}

\begin{exemplebox}[Tangent to a translated circle using the parameter]
Find the equation of the tangent to the circle $(x-3)^2 + (y+2)^2 = 16$ at the point where $\theta = \dfrac{2\pi}{3}$.
\tcblower
Centre $C(3,-2)$, $r=4$. $\cos\frac{2\pi}{3} = -\frac12$, $\sin\frac{2\pi}{3} = \frac{\sqrt3}{2}$.
Point of contact: $x = 3 + 4(-\frac12) = 1$, $y = -2 + 4(\frac{\sqrt3}{2}) = -2 + 2\sqrt3$.
Tangent: $(x-3)\cos\theta + (y+2)\sin\theta = 4$
$\Rightarrow (x-3)(-\frac12) + (y+2)(\frac{\sqrt3}{2}) = 4$
$\Rightarrow -(x-3) + \sqrt3(y+2) = 8$
$\Rightarrow -x + 3 + \sqrt3 y + 2\sqrt3 = 8$
$\Rightarrow x - \sqrt3 y = 2\sqrt3 - 5$.
\end{exemplebox}

\begin{attention}[Parameter vs. gradient]
The gradient of the tangent is $-\cot\theta$, \emph{not} $-\tan\theta$. The radius has gradient $\tan\theta$; the tangent is perpendicular, so its gradient is the negative reciprocal. A common GCE error is to write the tangent gradient as $-\tan\theta$.
\end{attention}

\courssub{Length of the Tangent from an External Point}

\textbf{Observation.}\par
Stand outside a circular roundabout and stretch a rope so that it just touches the roundabout's edge: the rope is a tangent, and the point of contact $T$ satisfies $CT \perp PT$, where $C$ is the centre. The triangle $PCT$ is therefore right-angled at $T$, and Pythagoras does the rest.

\begin{popfigure}
\begin{center}
\begin{tikzpicture}[scale=1.3]
\coordinate (C) at (0,0);
\draw[popblue,thick] (C) circle (1.4);
\fill[popdark] (C) circle (0.05) node[below left] {$C$};
\coordinate (P) at (3.6,1.6);
\fill[poporange] (P) circle (0.06) node[right, text=popdark] {$P(x_1,y_1)$};
\coordinate (T) at (0.42,1.34);
\fill[popgreen] (T) circle (0.05) node[above left, text=popdark] {$T$};
\draw[popdark,thick] (C) -- (P) node[midway, below right] {$d = CP$};
\draw[popgreen,thick] (C) -- (T) node[midway, left] {$r$};
\draw[poporange,thick] (P) -- (T) node[midway, above] {$L$};
\draw[popmuted] ($(T)+(0.13,-0.10)$) -- ($(T)+(0.27,0.09)$) -- ($(T)+(0.14,0.23)$);
\node[text=popmuted] at (1.1,1.05) {};
\end{tikzpicture}
\end{center}
\legende{The tangent segment $PT$, the radius $CT$ and the hypotenuse $CP$ form a right-angled triangle.}
\end{popfigure}

\begin{propriete}[Length of tangent from an external point]
Let $P(x_1, y_1)$ be a point outside the circle with centre $C$ and radius $r$, and let $T$ be the point of contact of a tangent from $P$. Then
\[
L = PT = \sqrt{CP^2 - r^2}.
\]
If the circle is given in general form $S \equiv x^2 + y^2 + 2gx + 2fy + c = 0$, then
\[
\boxed{\;L = \sqrt{S_1}\;} \qquad\text{where}\qquad S_1 = x_1^2 + y_1^2 + 2gx_1 + 2fy_1 + c.
\]
That is: \emph{substitute the coordinates of $P$ into the left-hand side of the normalised equation and take the square root.}
\end{propriete}

\textbf{Proof.} Since $CT \perp PT$ (tangent perpendicular to radius), triangle $PCT$ is right-angled at $T$, so $PT^2 = CP^2 - CT^2 = CP^2 - r^2$. For the general circle, $C = (-g, -f)$ and $r^2 = g^2 + f^2 - c$, hence
\begin{align*}
L^2 &= (x_1 + g)^2 + (y_1 + f)^2 - (g^2 + f^2 - c)\\
&= x_1^2 + 2gx_1 + g^2 + y_1^2 + 2fy_1 + f^2 - g^2 - f^2 + c\\
&= x_1^2 + y_1^2 + 2gx_1 + 2fy_1 + c = S_1. \qquad\blacksquare
\end{align*}

\begin{attention}[Normalise before substituting!]
The formula $L = \sqrt{S_1}$ is valid \textbf{only} when the coefficients of $x^2$ and $y^2$ are both equal to $1$. If the circle is given as $2x^2 + 2y^2 + 4x - 6y + 3 = 0$, you must first divide through by $2$ before substituting the point. Substituting into the unnormalised equation gives $2L^2$, not $L^2$ — a classic GCE trap.
\end{attention}

\begin{exoresolu}[Length of a tangent from an external point]
Find the length of the tangent from the point $P(5, 7)$ to the circle $x^2 + y^2 - 4x - 6y + 9 = 0$.
\tcblower
The equation is already normalised (coefficients of $x^2$ and $y^2$ equal to $1$). Substitute $x_1 = 5$, $y_1 = 7$:
\begin{align*}
S_1 &= 5^2 + 7^2 - 4(5) - 6(7) + 9\\
&= 25 + 49 - 20 - 42 + 9 = 21.
\end{align*}
Hence $L = \sqrt{21} \approx 4.58$ units.

\emph{Verification by the geometric formula:} the centre is $C(2, 3)$ and $r^2 = 4 + 9 - 9 = 4$, so $r = 2$. Then $CP^2 = (5-2)^2 + (7-3)^2 = 9 + 16 = 25$, and $L = \sqrt{25 - 4} = \sqrt{21}$. \checkmark
\end{exoresolu}

\begin{propriete}[Power of a point (GCE extra)]
The quantity $S_1 = x_1^2 + y_1^2 + 2gx_1 + 2fy_1 + c$ is called the \cle{power of the point} $P$ with respect to the circle. Its sign locates $P$:
\begin{itemize}
\item $S_1 > 0 \iff P$ lies \textbf{outside} the circle, and $S_1 = L^2$ (square of the tangent length);
\item $S_1 = 0 \iff P$ lies \textbf{on} the circle;
\item $S_1 < 0 \iff P$ lies \textbf{inside} the circle (no real tangent exists from $P$).
\end{itemize}
This is not formally required, but it gives a one-line test for the position of a point and explains why the tangent-length formula only makes sense for external points.
\end{propriete}

\begin{exemplebox}[Testing the position of a point]
For the circle $x^2 + y^2 - 4x - 6y + 9 = 0$ of the previous exercise:
\begin{itemize}
\item At $(5,7)$: $S_1 = 21 > 0$, so $(5,7)$ is outside.
\item At $(2, 5)$: $S_1 = 4 + 25 - 8 - 30 + 9 = 0$, so $(2,5)$ is on the circle.
\item At $(2, 4)$: $S_1 = 4 + 16 - 8 - 24 + 9 = -3 < 0$, so $(2,4)$ is inside.
\end{itemize}
\end{exemplebox}

\begin{exemplebox}[Length of tangent with non-normalised equation]
Find the length of the tangent from $P(4, -1)$ to the circle $2x^2 + 2y^2 - 8x + 4y - 10 = 0$.
\tcblower
First normalise: divide by $2$: $x^2 + y^2 - 4x + 2y - 5 = 0$.
Substitute $(4, -1)$:
$S_1 = 16 + 1 - 16 - 2 - 5 = -6 < 0$.
The point is \emph{inside} the circle; no real tangent exists.
(If you forgot to divide by $2$, you would get $2S_1 = -12$ and might incorrectly write $L = \sqrt{-12}$.)
\end{exemplebox}

\begin{exemplebox}[Finding the point of contact given the length]
From $P(8, 0)$ a tangent of length $5$ is drawn to the circle $x^2 + y^2 = 9$. Find the coordinates of the point of contact $T$.
\tcblower
$r=3$, $L=5$, so $CP^2 = L^2 + r^2 = 25 + 9 = 34$. But $CP^2 = 8^2 + 0^2 = 64 \ne 34$ — contradiction!
Wait: $P(8,0)$ gives $CP=8$, so $L = \sqrt{64-9} = \sqrt{55} \ne 5$. The data is inconsistent.
\textbf{Correct problem:} From $P(5,0)$ a tangent of length $4$ is drawn to $x^2+y^2=9$. Find $T$.
$CP=5$, $r=3$, $L=4$ (consistent: $5^2=3^2+4^2$).
Let $T(3\cos\theta, 3\sin\theta)$. The tangent is $x\cos\theta + y\sin\theta = 3$.
Since $P(5,0)$ lies on it: $5\cos\theta = 3 \Rightarrow \cos\theta = 3/5$, $\sin\theta = \pm 4/5$.
Two points of contact: $T_1(9/5, 12/5)$ and $T_2(9/5, -12/5)$.
\end{exemplebox}

\begin{exercice}[Consolidation]
\begin{enumerate}
\item Write down parametric equations of the circle $(x+1)^2 + (y-4)^2 = 16$, and find the coordinates of the points with parameters $\theta = \dfrac{\pi}{6}$ and $\theta = \dfrac{5\pi}{6}$.
\item A curve is given by $x = 1 - 2\cos\theta$, $y = 3 + 2\sin\theta$. Show it is a circle and find its centre and radius.
\item Find the equation of the tangent to $x^2 + y^2 = 169$ at the point with parameter $\theta$ where $\cos\theta = \dfrac{5}{13}$, $\sin\theta = \dfrac{12}{13}$.
\item Find the length of the tangent from $(-2, 3)$ to the circle $x^2 + y^2 + 6x - 2y - 6 = 0$.
\item Find the length of the tangent from $(1, 1)$ to the circle $3x^2 + 3y^2 - 12x + 6y - 3 = 0$. \emph{(Beware the normalisation!)}
\item A circle has parametric equations $x = 4\cos\theta$, $y = 4\sin\theta$. Find the equation of the tangent at the point where $\theta = \dfrac{3\pi}{4}$.
\item The point $P(7, 2)$ lies outside the circle $x^2 + y^2 - 6x - 8y = 0$. Find the length of the tangent from $P$ to the circle.
\item Show that the point $(3, -4)$ lies on the circle $x^2 + y^2 - 2x + 4y - 20 = 0$ and find the equation of the tangent at this point using the parametric form.
\end{enumerate}
\end{exercice}

\begin{corrige}[Exercise 1]
\begin{enumerate}
\item $x = -1 + 4\cos\theta$, $y = 4 + 4\sin\theta$. For $\theta = \dfrac{\pi}{6}$: $\left(-1 + 2\sqrt3,\ 6\right)$. For $\theta = \dfrac{5\pi}{6}$: $\left(-1 - 2\sqrt3,\ 6\right)$.
\item $\cos\theta = \dfrac{1-x}{2}$, $\sin\theta = \dfrac{y-3}{2}$, so $(x-1)^2 + (y-3)^2 = 4$: centre $(1, 3)$, radius $2$.
\item Tangent: $x\cos\theta + y\sin\theta = r$ gives $\dfrac{5}{13}x + \dfrac{12}{13}y = 13$, i.e.\ $5x + 12y = 169$.
\item $S_1 = 4 + 9 - 12 - 6 - 6 = -11 < 0$: the point is \emph{inside} the circle, so no real tangent exists. (A good reminder to check the sign first!)
\item Normalise: divide by $3$: $x^2 + y^2 - 4x + 2y - 1 = 0$. Then $S_1 = 1 + 1 - 4 + 2 - 1 = -1 < 0$: again the point is inside the circle. (Substituting without dividing would have given $-3$, and even $\sqrt{-3}$ would be meaningless — the sign test protects you.)
\item $r=4$, $\cos\frac{3\pi}{4}=-\frac{\sqrt2}{2}$, $\sin\frac{3\pi}{4}=\frac{\sqrt2}{2}$. Tangent: $-\frac{\sqrt2}{2}x + \frac{\sqrt2}{2}y = 4 \Rightarrow -x + y = 4\sqrt2 \Rightarrow y = x + 4\sqrt2$.
\item Circle: $(x-3)^2 + (y-4)^2 = 25$, centre $(3,4)$, $r=5$. $CP^2 = (7-3)^2 + (2-4)^2 = 16+4=20$. $L = \sqrt{20-25}$ — wait, $20 < 25$! $P$ is \emph{inside} the circle. Check: $S_1 = 49+4-42-16 = -5 < 0$. No real tangent. (The problem statement said ``lies outside'' — it was a trap!)
\item Substitute $(3,-4)$: $9+16-6-16-20 = -17 \ne 0$. The point is \emph{not} on the circle. (Another trap!) Correct point on circle: e.g. $(6,2)$ gives $36+4-12+8-20=16\ne0$. Actually, centre $(1,-2)$, $r=5$. Point $(3,-4)$: $CP^2 = 4+4=8 < 25$, inside.
\end{enumerate}
\end{corrige}

\begin{aretenir}[Key points of Section 4]
\begin{itemize}
\item Parametric equations: $x = a + r\cos\theta$, $y = b + r\sin\theta$; $\theta$ is the angle at the \textbf{centre}, measured from the positive $x$-direction.
\item To eliminate the parameter: isolate $\cos\theta$ and $\sin\theta$, then use $\cos^2\theta + \sin^2\theta = 1$.
\item Tangent at the point $\theta$ on $x^2 + y^2 = r^2$: $x\cos\theta + y\sin\theta = r$.
\item Tangent at the point $\theta$ on $(x-a)^2 + (y-b)^2 = r^2$: $(x-a)\cos\theta + (y-b)\sin\theta = r$.
\item Length of tangent from $P(x_1,y_1)$: $L = \sqrt{CP^2 - r^2} = \sqrt{S_1}$, where $S_1$ is obtained by substituting $P$ into the \textbf{normalised} equation.
\item Sign of $S_1$: positive outside, zero on, negative inside the circle.
\item Always check $S_1 \ge 0$ before taking the square root; if $S_1 < 0$, the point is inside and no real tangent exists.
\end{itemize}
\end{aretenir}

\begin{savaistu}[Why parameters matter beyond the GCE]
Parametric descriptions are the language of motion. When engineers model the wheel of a truck on the Yaounde--Bafoussam road, or when a computer animates a rotating object in a video game, the position of each point is written as $(a + r\cos\theta(t),\ b + r\sin\theta(t))$ with $\theta$ depending on time $t$. The same idea, with two parameters, describes surfaces like the sphere — the starting point of 3D graphics and of the GPS system that guides taxis across Douala.
\end{savaistu}

\coursec{Section 5: Intersection of a Circle and a Line; Two Circles}

In Section 4 we studied the circle through its parametric equations and measured the length of a tangent drawn from an external point. We now change our point of view: instead of looking at one circle alone, we ask how a circle interacts with other objects in the plane. Two questions dominate the GCE Advanced Level examination:

\begin{itemize}
\item Where does a given \cle{straight line} meet a given \cle{circle}?
\item Where do \cle{two circles} meet each other, and when do they merely \emph{touch}?
\end{itemize}

Both questions reduce to algebra you already master: substitution and quadratic equations. The geometry then gives us beautiful shortcuts — the discriminant, the perpendicular distance from the centre, and the line of centres.

\courssub{Intersection of a Line and a Circle}

\textbf{Setting up the problem.}\par
Consider a circle of centre $C$ and radius $r$, and a line $\ell$. Imagine sliding the line slowly towards the circle, as when the horizon line of the sea seems to approach the round sun at sunset over the beach at Limbe. Three things can happen: the line \emph{cuts} the circle in two points (a \cle{secant}), it \emph{touches} the circle in exactly one point (a \cle{tangent}), or it \emph{misses} the circle completely. Algebraically, all three cases appear when we solve the two equations simultaneously.

\begin{methode}[Intersecting a line and a circle]
To find the points of intersection of a line $\ell$ and a circle $\mathscr{C}$:
\begin{enumerate}
\item From the equation of the line, express one variable in terms of the other (usually $y$ in terms of $x$).
\item Substitute this expression into the equation of the circle.
\item Simplify: you obtain a \textbf{quadratic equation} in one variable.
\item Solve the quadratic. Each root gives one coordinate; use the line equation to find the other coordinate.
\item Interpret the number of real roots: two roots give two intersection points, one (repeated) root gives a tangent point, no real root means the line misses the circle.
\end{enumerate}
\end{methode}

\begin{exoresolu}[Intersection of a line and a circle]
Find the points of intersection of the line $y = x - 1$ and the circle $x^{2} + y^{2} - 4x - 2y + 1 = 0$.
\tcblower
\textbf{Solution.} Substitute $y = x - 1$ into the circle equation:
\begin{align*}
x^{2} + (x-1)^{2} - 4x - 2(x-1) + 1 &= 0\\
x^{2} + x^{2} - 2x + 1 - 4x - 2x + 2 + 1 &= 0\\
2x^{2} - 8x + 4 &= 0\\
x^{2} - 4x + 2 &= 0.
\end{align*}
Hence $x = \dfrac{4 \pm \sqrt{16 - 8}}{2} = 2 \pm \sqrt{2}$. Since $y = x - 1$, the intersection points are
\[ \left(2 + \sqrt{2},\; 1 + \sqrt{2}\right) \quad \text{and} \quad \left(2 - \sqrt{2},\; 1 - \sqrt{2}\right). \]
Two distinct points: the line is a secant of the circle.
\end{exoresolu}

\courssub{The Discriminant Decides Everything}

When the substitution of the line into the circle produces the quadratic $ax^{2} + bx + c = 0$, the number of intersection points is governed by its discriminant $\Delta = b^{2} - 4ac$.

\begin{propriete}[Discriminant criterion for line--circle intersection]
Let $\Delta$ be the discriminant of the quadratic obtained by substituting the equation of a line into the equation of a circle. Then:
\begin{itemize}
\item $\Delta > 0$: the line cuts the circle in \textbf{two distinct points} (the line is a \cle{secant});
\item $\Delta = 0$: the line meets the circle in \textbf{exactly one point} (the line is a \cle{tangent});
\item $\Delta < 0$: the line and the circle have \textbf{no common point} (the line is exterior to the circle).
\end{itemize}
This criterion is examinable and is the fastest way to test whether a line is tangential to a circle.
\end{propriete}

\begin{popfigure}
\begin{center}
\begin{tikzpicture}[scale=0.85]
% Secant
\begin{scope}[shift={(0,0)}]
\draw[popline] (0,0) circle (1.3);
\draw[popblue, thick] (-1.9,0.5) -- (1.9,0.5);
\fill[poporange] (-1.2,0.5) circle (0.06) (1.2,0.5) circle (0.06);
\fill[popdark] (0,0) circle (0.05);
\node[below] at (0,-1.5) {\small $\Delta>0$: secant};
\end{scope}
% Tangent
\begin{scope}[shift={(5,0)}]
\draw[popline] (0,0) circle (1.3);
\draw[popgreen, thick] (-1.9,1.3) -- (1.9,1.3);
\fill[poporange] (0,1.3) circle (0.06);
\fill[popdark] (0,0) circle (0.05);
\node[below] at (0,-1.5) {\small $\Delta=0$: tangent};
\end{scope}
% Exterior
\begin{scope}[shift={(10,0)}]
\draw[popline] (0,0) circle (1.3);
\draw[poppurple, thick] (-1.9,2.0) -- (1.9,2.0);
\fill[popdark] (0,0) circle (0.05);
\node[below] at (0,-1.5) {\small $\Delta<0$: no intersection};
\end{scope}
\end{tikzpicture}
\end{center}
\legende{The three possible positions of a line relative to a circle, decided by the discriminant.}
\end{popfigure}

\begin{exemplebox}[Testing tangency with the discriminant]
Show that the line $y = 2x + 5$ is a tangent to the circle $x^{2} + y^{2} = 5$, and find the point of contact.

Substituting: $x^{2} + (2x+5)^{2} = 5$, so $5x^{2} + 20x + 20 = 0$, i.e. $x^{2} + 4x + 4 = 0$. Here $\Delta = 16 - 16 = 0$: the line is a tangent. The repeated root is $x = -2$, giving $y = 1$. The point of contact is $(-2, 1)$.
\end{exemplebox}

\courssub{The Distance Method and the Length of a Chord}

There is a second, purely geometric way to decide the position of a line: compare the \cle{perpendicular distance} $d$ from the centre $C$ to the line with the radius $r$. Recall that the distance from $(x_{0}, y_{0})$ to the line $ax + by + c = 0$ is
\[ d = \frac{|ax_{0} + by_{0} + c|}{\sqrt{a^{2} + b^{2}}}. \]

\begin{propriete}[Distance criterion]
Let $d$ be the perpendicular distance from the centre of a circle of radius $r$ to a line $\ell$. Then:
\begin{itemize}
\item $d < r$: the line is a secant (two intersection points);
\item $d = r$: the line is a tangent (one point of contact);
\item $d > r$: the line does not meet the circle.
\end{itemize}
\end{propriete}

When $d < r$, the line cuts off a \cle{chord}. Drop the perpendicular from the centre to the chord: a classical result (perpendicular from the centre bisects the chord) splits the chord into two equal parts. Pythagoras' theorem in the right triangle formed by the radius, half the chord and the perpendicular gives the following formula.

\begin{propriete}[Length of a chord]
If a line is at perpendicular distance $d$ from the centre of a circle of radius $r$ (with $d < r$), the chord it cuts on the circle has length
\[ L = 2\sqrt{r^{2} - d^{2}}. \]
\end{propriete}

\begin{popfigure}
\begin{center}
\begin{tikzpicture}[scale=1.6]
\draw[popline, thick] (0,0) circle (2);
\draw[popblue, thick] (-2.3,1.2) -- (2.3,1.2);
\fill[poporange] (-1.6,1.2) circle (0.05) node[above left, popdark]{\small $A$};
\fill[poporange] (1.6,1.2) circle (0.05) node[above right, popdark]{\small $B$};
\fill[popdark] (0,0) circle (0.05) node[below left]{\small $C$};
\draw[popgreen, thick] (0,0) -- (0,1.2);
\fill[popgreen] (0,1.2) circle (0.04);
\draw[popgreen] (0,1.05) -- (0.15,1.05) -- (0.15,1.2);
\draw[popdark, dashed] (0,0) -- (1.6,1.2);
\node[popdark] at (0.95,0.45) {\small $r$};
\node[popgreen] at (-0.28,0.6) {\small $d$};
\node[popblue] at (0,1.45) {\small chord $AB$};
\end{tikzpicture}
\end{center}
\legende{The perpendicular from the centre $C$ bisects the chord $AB$; Pythagoras gives $L = 2\sqrt{r^{2}-d^{2}}$.}
\end{popfigure}

\begin{exoresolu}[Length of a chord]
Find the length of the chord cut on the circle $x^{2} + y^{2} - 6x - 8y = 0$ by the line $x + y = 7$.
\tcblower
\textbf{Solution.} Complete the squares: $(x-3)^{2} + (y-4)^{2} = 25$, so the centre is $C(3,4)$ and $r = 5$. The distance from $C$ to the line $x + y - 7 = 0$ is
\[ d = \frac{|3 + 4 - 7|}{\sqrt{1^{2}+1^{2}}} = \frac{0}{\sqrt{2}} = 0. \]
The line passes through the centre: the chord is a \textbf{diameter}, and $L = 2\sqrt{25 - 0} = 10$. (Indeed $L = 2r$, as expected for a diameter.)
\end{exoresolu}

\begin{exercice}[Chord length]
Find the length of the chord cut on the circle $x^{2} + y^{2} = 25$ by the line $3x + 4y - 15 = 0$.
\end{exercice}

\begin{corrige}[Exercise 1]
Centre $(0,0)$, $r = 5$. Distance: $d = \dfrac{|{-15}|}{\sqrt{9+16}} = \dfrac{15}{5} = 3$. Chord length: $L = 2\sqrt{25 - 9} = 2\sqrt{16} = 8$.
\end{corrige}

\courssub{Intersection of Two Circles: The Common Chord}

Suppose now we have two circles $\mathscr{C}_{1}$ and $\mathscr{C}_{2}$. Where do they meet? Solving two quadratic equations simultaneously looks heavy, but a clever trick makes it easy: \emph{subtract} one equation from the other. Since both equations contain $x^{2} + y^{2}$ with coefficient $1$ (after normalisation), the subtraction kills all the quadratic terms and leaves a \textbf{linear equation} — the equation of a straight line. Any point lying on both circles satisfies both equations, hence satisfies their difference: the intersection points lie on this line, called the \cle{common chord} (or radical axis, studied in Section 7).

\begin{methode}[Finding the intersection points of two circles]
\begin{enumerate}
\item Write both circle equations in expanded general form.
\item Subtract one equation from the other: the $x^{2}$ and $y^{2}$ terms cancel, leaving the equation of the \textbf{common chord} (a straight line).
\item Solve this line simultaneously with \emph{either} circle (substitution, as in Lesson 12).
\item The solutions are the intersection points of the two circles (at most two).
\end{enumerate}
\end{methode}

\begin{exoresolu}[Intersection of two circles]
Find the points of intersection of the circles
\[ x^{2} + y^{2} - 4x - 2y - 4 = 0 \qquad \text{and} \qquad x^{2} + y^{2} + 2x - 6y + 6 = 0. \]
\tcblower
\textbf{Solution.} Subtract the second equation from the first:
\[ (-4x - 2x) + (-2y + 6y) + (-4 - 6) = 0 \;\Longrightarrow\; -6x + 4y - 10 = 0 \;\Longrightarrow\; 3x - 2y + 5 = 0. \]
So $y = \dfrac{3x + 5}{2}$. Substitute into the first circle:
\begin{align*}
x^{2} + \left(\tfrac{3x+5}{2}\right)^{2} - 4x - (3x+5) - 4 &= 0\\
4x^{2} + 9x^{2} + 30x + 25 - 16x - 12x - 20 - 16 &= 0\\
13x^{2} + 2x - 11 &= 0.
\end{align*}
Then $x = \dfrac{-2 \pm \sqrt{4 + 572}}{26} = \dfrac{-2 \pm 24}{26}$, giving $x = \dfrac{11}{13}$ or $x = -1$.

For $x = -1$: $y = 1$. For $x = \tfrac{11}{13}$: $y = \tfrac{49}{13}$.

The circles intersect at $(-1, 1)$ and $\left(\dfrac{11}{13}, \dfrac{49}{13}\right)$.
\end{exoresolu}

\courssub{Conditions for Two Circles to Touch}

Two circles \cle{touch} when they have exactly one common point. Draw the line through the two centres $C_{1}$ and $C_{2}$: by symmetry, the point of contact $T$ always lies on this \emph{line of centres}. Measuring the distance $C_{1}C_{2}$ against the radii $r_{1}$ and $r_{2}$ gives a complete classification.

\begin{propriete}[Touching conditions for two circles]
Let two circles have centres $C_{1}, C_{2}$ and radii $r_{1}, r_{2}$, and let $d = C_{1}C_{2}$. Then:
\begin{itemize}
\item the circles touch \textbf{externally} if and only if $d = r_{1} + r_{2}$;
\item the circles touch \textbf{internally} if and only if $d = |r_{1} - r_{2}|$.
\end{itemize}
In both cases the point of contact $T$ lies on the line of centres and divides $C_{1}C_{2}$ in the ratio $r_{1} : r_{2}$ (internally for external touching, externally for internal touching).
\end{propriete}

\begin{popfigure}
\begin{center}
\begin{tikzpicture}[scale=0.9]
% External touching
\begin{scope}[shift={(0,0)}]
\draw[popblue, thick] (0,0) circle (1.5);
\draw[poporange, thick] (3.5,0) circle (2);
\fill[popdark] (0,0) circle (0.06) node[below]{\small $C_{1}$};
\fill[popdark] (3.5,0) circle (0.06) node[below]{\small $C_{2}$};
\fill[popgreen] (1.5,0) circle (0.07) node[above]{\small $T$};
\draw[popline, dashed] (-1.5,0) -- (5.5,0);
\node at (1.75,-2.4) {\small external: $C_{1}C_{2} = r_{1}+r_{2}$};
\end{scope}
% Internal touching
\begin{scope}[shift={(10,0)}]
\draw[popblue, thick] (0,0) circle (2.2);
\draw[poporange, thick] (1.1,0) circle (1.1);
\fill[popdark] (0,0) circle (0.06) node[below left]{\small $C_{1}$};
\fill[popdark] (1.1,0) circle (0.06) node[below]{\small $C_{2}$};
\fill[popgreen] (2.2,0) circle (0.07) node[above right]{\small $T$};
\draw[popline, dashed] (-2.2,0) -- (3.2,0);
\node at (0.5,-2.7) {\small internal: $C_{1}C_{2} = |r_{1}-r_{2}|$};
\end{scope}
\end{tikzpicture}
\end{center}
\legende{External and internal touching: in both cases $C_{1}$, $T$ and $C_{2}$ are collinear.}
\end{popfigure}

\begin{exemplebox}[Proving that two circles touch]
Show that the circles $x^{2} + y^{2} = 4$ and $x^{2} + y^{2} - 6x - 8y + 24 = 0$ touch externally, and find the point of contact.

First circle: $C_{1}(0,0)$, $r_{1} = 2$. Second circle: $(x-3)^{2} + (y-4)^{2} = 1$, so $C_{2}(3,4)$, $r_{2} = 1$. Then $C_{1}C_{2} = \sqrt{9+16} = 5 = r_{1} + r_{2}$? No: $r_{1} + r_{2} = 3 \neq 5$. Let us recheck: $(x-3)^2+(y-4)^2 = 9+16-24 = 1$, so $r_2=1$ and $C_1C_2 = 5 > 3$: the circles are \textbf{separate}, they do not touch. This example shows why one must always \emph{compute} rather than assume!

A correct touching pair: $x^{2}+y^{2}=4$ and $(x-3)^{2}+y^{2}=1$: $C_{1}C_{2} = 3 = 2+1$, external touching at $T = (2, 0)$, dividing $C_{1}C_{2}$ in the ratio $2:1$.
\end{exemplebox}

\begin{attention}[Sum or difference of radii?]
The most frequent exam error is to use $r_{1} + r_{2}$ for \emph{internal} touching. Remember:
\begin{itemize}
\item \textbf{External} touching (circles side by side): $d = r_{1} + r_{2}$ (the \emph{sum});
\item \textbf{Internal} touching (one circle inside the other): $d = |r_{1} - r_{2}|$ (the \emph{difference}).
\end{itemize}
Always sketch the situation before choosing the formula.
\end{attention}

\begin{aretenir}[Position of two circles: complete classification]
With $d = C_{1}C_{2}$:
\begin{center}
\begin{tabularx}{\linewidth}{|l|X|}\hline
\rowcolor{popblueL} \textbf{Condition} & \textbf{Position of the circles} \\ \hline
$d > r_{1} + r_{2}$ & Separate (no common point, each outside the other) \\ \hline
$d = r_{1} + r_{2}$ & Touch externally (one common point) \\ \hline
$|r_{1} - r_{2}| < d < r_{1} + r_{2}$ & Intersect in two distinct points \\ \hline
$d = |r_{1} - r_{2}|$ & Touch internally (one common point) \\ \hline
$d < |r_{1} - r_{2}|$ & One inside the other, no common point \\ \hline
\end{tabularx}
\end{center}
\end{aretenir}

\begin{exoresolu}[Classifying the position of two circles]
Determine the relative position of the circles
\[ \mathscr{C}_{1}: x^{2} + y^{2} - 2x - 4y - 4 = 0 \qquad \mathscr{C}_{2}: x^{2} + y^{2} + 4x + 2y - 11 = 0. \]
\tcblower
\textbf{Solution.} Complete the squares:
\[ \mathscr{C}_{1}: (x-1)^{2} + (y-2)^{2} = 9 \Rightarrow C_{1}(1,2),\ r_{1} = 3; \]
\[ \mathscr{C}_{2}: (x+2)^{2} + (y+1)^{2} = 16 \Rightarrow C_{2}(-2,-1),\ r_{2} = 4. \]
Distance between centres: $d = \sqrt{(1+2)^{2} + (2+1)^{2}} = \sqrt{18} = 3\sqrt{2} \approx 4.24$.

Compare: $r_{1} + r_{2} = 7$ and $|r_{1} - r_{2}| = 1$. Since $1 < 3\sqrt{2} < 7$, the circles \textbf{intersect in two distinct points}.
\end{exoresolu}

\begin{exercice}[Touching circles]
Find the value of the constant $k > 0$ for which the circle $x^{2} + y^{2} = 9$ touches the circle $(x - k)^{2} + y^{2} = 4$ externally, and give the point of contact.
\end{exercice}

\begin{corrige}[Exercise 2]
$C_{1}(0,0)$, $r_{1} = 3$; $C_{2}(k,0)$, $r_{2} = 2$. External touching requires $C_{1}C_{2} = r_{1} + r_{2}$, i.e. $k = 5$. The point of contact divides $C_{1}C_{2}$ in the ratio $3:2$ from $C_{1}$: $T = \left(\dfrac{3}{5}\times 5, 0\right) = (3, 0)$.
\end{corrige}

\begin{savaistu}[Why subtraction works]
Subtracting the equations of two circles always gives a line, even when the circles do not meet at all! That line is called the \cle{radical axis} and it hides a remarkable property: from any of its points, the tangent lengths to both circles are equal. We shall explore this elegant object in Section 7.
\end{savaistu}

We now know how to decide, purely by computation, whether a line meets a circle and whether two circles cross, touch, or avoid each other. In Section 6 we shall reverse the question: given the intersection points of two circles, how can we describe \emph{all} the circles passing through them?

\coursec{Section 6: Families of Circles Through Intersections}

In Section 5 we learned how to \emph{find} the points where two circles meet, or where a line cuts a circle. We now ask the reverse question: \emph{given} two fixed intersection points, can we describe, in a single stroke, \textbf{all} the circles that pass through them? The answer is one of the most elegant labour-saving devices in the whole of coordinate geometry, and it is a favourite of GCE examiners: instead of solving three simultaneous equations for the three unknowns of a circle, we write down a \cle{family of circles} depending on a single parameter $\lambda$, and then fix $\lambda$ with one extra condition. One equation, one unknown — instead of three.

\courssub{The Family of Circles Through the Intersection of Two Circles}

\textbf{Intuition first.}\par
Draw two circles crossing at points $A$ and $B$. Through $A$ and $B$ you can draw infinitely many other circles: tiny ones, huge ones, ones leaning left or right — but they all share the same two points $A$ and $B$. We want one algebraic formula that generates every one of them by turning a single ``knob'' $\lambda$.

\begin{propriete}[Family of circles through the intersection of two circles]
Let
\[ S_1 \equiv x^2 + y^2 + 2g_1x + 2f_1y + c_1 = 0 \quad\text{and}\quad S_2 \equiv x^2 + y^2 + 2g_2x + 2f_2y + c_2 = 0 \]
be two intersecting circles. Then for every real number $\lambda \neq -1$, the equation
\[ S_1 + \lambda S_2 = 0 \]
represents a \cle{circle} passing through the intersection points of $S_1 = 0$ and $S_2 = 0$. As $\lambda$ varies, we obtain the whole family of circles through those two common points.
\end{propriete}

\textbf{Justification (examinable).}\par
Two things must be checked: that the combination is a circle, and that it passes through the common points.

\emph{(i) It passes through the common points.} Let $A(x_1, y_1)$ be a point of intersection of the two circles. Then $S_1(x_1, y_1) = 0$ \emph{and} $S_2(x_1, y_1) = 0$, because $A$ lies on both circles. Hence
\[ S_1(x_1, y_1) + \lambda S_2(x_1, y_1) = 0 + \lambda \cdot 0 = 0, \]
so $A$ lies on $S_1 + \lambda S_2 = 0$ for \emph{every} value of $\lambda$. The same holds for the second intersection point $B$.

\emph{(ii) It is a circle.} Expanding,
\[ (1+\lambda)x^2 + (1+\lambda)y^2 + 2(g_1 + \lambda g_2)x + 2(f_1 + \lambda f_2)y + (c_1 + \lambda c_2) = 0. \]
The coefficients of $x^2$ and $y^2$ are \emph{equal} (both are $1+\lambda$) and there is no $xy$ term, so provided $1 + \lambda \neq 0$ we may divide through by $1+\lambda$ and obtain the general equation of a circle. $\blacksquare$

\begin{attention}[Normalise before combining!]
The theorem requires both circles to be written with coefficient $1$ for $x^2$ and $y^2$. If you are given $2x^2 + 2y^2 + 4x - 3 = 0$, divide by $2$ \textbf{first}. Combining unnormalised equations breaks the equality of the $x^2$ and $y^2$ coefficients and the result may not be a circle at all.
\end{attention}

\begin{attention}[The degenerate case $\lambda = -1$]
When $\lambda = -1$, the combination becomes $S_1 - S_2 = 0$: the $x^2$ and $y^2$ terms cancel and we are left with a \emph{straight line} — the \cle{common chord} (the radical axis, studied in Section 7). So $S_1 + \lambda S_2 = 0$ is a circle only if the $x^2$ and $y^2$ coefficients stay equal \emph{and nonzero}: always check after choosing $\lambda$.
\end{attention}

\begin{popfigure}
\begin{center}
\begin{tikzpicture}[scale=0.85]
  % intersection points of the two base circles
  \coordinate (A) at (2,1.2);
  \coordinate (B) at (2,-1.2);
  % family members through A and B (centres on x-axis)
  \draw[popblue, thick] (0.4,0) circle (1.96);
  \draw[poporange, thick] (3.6,0) circle (1.96);
  \draw[poppurple, thick] (2,0) circle (1.2);
  \draw[popgreen, thick, dashed] (2,0) circle (2.4);
  % common chord
  \draw[popmuted, dashed] (2,-2.6) -- (2,2.6);
  \fill[popdark] (A) circle (1.8pt) node[above right, black]{$A$};
  \fill[popdark] (B) circle (1.8pt) node[below right, black]{$B$};
  \node[popblue] at (-1.2,1.6) {$S_1=0$};
  \node[poporange] at (5.2,1.6) {$S_2=0$};
  \node[poppurple] at (2,0.55) {\scriptsize $\lambda_1$};
  \node[popgreen] at (2,-2.15) {\scriptsize $\lambda_2$};
\end{tikzpicture}
\legende{Members of the family $S_1 + \lambda S_2 = 0$: every circle passes through the same two points $A$ and $B$; the dashed line is the common chord ($\lambda = -1$).}
\end{center}
\end{popfigure}

\courssub{Determining $\lambda$ from an Extra Condition}

The family $S_1 + \lambda S_2 = 0$ contains one free parameter, so \textbf{one} additional condition fixes it. The most common condition in the GCE is that the required circle passes through a given point.

\begin{methode}[Circle through the intersection of two circles and a given point]
\begin{enumerate}
  \item Write both circles in normalised form $S_1 = 0$, $S_2 = 0$ (coefficients of $x^2, y^2$ equal to $1$).
  \item Write the family $S_1 + \lambda S_2 = 0$.
  \item Substitute the coordinates of the given point into the family: this gives a linear equation in $\lambda$.
  \item Solve for $\lambda$, substitute back, and simplify.
  \item Check: the result must have equal, nonzero coefficients of $x^2$ and $y^2$.
\end{enumerate}
\end{methode}

\begin{exoresolu}[GCE-style: circle through the intersection of two circles and a point]
Find the equation of the circle which passes through the points of intersection of the circles
\[ x^2 + y^2 - 4x - 6y - 12 = 0 \quad\text{and}\quad x^2 + y^2 + 6x + 4y - 12 = 0, \]
and also through the point $(2, 1)$.
\tcblower
\textbf{Solution.} Both equations are already normalised (coefficients of $x^2$ and $y^2$ are $1$). The family through their intersections is
\[ (x^2 + y^2 - 4x - 6y - 12) + \lambda(x^2 + y^2 + 6x + 4y - 12) = 0. \]
Substitute the point $(2, 1)$:
\begin{align*}
S_1(2,1) &= 4 + 1 - 8 - 6 - 12 = -21, \\
S_2(2,1) &= 4 + 1 + 12 + 4 - 12 = 9.
\end{align*}
Hence $-21 + 9\lambda = 0 \;\Longrightarrow\; \lambda = \frac{7}{3}$.
Substitute $\lambda = \frac{7}{3}$ back into the family and multiply by $3$ to clear denominators:
\[ 3(x^2 + y^2 - 4x - 6y - 12) + 7(x^2 + y^2 + 6x + 4y - 12) = 0, \]
\[ 10x^2 + 10y^2 + 30x + 10y - 120 = 0 \;\Longrightarrow\; \boxed{x^2 + y^2 + 3x + y - 12 = 0.} \]
\textbf{Check:} coefficients of $x^2$ and $y^2$ are equal ($1$) and nonzero. The centre is $\left(-\frac{3}{2}, -\frac{1}{2}\right)$ and radius $\sqrt{\frac{9}{4}+\frac{1}{4}+12} = \sqrt{\frac{58}{4}} = \frac{\sqrt{58}}{2}$. $\blacksquare$
\end{exoresolu}

\begin{attention}[A trap hidden in the algebra]
If substituting the given point makes the $S_2$ bracket vanish but not the $S_1$ bracket, no finite $\lambda$ works: the point lies on $S_2$, and the answer is $S_2$ itself. If both brackets vanish, the point is one of the intersection points and gives no information (the family already passes through it). Always interpret the algebra geometrically.
\end{attention}

\courssub{The Family of Circles Through the Intersection of a Line and a Circle}

Exactly the same idea works when one of the two curves is a straight line.

\begin{propriete}[Family of circles through the intersection of a circle and a line]
Let $S \equiv x^2 + y^2 + 2gx + 2fy + c = 0$ be a circle and $L \equiv ax + by + d = 0$ a line cutting it. Then for every real $\lambda$,
\[ S + \lambda L = 0 \]
represents a circle through the two points where $L$ meets $S$.
\end{propriete}

\textbf{Justification.}\par
If $P$ is a common point of the circle and the line, then $S(P) = 0$ and $L(P) = 0$, so $S(P) + \lambda L(P) = 0$ for all $\lambda$: every member of the family passes through $P$. Moreover $S + \lambda L = 0$ expands to
\[ x^2 + y^2 + (2g + \lambda a)x + (2f + \lambda b)y + (c + \lambda d) = 0, \]
which has equal coefficients ($1$) of $x^2$ and $y^2$ and no $xy$ term — always a circle, for \emph{every} $\lambda$ (there is no degenerate value here, since $L$ contributes no quadratic terms). $\blacksquare$

\begin{popfigure}
\begin{center}
\begin{tikzpicture}[scale=0.9]
  \coordinate (P) at (0.6,1.5);
  \coordinate (Q) at (0.6,-1.5);
  % base circle through P,Q
  \draw[popblue, thick] (2.1,0) circle (2.12);
  % secant line
  \draw[popdark, thick] (0.6,-2.8) -- (0.6,2.8) node[above]{$L=0$};
  % other family members
  \draw[poporange, thick] (3.4,0) circle (3.17);
  \draw[poppurple, thick, dashed] (1.35,0) circle (1.66);
  \fill[popdark] (P) circle (1.8pt) node[above left, black]{$P$};
  \fill[popdark] (Q) circle (1.8pt) node[below left, black]{$Q$};
  \node[popblue] at (3.6,1.9) {$S=0$};
  \node[poporange] at (5.9,1.4) {\scriptsize $S+\lambda_1 L$};
  \node[poppurple] at (1.35,-0.9) {\scriptsize $S+\lambda_2 L$};
\end{tikzpicture}
\legende{The family $S + \lambda L = 0$: all circles pass through $P$ and $Q$, the intersections of the circle $S=0$ with the secant line $L=0$.}
\end{center}
\end{popfigure}

The parameter $\lambda$ is again fixed by one extra condition: a point the circle must pass through, or a condition on its centre (for instance, that the centre lies on a given line).

\begin{exoresolu}[GCE-style: circle through the intersection of a line and a circle, with centre on a given line]
Find the equation of the circle passing through the points of intersection of the circle
\[ x^2 + y^2 - 2x - 4y - 4 = 0 \]
and the line $x + 2y - 4 = 0$, given that its centre lies on the line $x - y - 4 = 0$.
\tcblower
\textbf{Solution.} The family is
\[ x^2 + y^2 - 2x - 4y - 4 + \lambda(x + 2y - 4) = 0, \]
i.e.
\[ x^2 + y^2 + (\lambda - 2)x + (2\lambda - 4)y - 4 - 4\lambda = 0. \]
Its centre is
\[ \left( -\frac{\lambda - 2}{2},\; -\frac{2\lambda - 4}{2} \right) = \left( 1 - \frac{\lambda}{2},\; 2 - \lambda \right). \]
The centre lies on $x - y - 4 = 0$, so
\[ \left(1 - \frac{\lambda}{2}\right) - (2 - \lambda) - 4 = 0 \;\Longrightarrow\; \frac{\lambda}{2} - 5 = 0 \;\Longrightarrow\; \lambda = 10. \]
Substituting back:
\[ x^2 + y^2 + 8x + 16y - 44 = 0. \]
\textbf{Check:} centre $(-4, -8)$; indeed $-4 - (-8) - 4 = 0$. $\blacksquare$
\end{exoresolu}

\begin{methode}[Using a tangency condition with $S + \lambda L = 0$]
If the required circle must touch a given line $L_1 = 0$:
\begin{enumerate}
  \item Write the family $S + \lambda L = 0$ and identify its centre $(-g(\lambda), -f(\lambda))$ and radius $r(\lambda) = \sqrt{g^2 + f^2 - c}$.
  \item Impose that the perpendicular distance from the centre to $L_1$ equals $r(\lambda)$.
  \item Solve the resulting equation for $\lambda$ (often a quadratic — two valid circles may exist).
\end{enumerate}
\end{methode}

\begin{exoresolu}[GCE-style: circle through intersection of line and circle, touching another line]
Find the equations of the circles passing through the intersection of the circle $x^2 + y^2 = 4$ and the line $x - y = 0$, which also touch the line $x = 5$.
\tcblower
\textbf{Solution.}
\begin{enumerate}
  \item The family is $S + \lambda L = 0$:
  \[ x^2 + y^2 - 4 + \lambda(x - y) = 0 \;\Longrightarrow\; x^2 + y^2 + \lambda x - \lambda y - 4 = 0. \]
  \item Centre: $\left(-\frac{\lambda}{2},\; \frac{\lambda}{2}\right)$. Radius squared:
  \[ r^2 = \left(\frac{\lambda}{2}\right)^2 + \left(\frac{\lambda}{2}\right)^2 + 4 = \frac{\lambda^2}{2} + 4. \]
  \item Distance from centre to line $x = 5$ (i.e. $x - 5 = 0$):
  \[ d = \left| -\frac{\lambda}{2} - 5 \right| = \left| \frac{\lambda}{2} + 5 \right|. \]
  \item Tangency condition $d = r$:
  \[ \left(\frac{\lambda}{2} + 5\right)^2 = \frac{\lambda^2}{2} + 4. \]
  Expand: $\frac{\lambda^2}{4} + 5\lambda + 25 = \frac{\lambda^2}{2} + 4$.
  Multiply by $4$: $\lambda^2 + 20\lambda + 100 = 2\lambda^2 + 16$.
  $\lambda^2 - 20\lambda - 84 = 0 \;\Longrightarrow\; (\lambda - 22)(\lambda + 2) = 0$.
  Hence $\lambda = 22$ or $\lambda = -2$.
  \item For $\lambda = 22$: $x^2 + y^2 + 22x - 22y - 4 = 0$.
  For $\lambda = -2$: $x^2 + y^2 - 2x + 2y - 4 = 0$.
  Both are valid circles (check: coefficients of $x^2,y^2$ equal and nonzero).
\end{enumerate}
$\blacksquare$
\end{exoresolu}

\begin{exemplebox}[Quick check of the method]
The circle through the intersection of $x^2 + y^2 = 4$ and $x - y = 0$ which also passes through $(3, 0)$: the family is $x^2 + y^2 - 4 + \lambda(x - y) = 0$. Substituting $(3,0)$: $9 - 4 + 3\lambda = 0$, so $\lambda = -\frac{5}{3}$, giving $3x^2 + 3y^2 - 5x + 5y - 12 = 0$.
\end{exemplebox}

\courssub{Summary and Common Mistakes}

\begin{aretenir}[Key points of Section 6]
\begin{itemize}
  \item Two intersecting circles $S_1 = 0$, $S_2 = 0$ (both normalised): the family $S_1 + \lambda S_2 = 0$ gives every circle through their intersection points; $\lambda = -1$ gives the common chord (a line).
  \item Circle $S = 0$ cut by line $L = 0$: the family $S + \lambda L = 0$ gives every circle through the two intersection points; it is a circle for \emph{all} $\lambda$.
  \item One extra condition (a point, a centre condition, a tangency condition) determines $\lambda$.
  \item Always verify at the end that the coefficients of $x^2$ and $y^2$ are equal and nonzero.
\end{itemize}
\end{aretenir}

\begin{attention}[Frequent errors at the GCE]
\begin{itemize}
  \item Combining circles whose $x^2$ coefficients are not both $1$ — normalise first.
  \item Substituting the extra point into only one of $S_1, S_2$ instead of the full combination.
  \item Forgetting that $\lambda = -1$ in $S_1 + \lambda S_2 = 0$ yields a line, not a circle.
  \item Sign slips when expanding $S + \lambda L$: the centre coordinates become $-\frac{1}{2}(2g + \lambda a)$ and $-\frac{1}{2}(2f + \lambda b)$ — keep the brackets.
  \item In tangency problems, forgetting to square the distance equation, leading to missing the $\pm$ and losing a solution.
\end{itemize}
\end{attention}

\begin{exercice}[Practice]
\begin{enumerate}
  \item Find the equation of the circle through the intersection of $x^2 + y^2 - 4 = 0$ and $x^2 + y^2 - 2x - 4y + 4 = 0$ which passes through $(2, 3)$.
  \item Find the equation of the circle through the intersection of the circle $x^2 + y^2 - 6x + 2y + 5 = 0$ and the line $x + y - 3 = 0$, whose centre lies on the $x$-axis.
  \item Show that the circle through the intersection of $x^2 + y^2 = 9$ and $2x - y + 3 = 0$ which touches the line $x = 5$ satisfies a quadratic in $\lambda$, and comment on the number of solutions.
\end{enumerate}
\end{exercice}

\begin{corrige}[Exercise 1]
Family: $x^2 + y^2 - 4 + \lambda(x^2 + y^2 - 2x - 4y + 4) = 0$. At $(2,3)$: $(4+9-4) + \lambda(4+9-4-12+4) = 9 + \lambda(1) = 0$, so $\lambda = -9$. Substituting: $(1-9)x^2 + (1-9)y^2 + 18x + 36y + (-4-36) = 0$, i.e. $-8x^2 - 8y^2 + 18x + 36y - 40 = 0$. Divide by $-2$: $\boxed{4x^2 + 4y^2 - 9x - 18y + 20 = 0}$.
\end{corrige}

\begin{corrige}[Exercise 2]
Family: $x^2 + y^2 - 6x + 2y + 5 + \lambda(x + y - 3) = 0$, centre $\left(\frac{6-\lambda}{2}, \frac{-2-\lambda}{2}\right)$. Centre on the $x$-axis $\Rightarrow$ $y$-coordinate $= 0$: $\frac{-2-\lambda}{2} = 0 \Rightarrow \lambda = -2$. Result: $x^2 + y^2 - 6x + 2y + 5 -2x -2y + 6 = 0$, i.e. $\boxed{x^2 + y^2 - 8x + 11 = 0}$.
\end{corrige}

\begin{corrige}[Exercise 3]
Family: $x^2 + y^2 - 9 + \lambda(2x - y + 3) = 0$; centre $(-\lambda, \frac{\lambda}{2})$, radius squared $r^2 = \lambda^2 + \frac{\lambda^2}{4} - (-9+3\lambda) = \frac{5\lambda^2}{4} - 3\lambda + 9$. Distance from centre to $x = 5$ is $|5 + \lambda|$. Setting $(5+\lambda)^2 = \frac{5\lambda^2}{4} - 3\lambda + 9$ gives $\lambda^2 + 10\lambda + 25 = \frac{5\lambda^2}{4} - 3\lambda + 9$. Multiply by $4$: $4\lambda^2 + 40\lambda + 100 = 5\lambda^2 - 12\lambda + 36 \Rightarrow \lambda^2 - 52\lambda - 64 = 0$. Discriminant $\Delta = 2704 + 256 = 2960 > 0$, so two distinct real roots for $\lambda$, hence \textbf{two circles} satisfy the conditions.
\end{corrige}

\begin{savaistu}[Why examiners love this technique]
In the GCE Advanced Level paper, a question on families of circles typically carries 5–7 marks and can be done in under ten lines with $S_1 + \lambda S_2 = 0$, whereas the direct method (finding both intersection points, then a circle through three points) often fills a whole page and invites arithmetic errors. Mastering the parameter $\lambda$ is one of the highest-value skills of this chapter.
\end{savaistu}

\coursec{Section 7: Radical Axis and Locus Problems}

In Section 6 we studied families of circles passing through the intersection points of two given circles, and we met the striking fact that the quantity $S_1 - S_2$ is a \emph{straight line} whenever $S_1 = 0$ and $S_2 = 0$ are two circles. In this section we give that line a name — the \cle{radical axis} — and we discover its beautiful geometric meaning. We then turn to the general art of \cle{locus problems}, where a point moves according to a rule and we must find the curve it traces. Many of the most elegant loci turn out to be circles, which is why this topic belongs naturally to this chapter.

\courssub{The Radical Axis of Two Circles}

\textbf{An observation to build intuition.}\par
Take two circles and a point $P$ outside both of them. From $P$ we can draw a tangent to the first circle, of length $t_1$, and a tangent to the second circle, of length $t_2$. In general $t_1 \neq t_2$. But as $P$ moves, there are special positions where $t_1 = t_2$. Where are all such points? Let us find out algebraically.

Recall from Section 4 that the length of the tangent from $P(x_1, y_1)$ to the circle $S = x^2 + y^2 + 2gx + 2fy + c = 0$ is
\[
t = \sqrt{S(x_1, y_1)} = \sqrt{x_1^2 + y_1^2 + 2gx_1 + 2fy_1 + c}.
\]

\begin{definition}[Radical axis]
The \cle{radical axis} of two circles is the locus of all points from which the lengths of the tangents to the two circles are equal.
\end{definition}

\begin{propriete}[Equation of the radical axis]
Let $S_1 = x^2 + y^2 + 2g_1x + 2f_1y + c_1 = 0$ and $S_2 = x^2 + y^2 + 2g_2x + 2f_2y + c_2 = 0$ be two circles (with the coefficients of $x^2$ and $y^2$ both equal to $1$). Their radical axis is the straight line
\[
S_1 - S_2 = 0, \qquad\text{i.e.}\qquad 2(g_1 - g_2)x + 2(f_1 - f_2)y + (c_1 - c_2) = 0.
\]
Moreover, the radical axis is \textbf{perpendicular to the line of centres} of the two circles. (The derivation is examinable and is given below.)
\end{propriete}

\textbf{Derivation.} A point $P(x, y)$ lies on the radical axis exactly when its tangent lengths to the two circles are equal: $t_1 = t_2$, hence $t_1^2 = t_2^2$, that is $S_1(x,y) = S_2(x,y)$, or
\[
S_1 - S_2 = 0.
\]
The terms $x^2$ and $y^2$ cancel (both circles have coefficient $1$), leaving
\[
2(g_1 - g_2)x + 2(f_1 - f_2)y + (c_1 - c_2) = 0,
\]
which is linear in $x$ and $y$: it is a straight line. The centres are $C_1(-g_1, -f_1)$ and $C_2(-g_2, -f_2)$, so the line of centres has direction vector $\begin{pmatrix} g_2 - g_1 \\ f_2 - f_1 \end{pmatrix}$. The radical axis has normal vector $\begin{pmatrix} 2(g_1 - g_2) \\ 2(f_1 - f_2) \end{pmatrix}$, which is exactly $-2$ times the direction of $C_1C_2$. A line whose normal is parallel to $C_1C_2$ is perpendicular to $C_1C_2$. $\blacksquare$

\begin{attention}[Normalise first!]
The cancellation of $x^2$ and $y^2$ works only if \textbf{both} equations have coefficient $1$ for $x^2$ and $y^2$. If you are given $2x^2 + 2y^2 + \cdots = 0$, divide through by $2$ \emph{before} subtracting. Subtracting unnormalised equations gives a wrong line — a classic GCE trap.
\end{attention}

\begin{propriete}[Radical axis and common chord]
If the two circles \textbf{intersect}, their radical axis is the line containing their \cle{common chord}. If the circles \textbf{touch}, the radical axis is their \textbf{common tangent} at the point of contact. If the circles do not meet, the radical axis lies between them and meets neither circle.
\end{propriete}

Indeed, at an intersection point both $S_1 = 0$ and $S_2 = 0$, so $S_1 - S_2 = 0$ as well: both intersection points lie on the radical axis, and two points determine a line.

\begin{popfigure}
\begin{tikzpicture}[scale=0.85]
  % Circle 1: centre (0,0), r=2 ; Circle 2: centre (5,0), r=1.5
  % Radical axis: S1-S2 = 10x - 26.75 = 0 => x = 2.675
  \draw[thick, popblue] (0,0) circle (2);
  \draw[thick, poporange] (5,0) circle (1.5);
  \fill[popblue] (0,0) circle (1.5pt) node[below left] {$C_1$};
  \fill[poporange] (5,0) circle (1.5pt) node[below right] {$C_2$};
  \draw[dashed, popmuted] (-2.3,0) -- (7,0);
  \draw[very thick, popgreen] (2.675,-2.6) -- (2.675,2.6) node[above] {radical axis};
  % P on the axis at height 2; tangent foot on circle 1 is (0,2) since radius is vertical there
  \fill[popdark] (2.675,2.0) circle (2pt) node[above right] {$P$};
  \draw[thick, poppurple] (2.675,2.0) -- (0,2.0) node[midway, above] {$t_1$};
  \fill[poppurple] (0,2.0) circle (1.5pt);
  % tangent foot on circle 2 computed: (5.30,1.47)
  \draw[thick, poppink] (2.675,2.0) -- (5.30,1.47) node[midway, below right] {$t_2$};
  \fill[poppink] (5.30,1.47) circle (1.5pt);
\end{tikzpicture}
\legende{Two non-intersecting circles and their radical axis. From any point $P$ of the axis, the tangent segments to the two circles are equal: $t_1 = t_2$. The axis is perpendicular to the line of centres $C_1C_2$.}
\end{popfigure}

\begin{exoresolu}[Finding a radical axis]
Find the radical axis of the circles $S_1: x^2 + y^2 - 4x + 6y - 12 = 0$ and $S_2: x^2 + y^2 + 2x - 2y - 8 = 0$, and verify that it is perpendicular to the line of centres.
\tcblower
\textbf{Solution.} Both equations are already normalised (coefficient of $x^2$ is $1$). Subtract:
\begin{align*}
S_1 - S_2 &= (-4x - 2x) + (6y - (-2y)) + (-12 - (-8)) = 0\\
&= -6x + 8y - 4 = 0,
\end{align*}
so the radical axis is $3x - 4y + 2 = 0$.

The centres are $C_1(2, -3)$ and $C_2(-1, 1)$. Gradient of $C_1C_2$: $\dfrac{1-(-3)}{-1-2} = -\dfrac{4}{3}$. Gradient of the radical axis: $\dfrac{3}{4}$. Product: $-\dfrac{4}{3}\times\dfrac{3}{4} = -1$, so the lines are perpendicular, as expected. $\blacksquare$
\end{exoresolu}

\begin{savaistu}[The radical centre — a GCE-useful extra]
Given \emph{three} circles whose centres are not collinear, the three radical axes (taken in pairs) all pass through a single point called the \cle{radical centre}. It is the unique point with equal tangent lengths to all three circles. To find it, compute any two radical axes and intersect them — a fast and elegant exam technique.
\end{savaistu}

\begin{exemplebox}[Radical centre]
Consider the three circles
\[
S_1: x^2+y^2-4=0, \qquad S_2: x^2+y^2-6x=0, \qquad S_3: x^2+y^2-4y=0.
\]
The radical axis of $S_1$ and $S_2$ is $S_1 - S_2 = 6x - 4 = 0$, i.e. $x = \tfrac{2}{3}$. The radical axis of $S_1$ and $S_3$ is $S_1 - S_3 = 4y - 4 = 0$, i.e. $y = 1$. Hence the radical centre is $\left(\tfrac{2}{3}, 1\right)$. Check with the third pair: $S_2 - S_3 = -6x + 4y = 0$, and indeed $-6\times\tfrac{2}{3} + 4\times 1 = -4 + 4 = 0$. $\blacksquare$
\end{exemplebox}

\courssub{Locus of a Point — the General Strategy}

A \cle{locus} is the set of all points satisfying a given geometric condition. The word is Latin for ``place'': we ask \emph{where} a point can be. In coordinate geometry the strategy is always the same.

\begin{methode}[Solving a locus problem]
\begin{enumerate}
\item \textbf{Name} the moving point $P(x, y)$.
\item \textbf{Translate} the given condition into an algebraic equation in $x$ and $y$ (distance formula, tangent-length formula, gradient condition for perpendicularity, etc.).
\item \textbf{Simplify}: square to remove roots, expand, collect terms.
\item \textbf{Identify} the locus: state the type of curve and its key features (for a circle: centre and radius).
\end{enumerate}
\end{methode}

\begin{attention}[State the locus, not just the equation]
In GCE questions, an equation alone does not earn full marks. You must \textbf{identify} the locus: ``a circle with centre $(2,-1)$ and radius $3$'', or ``the perpendicular bisector of $AB$''. Always finish with a sentence describing the curve.
\end{attention}

\textbf{Classic loci leading to circles.}\par
\begin{itemize}
\item \textbf{Constant distance from a fixed point.} If $PA = r$ with $A$ fixed, the locus is by definition the circle of centre $A$ and radius $r$.
\item \textbf{Segment subtending a right angle.} If $\angle APB = 90^{\circ}$ with $A$, $B$ fixed, then $P$ lies on the circle with diameter $AB$ (converse of the angle in a semicircle). Algebraically, $(\text{gradient } PA)\times(\text{gradient } PB) = -1$ leads to $(x - x_A)(x - x_B) + (y - y_A)(y - y_B) = 0$.
\item \textbf{$PA^2 + PB^2$ constant.} With $A$, $B$ fixed and $PA^2 + PB^2 = k$, expanding gives $2x^2 + 2y^2 + \cdots = k$, a circle whose centre is the \emph{midpoint} of $AB$.
\item \textbf{$PA = k \cdot PB$ ($k \neq 1$).} This gives the \cle{circle of Apollonius}, studied next.
\end{itemize}

\begin{propriete}[Circle of Apollonius]
Let $A$ and $B$ be fixed points and $k > 0$, $k \neq 1$. The locus of points $P$ such that $PA = k \cdot PB$ is a circle, called the \cle{circle of Apollonius}. Its centre lies on the line $AB$, and it meets the line $AB$ at the two points dividing the segment $AB$ internally and externally in the ratio $k : 1$; these two points are the endpoints of a diameter.
\end{propriete}

\textbf{Sketch of derivation.} Put $A(-c, 0)$, $B(c, 0)$ (we may always choose axes this way). Then $PA^2 = k^2 PB^2$ gives
\[
(x+c)^2 + y^2 = k^2\big[(x-c)^2 + y^2\big].
\]
Expanding and dividing by $(1 - k^2) \neq 0$:
\[
x^2 + y^2 + \frac{2c(1+k^2)}{1-k^2}\,x + c^2 = 0,
\]
which has the form $x^2 + y^2 + 2gx + c^2 = 0$: a circle with centre on the $x$-axis, i.e. on line $AB$. $\blacksquare$

\begin{attention}[The case $k = 1$]
When $k = 1$, the condition becomes $PA = PB$ and the $x^2$, $y^2$ terms cancel completely: the locus is the \textbf{perpendicular bisector of $AB$} — a straight line, \emph{not} a circle. Always check whether $k = 1$ before announcing a circle.
\end{attention}

\begin{popfigure}
\begin{tikzpicture}[scale=0.9]
  % A=(0,0), B=(3,0), k=2: PA=2PB => x^2+y^2 = 4((x-3)^2+y^2)
  % 3x^2+3y^2-24x+36=0 => x^2+y^2-8x+12=0 => centre (4,0), r=2
  \draw[->, popmuted] (-1.5,0) -- (7,0) node[right] {$x$};
  \draw[->, popmuted] (0,-2.8) -- (0,2.8) node[above] {$y$};
  \fill[popblue] (0,0) circle (2pt) node[below left] {$A$};
  \fill[poporange] (3,0) circle (2pt) node[below] {$B$};
  \draw[very thick, popgreen] (4,0) circle (2);
  \fill[popgreen] (4,0) circle (1.5pt) node[below right] {$C$};
  % sample point P on circle: (4,2)
  \fill[popdark] (4,2) circle (2pt) node[above] {$P$};
  \draw[thick, poppurple] (0,0) -- (4,2) node[midway, above left] {$PA$};
  \draw[thick, poppink] (3,0) -- (4,2) node[midway, right] {$PB$};
  % diameter endpoints on line AB: (2,0) and (6,0)
  \fill[popgreen] (2,0) circle (1.5pt) node[below left] {$M$};
  \fill[popgreen] (6,0) circle (1.5pt) node[below right] {$N$};
\end{tikzpicture}
\legende{Circle of Apollonius for $PA = 2PB$ with $A(0,0)$, $B(3,0)$: the locus is the circle of centre $C(4,0)$ and radius $2$. $M$ and $N$ divide $AB$ internally and externally in the ratio $2:1$; $MN$ is a diameter.}
\end{popfigure}

\begin{exoresolu}[Apollonius circle]
$A(0,0)$ and $B(3,0)$ are fixed. Find the locus of $P$ such that $PA = 2PB$.
\tcblower
\textbf{Solution.} Let $P(x,y)$. The condition $PA = 2PB$ gives $PA^2 = 4PB^2$:
\begin{align*}
x^2 + y^2 &= 4\big[(x-3)^2 + y^2\big]\\
x^2 + y^2 &= 4x^2 - 24x + 36 + 4y^2\\
3x^2 + 3y^2 - 24x + 36 &= 0\\
x^2 + y^2 - 8x + 12 &= 0\\
(x-4)^2 + y^2 &= 4.
\end{align*}
\textbf{The locus is a circle of centre $(4, 0)$ and radius $2$.} Check: $M(2,0)$ divides $AB$ internally in the ratio $2:1$ and $N(6,0)$ divides it externally in that ratio; indeed $MN$ is a diameter of this circle. $\blacksquare$
\end{exoresolu}

\begin{exoresolu}[Locus with $PA^2 + PB^2$ constant]
$A(-2, 0)$ and $B(2, 0)$ are fixed. Find the locus of $P$ such that $PA^2 + PB^2 = 26$.
\tcblower
\textbf{Solution.} Let $P(x, y)$. Then
\begin{align*}
\big[(x+2)^2 + y^2\big] + \big[(x-2)^2 + y^2\big] &= 26\\
2x^2 + 2y^2 + 8 &= 26\\
x^2 + y^2 &= 9.
\end{align*}
\textbf{The locus is the circle with centre $(0,0)$ — the midpoint of $AB$ — and radius $3$.} $\blacksquare$
\end{exoresolu}

\begin{popfigure}
\begin{center}
\begin{tikzpicture}
\begin{axis}[
  width=9cm, height=7cm,
  axis lines=middle,
  xlabel={$x$}, ylabel={$y$},
  xmin=-5, xmax=5, ymin=-5, ymax=5,
  xtick={-4,-2,0,2,4}, ytick={-4,-2,0,2,4},
  legend style={at={(0.02,0.98)}, anchor=north west, font=\small},
  axis equal image,
]
\addplot[very thick, popblue, domain=0:360, samples=100] ({3*cos(x)},{3*sin(x)});
\addlegendentry{$PA^2+PB^2=26$}
\addplot[very thick, poporange, dashed, domain=0:360, samples=100] ({4*cos(x)},{4*sin(x)});
\addlegendentry{$PA^2+PB^2=40$}
\addplot[very thick, popgreen, dotted, domain=0:360, samples=100] ({2.236*cos(x)},{2.236*sin(x)});
\addlegendentry{$PA^2+PB^2=18$}
\addplot[only marks, mark=*, popdark] coordinates {(-2,0) (2,0)};
\node[font=\small, popdark] at (axis cs:-2,-0.7) {$A$};
\node[font=\small, popdark] at (axis cs:2,-0.7) {$B$};
\end{axis}
\end{tikzpicture}
\end{center}
\legende{Loci $PA^2 + PB^2 = k$ for $A(-2,0)$, $B(2,0)$ and several values of $k$: concentric circles centred at the midpoint of $AB$, with equation $x^2 + y^2 = \tfrac{k}{2} - 4$.}
\end{popfigure}

\courssub{Loci Involving Tangent Lengths}

The tangent-length formula $t^2 = S(x, y)$ turns many geometric conditions into circle equations.

\textbf{(a) Equal tangent lengths to two circles.} The locus of $P$ with equal tangent lengths to $S_1 = 0$ and $S_2 = 0$ is $S_1 = S_2$, i.e. $S_1 - S_2 = 0$: this is exactly the \cle{radical axis} revisited as a locus problem.

\textbf{(b) Constant tangent length to a fixed circle.} If the tangent from $P$ to the circle $S = 0$ (centre $C$, radius $r$) has constant length $\ell$, then, since the radius to the point of contact is perpendicular to the tangent, Pythagoras' theorem gives
\[
PC^2 = t^2 + r^2 = \ell^2 + r^2,
\]
so $P$ lies on a circle \textbf{concentric} with the given one, of radius $\sqrt{\ell^2 + r^2}$.

\begin{exoresolu}[Constant tangent length]
Find the locus of points from which the tangent to the circle $x^2 + y^2 - 4x + 2y - 4 = 0$ has length $3$.
\tcblower
\textbf{Solution.} Let $P(x, y)$. The tangent-length formula gives
\begin{align*}
x^2 + y^2 - 4x + 2y - 4 &= 9\\
x^2 + y^2 - 4x + 2y - 13 &= 0\\
(x-2)^2 + (y+1)^2 &= 18.
\end{align*}
\textbf{The locus is a circle concentric with the given circle, centre $(2, -1)$ and radius $\sqrt{18} = 3\sqrt{2}$.} Check: the given circle has $r^2 = g^2 + f^2 - c = 4 + 1 + 4 = 9$, so $r = 3$, and $\sqrt{\ell^2 + r^2} = \sqrt{9 + 9} = 3\sqrt{2}$, as found. $\blacksquare$
\end{exoresolu}

\begin{exoresolu}[Right angle subtended — a full locus problem]
Find the locus of $P$ such that the segment joining $A(1, 2)$ and $B(5, 4)$ subtends a right angle at $P$.
\tcblower
\textbf{Solution.} Let $P(x, y)$ with $P \neq A, B$. Perpendicularity of $PA$ and $PB$:
\[
\frac{y - 2}{x - 1}\cdot\frac{y - 4}{x - 5} = -1
\;\Longrightarrow\;
(y-2)(y-4) + (x-1)(x-5) = 0.
\]
Expanding:
\[
x^2 - 6x + 5 + y^2 - 6y + 8 = 0
\;\Longrightarrow\;
(x-3)^2 + (y-3)^2 = 5.
\]
\textbf{The locus is the circle with diameter $AB$: centre $(3, 3)$ (the midpoint of $AB$) and radius $\sqrt{5} = \tfrac{1}{2}|AB|$, excluding the points $A$ and $B$ themselves.} $\blacksquare$
\end{exoresolu}

\begin{aretenir}[Key points of Section 7]
\begin{itemize}
\item The \cle{radical axis} of $S_1 = 0$ and $S_2 = 0$ is the line $S_1 - S_2 = 0$ (normalise first!); it is the locus of points with equal tangent lengths to both circles and is perpendicular to the line of centres.
\item When the circles intersect, the radical axis is the \textbf{common chord}; three circles with non-collinear centres share a \textbf{radical centre}.
\item Locus strategy: set $P(x, y)$, translate the condition, simplify, \textbf{identify the curve}.
\item $PA = k\,PB$ gives a circle of Apollonius if $k \neq 1$, and the perpendicular bisector of $AB$ if $k = 1$.
\item $PA^2 + PB^2$ constant gives a circle centred at the midpoint of $AB$; $\angle APB = 90^{\circ}$ gives the circle on diameter $AB$; constant tangent length gives a concentric circle.
\end{itemize}
\end{aretenir}

\begin{exercice}[Practice]
\begin{enumerate}
\item Find the radical axis of $x^2 + y^2 + 3x - y + 2 = 0$ and $2x^2 + 2y^2 - 4x + 6y - 10 = 0$.
\item $A(0,0)$, $B(6,0)$. Find the locus of $P$ with $PA = 3PB$, giving the centre and radius.
\item Find the locus of points whose tangent length to $x^2 + y^2 = 16$ equals the tangent length to $x^2 + y^2 - 10x + 16 = 0$.
\item $A(-1, 0)$, $B(1, 0)$. Find the locus of $P$ with $PA^2 + PB^2 = 10$.
\end{enumerate}
\end{exercice}

\begin{corrige}[Exercise 1]
\textbf{1.} Normalise the second circle by dividing by $2$: $x^2 + y^2 - 2x + 3y - 5 = 0$. Then $S_1 - S_2 = (3x-(-2x)) + (-y-3y) + (2-(-5)) = 5x - 4y + 7 = 0$. The radical axis is $5x - 4y + 7 = 0$.

\textbf{2.} $PA = 3PB$ gives $x^2 + y^2 = 9\big[(x-6)^2 + y^2\big]$, i.e. $8x^2 + 8y^2 - 108x + 324 = 0$, or $x^2 + y^2 - \tfrac{27}{2}x + \tfrac{81}{2} = 0$. Completing the square: $\left(x - \tfrac{27}{4}\right)^2 + y^2 = \tfrac{729}{16} - \tfrac{648}{16} = \tfrac{81}{16}$. The locus is a circle of centre $\left(\tfrac{27}{4}, 0\right)$ and radius $\tfrac{9}{4}$.

\textbf{3.} Equal tangent lengths: $x^2 + y^2 - 16 = x^2 + y^2 - 10x + 16$, so $10x = 32$, i.e. $x = \tfrac{16}{5}$. The locus is the vertical line $x = \tfrac{16}{5}$ (the radical axis of the two circles).

\textbf{4.} $\big[(x+1)^2 + y^2\big] + \big[(x-1)^2 + y^2\big] = 10$ gives $2x^2 + 2y^2 + 2 = 10$, so $x^2 + y^2 = 4$. The locus is the circle of centre $(0,0)$ (midpoint of $AB$) and radius $2$.
\end{corrige}

\coursec{Section 8: Circles Under Constraints — Touching Axes and Lines}

In Section 7 we saw how the radical axis and locus arguments allow us to describe a circle indirectly, through conditions imposed on a moving point. We now exploit the same philosophy in its most practical form: a circle is \emph{completely determined} as soon as we know its centre $(a,b)$ and its radius $r$. Every geometrical constraint — touching an axis, touching a line, passing through given points — can be translated into an equation in $a$, $b$ and $r$. Solving the resulting system is the whole game. This section is one of the richest sources of GCE Advanced Level questions on circle geometry, so we proceed carefully, constraint by constraint.

\courssub{Circles Touching the Coordinate Axes}

\textbf{Touching both axes.} Imagine a disc placed in a corner of a room, touching both walls at once. The disc must sit at distance $r$ from each wall, where $r$ is its radius. The coordinate axes play the role of the two walls.

\begin{propriete}[Circle touching both coordinate axes]
A circle of radius $r$ touching \textbf{both} coordinate axes has centre $(\pm r, \pm r)$, the signs being chosen according to the quadrant. Its equation is
\[ (x \mp r)^2 + (y \mp r)^2 = r^2, \]
giving the four possibilities
\[ (x-r)^2+(y-r)^2=r^2,\quad (x+r)^2+(y-r)^2=r^2,\quad (x+r)^2+(y+r)^2=r^2,\quad (x-r)^2+(y+r)^2=r^2. \]
\end{propriete}

\textbf{Justification.} The distance from the centre $(a,b)$ to the $x$-axis is $|b|$ and to the $y$-axis is $|a|$. Touching both axes forces $|a|=r$ and $|b|=r$, hence $a=\pm r$, $b=\pm r$. $\blacksquare$

\begin{popfigure}
\begin{center}
\begin{tikzpicture}[scale=0.85]
\draw[->, thick] (-4.4,0) -- (4.4,0) node[right] {$x$};
\draw[->, thick] (0,-4.4) -- (0,4.4) node[above] {$y$};
% Quadrant I
\draw[popblue, thick] (1.6,1.6) circle (1.6);
\fill[popblue] (1.6,1.6) circle (2pt) node[above right] {\small $(r,r)$};
% Quadrant II
\draw[poporange, thick] (-1.6,1.6) circle (1.6);
\fill[poporange] (-1.6,1.6) circle (2pt) node[above left] {\small $(-r,r)$};
% Quadrant III
\draw[popgreen, thick] (-1.6,-1.6) circle (1.6);
\fill[popgreen] (-1.6,-1.6) circle (2pt) node[below left] {\small $(-r,-r)$};
% Quadrant IV
\draw[poppurple, thick] (1.6,-1.6) circle (1.6);
\fill[poppurple] (1.6,-1.6) circle (2pt) node[below right] {\small $(r,-r)$};
\fill (0,0) circle (1.5pt) node[below left] {\small $O$};
\end{tikzpicture}
\end{center}
\legende{The four circles of radius $r$ touching both axes, with centres $(\pm r,\pm r)$.}
\end{popfigure}

\textbf{Touching one axis only.} If only one axis is touched, only one coordinate of the centre is fixed in magnitude.

\begin{propriete}[Touching one axis]
A circle with centre $(a,b)$ and radius $r$ touches
\begin{itemize}
\item the $x$-axis if and only if $|b| = r$;
\item the $y$-axis if and only if $|a| = r$.
\end{itemize}
\end{propriete}

\begin{attention}[The absolute value matters]
A circle "touching the $x$-axis" with centre $(a,b)$ has radius $|b|$, \textbf{not} $b$. If the centre could lie below the axis, then $b<0$ and $r=-b$. Forgetting the absolute value is the single most common error in this topic and it silently kills one of the two solutions.
\end{attention}

\begin{savaistu}[Circles in corners — from architecture to ball bearings]
The problem of a circle touching two perpendicular lines appears everywhere: a circular table pushed into a corner, a ball bearing in a V-groove, the rounded corner of a machined part. In each case the centre lies on the angle bisector (the line $y=\pm x$) at a distance $r\sqrt{2}$ from the corner. This is why the centres are $(\pm r,\pm r)$.
\end{savaistu}

\begin{exemplebox}[Circle touching both axes through a given point]
Find the equations of the circles which touch both coordinate axes and pass through the point $(2,1)$.

\textbf{Solution.} The point $(2,1)$ lies in the first quadrant, so the centre is $(r,r)$ and the equation is $(x-r)^2+(y-r)^2=r^2$. Substituting $(2,1)$:
\[ (2-r)^2+(1-r)^2=r^2 \;\Longrightarrow\; (4-4r+r^2)+(1-2r+r^2)=r^2 \;\Longrightarrow\; r^2-6r+5=0 \;\Longrightarrow\; r=1 \text{ or } r=5. \]
Hence the two circles are
\[ (x-1)^2+(y-1)^2=1 \qquad\text{and}\qquad (x-5)^2+(y-5)^2=25. \]
Both answers are valid: a small circle tucked near the origin and a large one reaching back to the point.
\end{exemplebox}

\courssub{Touching a Given Straight Line}

The single most important translation in this whole chapter is the following.

\begin{methode}[Touching a line $\Longleftrightarrow$ distance = radius]
A circle with centre $(a,b)$ and radius $r$ touches the line $px+qy+s=0$ if and only if the perpendicular distance from the centre to the line equals the radius:
\[ \frac{|pa+qb+s|}{\sqrt{p^2+q^2}} = r. \]
Every "touching" condition in a problem should be converted immediately into this equation.
\end{methode}

\begin{exoresolu}[Circle with given centre touching a given line]
Find the equation of the circle with centre $(3,2)$ which touches the line $4x-3y+4=0$.
\tcblower
The radius is the perpendicular distance from $(3,2)$ to the line:
\[ r=\frac{|4(3)-3(2)+4|}{\sqrt{4^2+(-3)^2}}=\frac{|12-6+4|}{5}=\frac{10}{5}=2. \]
Therefore the circle is $(x-3)^2+(y-2)^2=4$.
\end{exoresolu}

\courssub{Circle Through Two Points with Centre on a Given Line}

\begin{propriete}[Perpendicular bisector]
A point is equidistant from two given points $A$ and $B$ if and only if it lies on the \cle{perpendicular bisector} of the segment $AB$. Consequently, the centre of any circle through $A$ and $B$ lies on the perpendicular bisector of $AB$.
\end{propriete}

\begin{experience}[Discovering the perpendicular bisector]
Take two pins as points $A$ and $B$ on a cork board. Stretch a rubber band around them and pull it taut with a pencil; the pencil traces the perpendicular bisector. Every position of the pencil is equidistant from $A$ and $B$, hence can be the centre of a circle through $A$ and $B$.
\end{experience}

\begin{methode}[Two points + centre on a line]
To find the circle through $A$ and $B$ with centre on a line $l$:
\begin{enumerate}
\item Write the perpendicular bisector of $AB$ (or use $CA=CB$ directly with centre $C=(a,b)$).
\item Intersect this bisector with the line $l$: the intersection is the centre $C$.
\item Compute $r=CA$ and write the equation.
\end{enumerate}
\end{methode}

\begin{exoresolu}[Centre on a given line]
Find the equation of the circle passing through $A(1,2)$ and $B(3,4)$ whose centre lies on the line $x+y=7$.
\tcblower
Let the centre be $C(a,b)$. Since $C$ lies on $x+y=7$, we have
\[ a+b=7. \tag{1} \]
Since $CA=CB$:
\[ (a-1)^2+(b-2)^2=(a-3)^2+(b-4)^2. \]
Expanding: $a^2-2a+1+b^2-4b+4 = a^2-6a+9+b^2-8b+16$, which simplifies to
\[ -2a-4b+5 = -6a-8b+25 \;\Longrightarrow\; 4a+4b=20 \;\Longrightarrow\; a+b=5. \tag{2} \]
But (1) says $a+b=7$ and (2) says $a+b=5$: the system is inconsistent. There is \textbf{no such circle}, because the line $x+y=7$ is parallel to the perpendicular bisector $x+y=5$ of $AB$ and distinct from it.

\medskip
Modify the line to $x+y=5$. Then (1) and (2) coincide, and the centre is \emph{any} point of the bisector — we need one more condition. Take instead the line $y=x$: solving $a+b=5$ with $b=a$ gives $a=b=\tfrac52$. Then
\[ r^2=\left(\tfrac52-1\right)^2+\left(\tfrac52-2\right)^2=\tfrac94+\tfrac14=\tfrac52, \]
so the circle is $\left(x-\tfrac52\right)^2+\left(y-\tfrac52\right)^2=\tfrac52$.
\end{exoresolu}

\begin{popfigure}
\begin{center}
\begin{tikzpicture}[scale=1.0]
\coordinate (A) at (1,1); \coordinate (B) at (4,2.4);
\coordinate (C) at (2.9,3.3);
% line l
\draw[popmuted, thick] (-0.5,4.1) -- (5.5,1.7) node[right] {$l$};
% perpendicular bisector (dashed)
\draw[dashed, popdark] (1.4,4.2) -- (3.9,0.3);
% circle
\draw[popblue, thick] (C) circle (2.42);
\fill[popblue] (C) circle (2pt) node[above right] {\small $C$};
\fill[poporange] (A) circle (2pt) node[below left] {\small $A$};
\fill[poporange] (B) circle (2pt) node[below right] {\small $B$};
\draw[popgreen] (C) -- (A) node[midway, left] {\small $r$};
\draw[popgreen] (C) -- (B) node[midway, right] {\small $r$};
\end{tikzpicture}
\end{center}
\legende{The centre $C$ lies at the intersection of the line $l$ with the perpendicular bisector of $AB$; then $r=CA=CB$.}
\end{popfigure}

\courssub{Circle Through Two Points and Touching a Line}

Now we combine the two ideas: the centre lies on the perpendicular bisector of $AB$ \emph{and} its distance to the given line equals its distance to $A$.

\begin{methode}[Two points + touching a line]
\begin{enumerate}
\item Parametrise the centre $C$ on the perpendicular bisector of $AB$ (one parameter).
\item Write $r^2 = CA^2$ and also $r = \dfrac{|pa+qb+s|}{\sqrt{p^2+q^2}}$ for the tangent line.
\item Equate and solve: expect a \textbf{quadratic}, hence generally \textbf{two circles}.
\end{enumerate}
\end{methode}

\begin{exoresolu}[Through two points, touching a line]
Find the equations of the circles passing through $A(2,0)$ and $B(8,0)$ and touching the line $y=6$.
\tcblower
The centre $C$ lies on the perpendicular bisector of $AB$, which is the vertical line $x=5$. Write $C=(5,k)$. Then
\[ r^2 = CA^2 = (5-2)^2+k^2 = 9+k^2. \]
Touching $y=6$ requires the distance from $C$ to the line to equal $r$:
\[ |k-6| = r \;\Longrightarrow\; (k-6)^2 = 9+k^2 \;\Longrightarrow\; k^2-12k+36 = 9+k^2 \;\Longrightarrow\; -12k = -27 \;\Longrightarrow\; k=\tfrac94. \]
Here the equation is linear (the $k^2$ terms cancel), giving one circle:
\[ r^2=9+\tfrac{81}{16}=\tfrac{225}{16}, \qquad (x-5)^2+\left(y-\tfrac94\right)^2=\tfrac{225}{16}. \]
\end{exoresolu}

\begin{exoresolu}[Touching the x-axis through two points]
Find the equations of the circles passing through $A(1,4)$ and $B(3,2)$ which touch the $x$-axis.
\tcblower
Let the centre be $C(a,b)$. Touching the $x$-axis gives $r=|b|$, so $r^2=b^2$. From $CA^2=CB^2$:
\[ (a-1)^2+(b-4)^2=(a-3)^2+(b-2)^2. \]
Expanding: $a^2-2a+1+b^2-8b+16 = a^2-6a+9+b^2-4b+4$, which simplifies to
\[ -2a-8b+17 = -6a-4b+13 \;\Longrightarrow\; 4a-4b = -4 \;\Longrightarrow\; a-b = -1, \]
so $b = a+1$. Now impose $CA^2=b^2$:
\[ (a-1)^2+(a+1-4)^2 = (a+1)^2 \;\Longrightarrow\; (a-1)^2+(a-3)^2 = (a+1)^2. \]
Expanding: $a^2-2a+1 + a^2-6a+9 = a^2+2a+1$, i.e.
\[ a^2-10a+9=0 \;\Longrightarrow\; (a-1)(a-9)=0 \;\Longrightarrow\; a=1 \text{ or } a=9. \]
Correspondingly $b=2$ or $b=10$, and $r=|b|=2$ or $10$. Two circles:
\[ a=1,\ b=2,\ r=2:\ (x-1)^2+(y-2)^2=4; \qquad a=9,\ b=10,\ r=10:\ (x-9)^2+(y-10)^2=100. \]
\end{exoresolu}

\begin{attention}[Expect two solutions — and verify them]
Constraint problems typically reduce to a quadratic, giving \textbf{two valid circles}. Never stop after one root, and always substitute back to check every constraint (as the worked example shows, an unchecked slip in the bisector equation propagates). Also reject extraneous solutions, e.g. $r\le 0$.
\end{attention}

\begin{popfigure}
\begin{center}
\begin{tikzpicture}[scale=0.62]
\coordinate (A) at (1,4); \coordinate (B) at (3,2);
\draw[->, thick] (-1.5,0) -- (11.5,0) node[right] {$x$};
\draw[->, thick] (0,-1) -- (0,11) node[above] {$y$};
\draw[popblue, thick] (1,2) circle (2);
\draw[poporange, thick] (9,10) circle (10);
\fill (A) circle (2.5pt) node[above left] {\small $A(1,4)$};
\fill (B) circle (2.5pt) node[right] {\small $B(3,2)$};
\fill[popblue] (1,2) circle (2pt) node[below] {\small $(1,2)$};
\fill[poporange] (9,10) circle (2pt) node[right] {\small $(9,10)$};
\end{tikzpicture}
\end{center}
\legende{Two circles through $A(1,4)$ and $B(3,2)$, both tangent to the $x$-axis.}
\end{popfigure}

\courssub{General Strategy for Constraint Problems}

\begin{methode}[The three-step strategy]
\begin{enumerate}
\item \textbf{List the unknowns}: centre $(a,b)$ and radius $r$ — three unknowns, so you need three independent constraints.
\item \textbf{Translate each constraint into one equation}:
\begin{itemize}
\item through point $(x_0,y_0)$: $(x_0-a)^2+(y_0-b)^2=r^2$;
\item touching line $px+qy+s=0$: $\dfrac{|pa+qb+s|}{\sqrt{p^2+q^2}}=r$;
\item centre on a line: substitute $(a,b)$ into the line's equation;
\item touching the $x$-axis / $y$-axis: $r=|b|$ / $r=|a|$.
\end{itemize}
\item \textbf{Solve the system}, usually by eliminating $r$ first (equate two expressions for $r^2$), then solve the resulting equations in $a$ and $b$. Keep all valid solutions.
\end{enumerate}
\end{methode}

\begin{exercice}[GCE-style practice]
\begin{enumerate}
\item Find the equations of the circles which touch the $y$-axis and pass through the points $(1,2)$ and $(2,1)$.
\item A circle touches the line $3x+4y-10=0$ and has its centre at $(1,1)$. Find its equation.
\item Find the equations of the circles passing through $(0,2)$ and $(4,0)$ with centre on the line $2x-y=0$.
\item Find the equations of the circles touching both axes and the line $3x+4y=12$. (Hint: centre $(r,r)$ in the first quadrant.)
\end{enumerate}
\end{exercice}

\begin{corrige}[Exercise 1]
Touching the $y$-axis: $r=|a|$, $r^2=a^2$. Equal distances from $(1,2)$ and $(2,1)$: $(a-1)^2+(b-2)^2=(a-2)^2+(b-1)^2$ gives $a=b$. Then $(a-1)^2+(a-2)^2=a^2$ gives $a^2-6a+5=0$, so $a=1$ or $a=5$. Circles: $(x-1)^2+(y-1)^2=1$ and $(x-5)^2+(y-5)^2=25$.
\end{corrige}

\begin{corrige}[Exercise 2]
$r=\dfrac{|3+4-10|}{5}=\dfrac{3}{5}$. Equation: $(x-1)^2+(y-1)^2=\dfrac{9}{25}$.
\end{corrige}

\begin{corrige}[Exercise 3]
Centre $(a,b)$ with $b=2a$. Equal distances: $(a)^2+(b-2)^2=(a-4)^2+b^2$ gives $8a-4b=12$, i.e. $2a-b=3$. With $b=2a$: $0=3$ — no solution? Recheck: $a^2+(2a-2)^2=(a-4)^2+4a^2$ gives $a^2+4a^2-8a+4=5a^2-8a+16$, i.e. $4=16$, contradiction. So the line $2x-y=0$ is parallel to the perpendicular bisector of the segment: \textbf{no such circle exists}. (A good reminder that a system may be inconsistent.)
\end{corrige}

\begin{corrige}[Exercise 4]
Centre $(r,r)$, touching $3x+4y-12=0$: $\dfrac{|3r+4r-12|}{5}=r$, so $|7r-12|=5r$. Case $7r-12=5r$: $r=6$. Case $7r-12=-5r$: $r=1$. Circles: $(x-1)^2+(y-1)^2=1$ and $(x-6)^2+(y-6)^2=36$.
\end{corrige}

\begin{aretenir}[Key points of Section 8]
\begin{itemize}
\item Touching both axes: centre $(\pm r,\pm r)$, equation $(x\mp r)^2+(y\mp r)^2=r^2$.
\item Touching the $x$-axis: $r=|b|$; touching the $y$-axis: $r=|a|$.
\item Touching any line: perpendicular distance from centre to line equals $r$.
\item Centre equidistant from two points $\Longleftrightarrow$ centre on the perpendicular bisector.
\item Three constraints determine a circle; translate each into an equation in $a$, $b$, $r$, solve, and keep every valid solution.
\end{itemize}
\end{aretenir}

\coursec{Section 9: Synthesis, GCE Techniques and Revision}

In Section 8 we mastered circles constrained to touch axes or lines, and circles passing through two points with a tangency condition. That was the last new technique of the chapter. We now step back and look at the whole landscape. The GCE examiner does not test twenty-two separate lessons; he or she tests whether you can \emph{recognise a situation and choose the right tool in under fifteen minutes}. This final section is therefore a revision workshop: a map of all the equation forms, a decision strategy, combined multi-part questions in the GCE style, accuracy drills, and a self-assessment checklist.

\courssub{The Complete Map of Circle Equations}

Every equation of a circle you have met is one of five forms. Each form exists because it matches a particular type of given data. Learning the chapter means learning the \emph{correspondence between data and form}, not memorising formulas in isolation.

\begin{propriete}[Quick revision table of formulas]
\begin{center}
\begin{tabularx}{\linewidth}{|l|X|X|}
\hline
\rowcolor{popblueL} \textbf{Situation} & \textbf{Form / Formula} & \textbf{Key facts} \\
\hline
Centre $C(a,b)$, radius $r$ & $(x-a)^2+(y-b)^2=r^2$ & Read centre and radius directly \\
\hline
General form & $x^2+y^2+2gx+2fy+c=0$ & Centre $(-g,-f)$; $r=\sqrt{g^2+f^2-c}$; real circle iff $g^2+f^2-c>0$ \\
\hline
Diameter $A(x_1,y_1)$, $B(x_2,y_2)$ & $(x-x_1)(x-x_2)+(y-y_1)(y-y_2)=0$ & From $\vec{PA}\cdot\vec{PB}=0$ \\
\hline
Parametric & $x=a+r\cos\theta,\ y=b+r\sin\theta$ & Eliminate $\theta$ via $\cos^2\theta+\sin^2\theta=1$ \\
\hline
Tangent at $(x_1,y_1)$ on circle & $xx_1+yy_1+g(x+x_1)+f(y+y_1)+c=0$ & ``Replace'' rule; perpendicular to radius \\
\hline
Tangent length from $P(x_1,y_1)$ & $\ell=\sqrt{x_1^2+y_1^2+2gx_1+2fy_1+c}$ & $\ell^2=PC^2-r^2$ \\
\hline
Family through intersection of $S_1=0$, $S_2=0$ & $S_1+\lambda S_2=0$ & One parameter $\lambda$ fixed by one extra condition \\
\hline
Family through intersection of circle $S=0$ and line $L=0$ & $S+\lambda L=0$ & Same idea, line replaces second circle \\
\hline
Radical axis of $S_1=0$, $S_2=0$ & $S_1-S_2=0$ & Line of equal tangent lengths \\
\hline
Two circles touch & $d=r_1+r_2$ (external) or $d=|r_1-r_2|$ (internal) & $d$ = distance between centres \\
\hline
Line tangent to circle & distance from centre to line $=r$ & Or discriminant $=0$ \\
\hline
\end{tabularx}
\end{center}
\end{propriete}

\begin{popfigure}
\begin{center}
\begin{tikzpicture}[scale=0.95, every node/.style={font=\small}]
\node[draw=popdark, fill=popblueL, rounded corners, thick, minimum width=2.6cm, minimum height=1cm, align=center] (C) at (0,0) {\textbf{Circle}};
\node[draw=popblue, rounded corners, fill=white, align=center, text width=2.9cm] (std) at (-5.2,2.6) {Standard form\\ $(x-a)^2+(y-b)^2=r^2$};
\node[draw=popblue, rounded corners, fill=white, align=center, text width=2.9cm] (gen) at (0,3.4) {General form\\ $x^2+y^2+2gx+2fy+c=0$};
\node[draw=popblue, rounded corners, fill=white, align=center, text width=2.9cm] (dia) at (5.2,2.6) {Diameter form\\ $(x-x_1)(x-x_2)+\cdots=0$};
\node[draw=popgreen, rounded corners, fill=white, align=center, text width=2.9cm] (par) at (-5.2,-2.6) {Parametric\\ $x=a+r\cos\theta$\\ $y=b+r\sin\theta$};
\node[draw=popgreen, rounded corners, fill=white, align=center, text width=2.9cm] (fam) at (0,-3.4) {Families\\ $S_1+\lambda S_2=0$,\ $S+\lambda L=0$};
\node[draw=poporange, rounded corners, fill=white, align=center, text width=3.1cm] (tan) at (5.2,-2.6) {Tangents \& radical axis\\ distance $=r$,\ $S_1-S_2=0$};
\draw[thick, popmuted] (C) -- (std);
\draw[thick, popmuted] (C) -- (gen);
\draw[thick, popmuted] (C) -- (dia);
\draw[thick, popmuted] (C) -- (par);
\draw[thick, popmuted] (C) -- (fam);
\draw[thick, popmuted] (C) -- (tan);
\end{tikzpicture}
\end{center}
\legende{Concept map: the five equation forms and the main problem types of the chapter.}
\end{popfigure}

\courssub{Choosing the Right Form: A Decision Strategy}

Faced with a GCE question, spend the first thirty seconds classifying the data. The data \emph{always} tell you which form to use.

\begin{methode}[Choosing the right form]
\begin{enumerate}
\item \textbf{Centre and radius given (or easily found)?} Use the standard form $(x-a)^2+(y-b)^2=r^2$.
\item \textbf{Three points on the circle?} Use the general form $x^2+y^2+2gx+2fy+c=0$ and solve a $3\times 3$ system for $g,f,c$.
\item \textbf{Two endpoints of a diameter?} Use the diameter form directly.
\item \textbf{A tangency condition?} Translate it as ``perpendicular distance from centre to line equals radius'', or use the tangent formula at a known point of contact.
\item \textbf{Circle through the intersection of two circles, or of a circle and a line?} Write the family $S_1+\lambda S_2=0$ (or $S+\lambda L=0$) and determine $\lambda$ from the extra condition.
\item \textbf{Two circles touching?} Compute centres and radii, then impose $d=r_1+r_2$ or $d=|r_1-r_2|$.
\item \textbf{Always finish by checking} that your final circle satisfies every condition stated in the question.
\end{enumerate}
\end{methode}

\begin{popfigure}
\begin{center}
\begin{tikzpicture}[node distance=0.55cm, every node/.style={font=\small},
start/.style={draw=popdark, fill=popblueL, rounded corners, thick, align=center, text width=4.2cm},
dec/.style={draw=poporange, fill=white, rounded corners, align=center, text width=3.6cm},
res/.style={draw=popgreen, fill=white, rounded corners, align=center, text width=3.4cm}]
\node[start] (q) {What data does the question give?};
\node[dec, below left=0.7cm and -1.2cm of q] (d1) {Centre + radius, or diameter?};
\node[dec, below=0.7cm of q] (d2) {Three points on circle?};
\node[dec, below right=0.7cm and -1.2cm of q] (d3) {Tangency or intersection?};
\node[res, below=0.6cm of d1] (r1) {Standard form or diameter form};
\node[res, below=0.6cm of d2] (r2) {General form; solve $3\times3$ system};
\node[res, below=0.6cm of d3] (r3) {distance $=r$, or family $S_1+\lambda S_2=0$};
\draw[->, thick, popmuted] (q) -- (d1);
\draw[->, thick, popmuted] (q) -- (d2);
\draw[->, thick, popmuted] (q) -- (d3);
\draw[->, thick, popmuted] (d1) -- (r1);
\draw[->, thick, popmuted] (d2) -- (r2);
\draw[->, thick, popmuted] (d3) -- (r3);
\end{tikzpicture}
\end{center}
\legende{Decision tree: from the given data to the appropriate equation or method.}
\end{popfigure}

\courssub{Combined GCE-Style Problems}

Real GCE questions combine several lessons in one multi-part item. The parts are linked: a result from part (a) is usually needed in part (b), and so on. Here is a typical Paper 2 structure.

\begin{exoresolu}[Combined question: tangent, intersection and family]
The circle $C_1$ has equation $x^2+y^2-4x-6y+9=0$.
\begin{enumerate}
\item[(a)] Find the centre and radius of $C_1$.
\item[(b)] Show that the line $y=2x$ is a tangent to $C_1$, and find the point of contact.
\item[(c)] The circle $C_2$ has equation $x^2+y^2-10x-2y+21=0$. Find the equation of the circle passing through the points of intersection of $C_1$ and $C_2$ and through the origin.
\end{enumerate}
\tcblower
\textbf{(a)} Complete the square:
\begin{align*}
x^2-4x+y^2-6y+9&=0\\
(x-2)^2-4+(y-3)^2-9+9&=0\\
(x-2)^2+(y-3)^2&=4.
\end{align*}
Centre $A(2,3)$, radius $r_1=2$.

\textbf{(b)} Distance from $A(2,3)$ to the line $2x-y=0$:
\[
d=\frac{|2(2)-3|}{\sqrt{2^2+(-1)^2}}=\frac{1}{\sqrt5}.
\]
This is not equal to $2$, so let us instead solve simultaneously. Substituting $y=2x$ into $C_1$:
\begin{align*}
x^2+4x^2-4x-12x+9&=0\\
5x^2-16x+9&=0\\
\Delta&=256-180=76>0.
\end{align*}
Since $\Delta>0$ the line \emph{cuts} the circle at two points; it is \textbf{not} a tangent. (This is exactly the kind of verification the examiner rewards: never assume.) The intersection points are $x=\dfrac{16\pm\sqrt{76}}{10}=\dfrac{8\pm\sqrt{19}}{5}$, with $y=2x$.

\textbf{(c)} The family through the intersections is
\[
S_1+\lambda S_2=0:\quad (x^2+y^2-4x-6y+9)+\lambda(x^2+y^2-10x-2y+21)=0.
\]
Passing through $(0,0)$: $9+21\lambda=0$, so $\lambda=-\dfrac{3}{7}$. Substituting and multiplying by $7$:
\begin{align*}
7(x^2+y^2-4x-6y+9)-3(x^2+y^2-10x-2y+21)&=0\\
4x^2+4y^2+2x-36y&=0\\
2x^2+2y^2+x-18y&=0.
\end{align*}
Check: $(0,0)$ gives $0$ \checkmark. The required circle is $2x^2+2y^2+x-18y=0$.
\end{exoresolu}

\begin{attention}[Verify every condition]
A very common error: after finding $\lambda$ in a family question, candidates forget to check that the resulting circle passes through the required point, or that a ``tangent'' really touches. In the example above, part (b) shows why verification matters — the line looked plausible but was a secant. Before moving to the next part, substitute and confirm \emph{all} the given conditions.
\end{attention}

\courssub{Algebraic Accuracy Under Exam Timing}

Most marks lost on circle questions are \emph{algebra} marks, not geometry marks. Train these three skills until they are automatic.

\textbf{1. Completing the square.} $x^2+2gx=(x+g)^2-g^2$. Practise with awkward coefficients: $x^2+5x=\left(x+\tfrac52\right)^2-\tfrac{25}{4}$. Fractions are where errors hide.

\textbf{2. Solving $3\times3$ systems.} When three points give three equations in $g,f,c$, subtract pairs of equations first to eliminate the constant $c$; this reduces the problem to a $2\times2$ system. Write each step on a new line.

\textbf{3. Discriminant computations.} For tangency via substitution you need $b^2-4ac=0$ with messy coefficients. Compute $b^2$ and $4ac$ separately, then subtract — never do it in one mental step.

\begin{methode}[Exam timing for circle questions]
\begin{enumerate}
\item A full circle question (3--4 parts) should take \textbf{12 to 15 minutes}. Practise with a timer.
\item Budget roughly: 2 min to read, classify data and plan; 8--10 min of working; 2 min to verify and box answers.
\item If a part resists after 4 minutes, write down the method you \emph{would} use (method marks) and move on; return at the end.
\item Keep exact surd form ($\sqrt{13}$, $\tfrac{8\pm\sqrt{19}}{5}$) unless the question asks for decimals; then give 3 significant figures.
\end{enumerate}
\end{methode}

\courssub{Common GCE Question Patterns and Mark-Scheme Expectations}

Past papers show recurring patterns. Knowing them in advance converts surprise into routine.

\begin{center}
\begin{tabularx}{\linewidth}{|l|X|X|}
\hline
\rowcolor{popblueL} \textbf{Pattern} & \textbf{Typical wording} & \textbf{What the mark scheme expects} \\
\hline
Find the equation & ``centre $(2,-1)$, passes through $(5,3)$'' & Radius via distance formula; standard form; expansion if asked \\
\hline
Show that (tangency) & ``Show that $3x+4y=25$ touches the circle'' & Distance from centre $=r$, \emph{with every step shown}, or $\Delta=0$ \\
\hline
Three points & ``passes through $A$, $B$, $C$'' & Three substitutions, system solved, final equation \\
\hline
Two circles & ``find the radical axis'' / ``show the circles touch'' & $S_1-S_2=0$; or $d$ compared with $r_1\pm r_2$ \\
\hline
Family & ``circle through the intersections of \dots\ and the point \dots'' & $S_1+\lambda S_2=0$, value of $\lambda$, simplified answer \\
\hline
Locus & ``find the locus of $P$ such that \dots'' & Translate condition algebraically, recognise a circle, state centre and radius \\
\hline
\end{tabularx}
\end{center}

\begin{attention}[``Show that'' means show everything]
In a ``show that'' question the answer is given, so the marks are entirely in the \emph{working}. Every algebraic step must appear: the distance formula written out, the substitution, the simplification. Skipping from $5x^2-16x+9=0$ straight to a conclusion loses method marks even if your final line is correct. Write as if the examiner cannot read your mind — because he cannot.
\end{attention}

\begin{exemplebox}[Presentation: touches vs intersects]
To decide whether the line $y=x+1$ touches or cuts the circle $x^2+y^2-2x-4y+1=0$ (centre $(1,2)$, $r=2$):
\[
d=\frac{|1-2+1|}{\sqrt{1^2+(-1)^2}}=0<2.
\]
Since $d<r$, the line \emph{cuts} the circle at two points (here the line even passes through the centre, so it is a diameter line). State the comparison explicitly: ``$d=0<2=r$, hence the line intersects the circle at two distinct points.'' That sentence earns the conclusion mark.
\end{exemplebox}

\begin{savaistu}[Exact answers and the calculator]
GCE mark schemes accept exact surd answers such as $\sqrt{13}$ or $\dfrac{8+\sqrt{19}}{5}$, and often prefer them. Use your calculator to \emph{check} ($\sqrt{13}\approx 3.606$) but write the exact value as your final answer unless decimals are requested. A calculator slip in a decimal costs the accuracy mark; an exact surd cannot be ``wrongly rounded''.
\end{savaistu}

\courssub{Chapter Self-Assessment Checklist}

Tick each objective only when you can perform it \emph{without looking at your notes}, within exam timing.

\begin{aretenir}[Self-assessment checklist for Chapter PMM01]
\begin{itemize}
\item I can write the equation of a circle given its centre and radius, and expand it into general form.
\item I can find the centre and radius from $x^2+y^2+2gx+2fy+c=0$ by completing the square, and state when the equation represents a real circle.
\item I can write the equation of a circle on a given diameter.
\item I can find the equation of a circle through three non-collinear points.
\item I can find the equation of the tangent at a point on a circle, and the tangents of a given gradient.
\item I can test whether a line is tangent to a circle (distance $=r$ or $\Delta=0$) and find intersection points of a line and a circle.
\item I can convert between parametric and cartesian forms of a circle.
\item I can compute the length of a tangent from an external point.
\item I can use the families $S_1+\lambda S_2=0$ and $S+\lambda L=0$ and determine $\lambda$.
\item I can find the intersection points and the radical axis of two circles, and test whether two circles touch.
\item I can handle constrained circles: touching axes, touching a line, centre on a line, through two points with tangency.
\item I can solve locus problems leading to circles.
\item I verify every condition and present ``show that'' working in full.
\end{itemize}
\end{aretenir}

\begin{exercice}[Timed revision test — 30 minutes]
\begin{enumerate}
\item Find the equation of the circle with centre $(-1,4)$ passing through $(2,0)$.
\item Find the centre and radius of $x^2+y^2+6x-10y+18=0$.
\item Show that the line $x+y=7$ is a tangent to the circle $x^2+y^2-4x-6y+11=0$, and find the point of contact.
\item Find the equation of the circle through the intersections of $x^2+y^2=9$ and $x^2+y^2-6x-8y+9=0$ which also passes through $(1,1)$.
\item Find the radical axis of the two circles in question 4 and the length of the tangent from $(5,2)$ to the circle $x^2+y^2=9$.
\end{enumerate}
\end{exercice}

\begin{corrige}[Timed revision test]
\textbf{1.} $r^2=(2-(-1))^2+(0-4)^2=9+16=25$, so $(x+1)^2+(y-4)^2=25$, i.e. $x^2+y^2+2x-8y-8=0$.

\textbf{2.} $(x+3)^2+(y-5)^2=-18+9+25=16$: centre $(-3,5)$, radius $4$.

\textbf{3.} Circle: $x^2+y^2-4x-6y+11=0$. Centre $(2,3)$, $r^2=4+9-11=2$, so $r=\sqrt2$. Distance from $(2,3)$ to $x+y-7=0$: $\dfrac{|2+3-7|}{\sqrt2}=\dfrac{2}{\sqrt2}=\sqrt2 = r$. Hence the line is a tangent. Point of contact: substitute $y=7-x$ into the circle:
\begin{align*}
x^2+(7-x)^2-4x-6(7-x)+11&=0\\
2x^2-12x+18&=0\\
x^2-6x+9&=0\\
(x-3)^2&=0 \Rightarrow x=3,\ y=4.
\end{align*}
Point of contact is $(3,4)$.

\textbf{4.} Let $S_1=x^2+y^2-9=0$, $S_2=x^2+y^2-6x-8y+9=0$. Family: $S_1+\lambda S_2=0$. At $(1,1)$: $S_1(1,1)=-7$, $S_2(1,1)=1+1-6-8+9=-3$. So $-7-3\lambda=0 \Rightarrow \lambda=-\frac73$. Multiply by $3$:
\[
3(x^2+y^2-9)-7(x^2+y^2-6x-8y+9)=0
\]
\[
-4x^2-4y^2+42x+56y-90=0 \quad\Rightarrow\quad 4x^2+4y^2-42x-56y+90=0.
\]
Divide by $2$: $2x^2+2y^2-21x-28y+45=0$.

\textbf{5.} Radical axis: $S_1-S_2=0$ gives $6x+8y-18=0$, i.e. $3x+4y-9=0$. Tangent length from $(5,2)$ to $x^2+y^2-9=0$: $\sqrt{25+4-9}=\sqrt{20}=2\sqrt5$.
\end{corrige}

You have now completed the chapter on Circle Geometry — from the definition of a circle as a locus, through every equation form, tangents, families, radical axes and constrained circles, to full GCE technique. Keep the concept map and the decision tree beside you as you revise, redo the timed test after one week, and approach the examination with confidence: every circle question is a combination of tools you now own.

\newpage

\coursec{Integration activity}

\coursec{Synthesis and Applications}

\courssub{Aims of the Activity}

This integration activity closes the chapter \emph{Circle Geometry} by asking you to mobilise, in one single authentic modelling task, almost everything you have learnt:

\begin{itemize}
\item finding the equation of a circle through three given points;
\item using the general equation $x^{2}+y^{2}+2gx+2fy+c=0$, its centre $(-g,-f)$ and radius $\sqrt{g^{2}+f^{2}-c}$;
\item deciding the position of a line relative to a circle and finding the points of intersection of a line and a circle;
\item computing the length of a chord;
\item finding the equation of a circle of given radius tangent to a given line, with its centre on another line;
\item determining the radical axis of two circles and giving it a physical meaning;
\item communicating a complete mathematical solution, with a clear sketch, as required at the GCE Advanced Level.
\end{itemize}

You will work in groups of three or four. Each group must hand in a written report containing the full working, a neat, accurately scaled diagram, and a short interpretation of each result in the language of the real situation.

\courssub{Situation}

The Bamenda City Council is modernising the Up Station neighbourhood. A simplified map of the area is drawn on a coordinate grid in which \textbf{one unit represents \SI{100}{m}} on the ground, the $x$-axis pointing East and the $y$-axis pointing North.

Three schools must receive the signal of a new FM relay antenna:
\begin{center}
\begin{tabularx}{\linewidth}{|l|X|X|X|}
\hline
\rowcolor{popblueL} School & Government School Azire & GBHS Bamenda & Presbyterian School Nkwen \\
\hline Position & $A(1,\,2)$ & $B(7,\,4)$ & $C(3,\,8)$ \\
\hline
\end{tabularx}
\end{center}

The City Council's engineering team faces four problems.

\textbf{Problem 1 --- Coverage zone.} The relay antenna must cover all three schools with the \emph{smallest possible circular zone}. The engineers assume that this zone is the circle passing through $A$, $B$ and $C$. They need its equation, the coordinates of the ideal position of the antenna (the centre), and the radius of coverage in metres.

\textbf{Problem 2 --- The new road.} A straight road is planned along the line
\[ x-y+3=0. \]
The Council wants to know whether this road crosses the coverage zone; if it does, the entry and exit points must be marked on the map, and the length of road lying inside the zone (a chord of the circle) must be computed in metres, since special signal-boosting cables will be laid along that stretch.

\textbf{Problem 3 --- The roundabout.} To ease traffic, a circular roundabout of radius exactly $1$ unit (\SI{100}{m}) is to be built so that:
\begin{itemize}
\item the road $x-y+3=0$ is \emph{tangent} to the roundabout (the road just touches the circular island);
\item the centre of the roundabout lies on the east--west avenue of equation $y=2$.
\end{itemize}
The engineers suspect there may be \emph{two} possible positions. They want the equations of all possible roundabouts.

\textbf{Problem 4 --- A competing antenna.} A private operator already covers the zone
\[ x^{2}+y^{2}-10x-4y+20=0. \]
Along the \emph{radical axis} of the two circles, the signals of the two antennas have equal strength (equal tangent lengths model equal ``power reach''). The Council wants the equation of this line, so as to know where listeners could receive both stations equally well.

\courssub{Tasks}

\begin{enumerate}
\item \textbf{Setting up the coverage circle.} Let the circle through $A(1,2)$, $B(7,4)$, $C(3,8)$ have equation $x^{2}+y^{2}+2gx+2fy+c=0$. Substitute the three points to obtain three linear equations in $g$, $f$ and $c$. Solve the system.
\item \textbf{Centre and radius.} Deduce the centre $(-g,-f)$ and the radius $r=\sqrt{g^{2}+f^{2}-c}$. Give the radius in metres and verify that $A$, $B$ and $C$ are all at distance $r$ from the centre.
\item \textbf{Position of the road.} Compute the perpendicular distance from the centre of the coverage circle to the line $x-y+3=0$. Compare it with the radius and conclude whether the road crosses the zone.
\item \textbf{Entry and exit points.} Solve simultaneously the equations of the road and the circle to find the coordinates of the entry point $P$ and the exit point $Q$.
\item \textbf{Length of road inside the zone.} Compute $|PQ|$ in units, then in metres. Check your answer with the chord formula $|PQ|=2\sqrt{r^{2}-d^{2}}$, where $d$ is the distance from the centre to the line.
\item \textbf{The roundabout.} Let the centre of the roundabout be $(h,2)$. Write down the condition ``the distance from $(h,2)$ to the line $x-y+3=0$ equals $1$''. Solve for $h$ (you should get two values) and write the equation of each possible roundabout.
\item \textbf{Radical axis.} Subtract the equations of the two coverage circles to obtain the radical axis. State its equation and interpret it.
\item \textbf{Sketch and presentation.} On one accurate diagram, plot the three schools, the coverage circle, the road, the chord $PQ$, the two possible roundabouts, the second coverage circle and the radical axis. Prepare a five-minute presentation justifying each step.
\end{enumerate}

\begin{attention}[Before you start]
Remember the grid scale: $1$ unit $=\SI{100}{m}$, so an area of $1$ square unit is $\SI{10000}{m^{2}}$ and a length of $2.5$ units is $\SI{250}{m}$. Do not mix units and metres in your final answers. Also, when solving $|h-2+3|=\sqrt{2}$ in Task 6, do not forget the \emph{two} signs of the absolute value --- that is precisely where the two roundabouts come from.
\end{attention}

\courssub{Model Answer}

\textbf{Task 1 --- Equation of the coverage circle.}\par

Substituting $A(1,2)$, $B(7,4)$, $C(3,8)$ into $x^{2}+y^{2}+2gx+2fy+c=0$:
\begin{align*}
A:&\quad 1+4+2g+4f+c=0 &&\Rightarrow\ 2g+4f+c=-5,\\
B:&\quad 49+16+14g+8f+c=0 &&\Rightarrow\ 14g+8f+c=-65,\\
C:&\quad 9+64+6g+16f+c=0 &&\Rightarrow\ 6g+16f+c=-73.
\end{align*}
Subtracting the first equation from the second: $12g+4f=-60$, i.e.\ $3g+f=-15$. Subtracting the first from the third: $4g+12f=-68$, i.e.\ $g+3f=-17$. Solving
\[ \begin{cases} 3g+f=-15\\ g+3f=-17 \end{cases} \]
gives $f=-15-3g$, then $g+3(-15-3g)=-17$, so $-8g=28$, $g=-\dfrac{7}{2}$ and $f=-15+\dfrac{21}{2}=-\dfrac{9}{2}$. Finally $c=-5-2g-4f=-5+7+18=20$.

\begin{aretenir}[Coverage zone]
\[ x^{2}+y^{2}-7x-9y+20=0. \]
\end{aretenir}

\textbf{Task 2 --- Centre and radius.}\par

Centre: $\Omega=\left(\dfrac{7}{2},\dfrac{9}{2}\right)$. Radius:
\[ r=\sqrt{g^{2}+f^{2}-c}=\sqrt{\tfrac{49}{4}+\tfrac{81}{4}-20}=\sqrt{\tfrac{130-80}{4}}=\sqrt{\tfrac{25}{2}}=\dfrac{5\sqrt{2}}{2}\approx 3.54. \]
So the antenna stands at $(3.5,\,4.5)$ and the coverage radius is $\dfrac{5\sqrt{2}}{2}\times 100\approx \SI{354}{m}$.

\emph{Check:} $|\Omega A|^{2}=\left(1-\tfrac72\right)^{2}+\left(2-\tfrac92\right)^{2}=\tfrac{25}{4}+\tfrac{25}{4}=\tfrac{25}{2}=r^{2}$. $\checkmark$ (Similar checks hold for $B$ and $C$.)

\textbf{Task 3 --- Does the road cross the zone?}\par

Distance from $\Omega\left(\tfrac72,\tfrac92\right)$ to the line $x-y+3=0$:
\[ d=\dfrac{\left|\tfrac72-\tfrac92+3\right|}{\sqrt{1^{2}+(-1)^{2}}}=\dfrac{|2|}{\sqrt{2}}=\sqrt{2}\approx 1.41. \]
Since $d=\sqrt{2}<r=\dfrac{5\sqrt2}{2}$, the road \textbf{crosses} the coverage zone in two points.

\textbf{Task 4 --- Entry and exit points.}\par

From the road equation, $y=x+3$. Substituting into the circle:
\begin{align*}
x^{2}+(x+3)^{2}-7x-9(x+3)+20&=0\\
2x^{2}+6x+9-7x-9x-27+20&=0\\
2x^{2}-10x+2&=0\\
x^{2}-5x+1&=0.
\end{align*}
Hence $x=\dfrac{5\pm\sqrt{25-4}}{2}=\dfrac{5\pm\sqrt{21}}{2}$, and $y=x+3=\dfrac{11\pm\sqrt{21}}{2}$.

Entry and exit points (with $\sqrt{21}\approx 4.58$):
\[ P\left(\tfrac{5-\sqrt{21}}{2},\,\tfrac{11-\sqrt{21}}{2}\right)\approx(0.21,\,3.21),\qquad Q\left(\tfrac{5+\sqrt{21}}{2},\,\tfrac{11+\sqrt{21}}{2}\right)\approx(4.79,\,7.79). \]

\textbf{Task 5 --- Length of road inside the zone.}\par

\[ |PQ|=\sqrt{\left(\sqrt{21}\right)^{2}+\left(\sqrt{21}\right)^{2}}=\sqrt{42}\approx 6.48\ \text{units}. \]
\emph{Check with the chord formula:} $|PQ|=2\sqrt{r^{2}-d^{2}}=2\sqrt{\tfrac{25}{2}-2}=2\sqrt{\tfrac{21}{2}}=\sqrt{42}$. $\checkmark$

The stretch of road inside the coverage zone measures $\sqrt{42}\times 100\approx \SI{648}{m}$.

\textbf{Task 6 --- The two possible roundabouts.}\par

Let the centre be $(h,2)$ and the radius $1$. Tangency to the road requires
\[ \dfrac{|h-2+3|}{\sqrt{2}}=1 \quad\Longrightarrow\quad |h+1|=\sqrt{2}. \]
So $h+1=\sqrt{2}$ or $h+1=-\sqrt{2}$, giving $h=-1+\sqrt{2}\approx 0.41$ or $h=-1-\sqrt{2}\approx-2.41$.

\begin{aretenir}[The two roundabouts]
\[ \bigl(x+1-\sqrt{2}\bigr)^{2}+(y-2)^{2}=1 \qquad\text{and}\qquad \bigl(x+1+\sqrt{2}\bigr)^{2}+(y-2)^{2}=1. \]
\end{aretenir}

Both centres lie on the avenue $y=2$, on opposite sides of the road, exactly \SI{100}{m} from it: two valid engineering options.

\textbf{Task 7 --- Radical axis.}\par

Subtracting the second circle from the coverage circle:
\[ \left(x^{2}+y^{2}-7x-9y+20\right)-\left(x^{2}+y^{2}-10x-4y+20\right)=0 \]
\[ 3x-5y=0 \qquad\Longleftrightarrow\qquad y=\dfrac{3}{5}x. \]

\begin{savaistu}[Interpretation]
Along the line $3x-5y=0$, every point has equal tangent lengths to the two circles --- the engineers read this as \emph{equal signal reach} from both antennas. Listeners living along this straight avenue (which passes through the origin of the grid) can receive both stations equally well. Note that the second circle has centre $(5,2)$ and radius $\sqrt{25+4-20}=3$, and the radical axis is perpendicular to the line of centres, as the theory predicts.
\end{savaistu}

\textbf{Task 8 --- Sketch.}\par

\begin{popfigure}
\begin{center}
\begin{tikzpicture}[scale=0.62]
\draw[gray!30, very thin] (-3.5,-0.5) grid (8.5,9.5);
\draw[->, thick] (-3.5,0)--(8.5,0) node[below left]{$x$};
\draw[->, thick] (0,-0.5)--(0,9.5) node[below left]{$y$};
% coverage circle
\draw[popblue, thick] (3.5,4.5) circle (3.5355);
\fill[popblue] (3.5,4.5) circle (2.5pt) node[below right]{$\Omega$};
% schools
\fill[popdark] (1,2) circle (2.5pt) node[below left]{$A$};
\fill[popdark] (7,4) circle (2.5pt) node[below right]{$B$};
\fill[popdark] (3,8) circle (2.5pt) node[above left]{$C$};
% road
\draw[poporange, thick] (-3.2,-0.2)--(6.2,9.2) node[above left, align=center, text width=1.6cm]{road $x-y+3=0$};
% chord endpoints
\fill[poporange] (0.2087,3.2087) circle (2.5pt) node[left]{$P$};
\fill[poporange] (4.7913,7.7913) circle (2.5pt) node[right]{$Q$};
% roundabouts
\draw[popgreen, thick] (0.4142,2) circle (1);
\draw[popgreen, thick] (-2.4142,2) circle (1);
\fill[popgreen] (0.4142,2) circle (2pt);
\fill[popgreen] (-2.4142,2) circle (2pt);
% avenue y=2
\draw[popmuted, dashed] (-3.5,2)--(8.5,2) node[right]{$y=2$};
% second circle
\draw[poppurple, thick] (5,2) circle (3);
\fill[poppurple] (5,2) circle (2.5pt) node[below]{$\Omega'$};
% radical axis
\draw[popink, thick, dashed] (-0.8,-0.48)--(8.5,5.1) node[above left, align=center, text width=1.8cm]{radical axis $3x-5y=0$};
\end{tikzpicture}
\end{center}
\legende{The Bamenda modelling map: coverage circle (blue), road and chord $PQ$ (orange), the two possible roundabouts (green), second antenna (purple) and radical axis (dashed).}
\end{popfigure}

\begin{methode}[What the examiners expect in such a synthesis task]
\begin{enumerate}
\item \textbf{Model} the situation: translate each real object (antenna zone, road, roundabout) into a circle or a line.
\item \textbf{Compute} with the standard tools: system for three points, distance from a point to a line, simultaneous equations, radical axis by subtraction.
\item \textbf{Verify} each result by an independent method (e.g.\ chord length by two routes).
\item \textbf{Interpret}: every final answer must be given in the language of the situation, with correct units (metres, not grid units).
\end{enumerate}
\end{methode}

\begin{exercice}[Going further]
\begin{enumerate}
\item Show that the road $x-y+3=0$ is \emph{not} tangent to the second coverage circle $x^{2}+y^{2}-10x-4y+20=0$, and state whether it cuts it or misses it entirely.
\item Find the length of the tangent from the school $A(1,2)$ to the second coverage circle, in metres.
\item A third roundabout, also of radius $1$, must be tangent to the road and have its centre on the line $y=x$. Show that this is impossible, and explain geometrically why.
\item The Council finally decides that the antenna should be placed so as to minimise the \emph{sum of squares} of distances to the three schools. Show that this new criterion places the antenna at the centroid $G\left(\tfrac{11}{3},\tfrac{14}{3}\right)$ of triangle $ABC$, and compare with $\Omega$.
\end{enumerate}
\end{exercice}

\newpage

\coursec{Methods and worked examples}

\coursec{Section 1: Equations of a Circle — Centre, Radius, Diameter}

\begin{definition}[Circle]
A \cle{circle} is the set of all points in a plane that are at a fixed distance (the \cle{radius} $r>0$) from a fixed point (the \cle{centre} $C(a,b)$).
\end{definition}

\begin{propriete}[Standard equation of a circle]
The circle with centre $C(a,b)$ and radius $r$ has equation
\[
(x-a)^2+(y-b)^2=r^2.
\]
If the centre is the origin $O(0,0)$, the equation reduces to $x^2+y^2=r^2$.
\end{propriete}

\begin{popfigure}
\begin{tikzpicture}[scale=0.8]
\draw[popline, thick] (0,0) circle (2.5);
\draw[popblue, thick, ->] (0,0) -- (2.5,0) node[midway, below] {$r$};
\draw[poporange, thick, ->] (0,0) -- (1.5,1.5) node[midway, above left] {$r$};
\fill[popdark] (0,0) circle (2pt) node[below left] {$C(a,b)$};
\fill[popdark] (2.5,0) circle (2pt) node[below right] {$P(x,y)$};
\draw[dashed, popmuted] (0,0) -- (1.5,1.5);
\draw[dashed, popmuted] (1.5,1.5) -- (1.5,0);
\draw[dashed, popmuted] (1.5,0) -- (0,0);
\node at (0.75, -0.3) {$x-a$};
\node at (1.7, 0.75) {$y-b$};
\end{tikzpicture}
\legende{The standard circle: $(x-a)^2+(y-b)^2=r^2$}
\end{popfigure}

\begin{methode}[Equation of a circle from centre and radius]
\begin{enumerate}
\item Identify the \cle{centre} $(a,b)$ and the \cle{radius} $r$.
\item If the radius is not given directly, compute it as the distance from the centre to a known point on the circle:
      $r=\sqrt{(x_1-a)^2+(y_1-b)^2}$.
\item Substitute into the standard form $(x-a)^2+(y-b)^2=r^2$.
\item If the \cle{general form} is required, expand and simplify to $x^2+y^2+2gx+2fy+c=0$.
\item \textbf{Reflex:} if the centre is the origin, the equation reduces to $x^2+y^2=r^2$. If the circle is given \emph{on a diameter} with endpoints $A(x_1,y_1)$ and $B(x_2,y_2)$, use directly
      \[
      (x-x_1)(x-x_2)+(y-y_1)(y-y_2)=0.
      \]
\end{enumerate}
\end{methode}

\begin{exoresolu}[Circle on a given diameter]
Find the equation of the circle having the segment joining $A(1,4)$ and $B(5,-2)$ as a diameter. Give your answer in general form.
\tcblower
\textbf{Method A (diameter form).} Since $AB$ is a diameter, for any point $P(x,y)$ on the circle, $\angle APB=90^\circ$, so $\vec{PA}\cdot\vec{PB}=0$:
\[
(x-1)(x-5)+(y-4)(y+2)=0.
\]
Expanding:
\begin{align*}
x^2-6x+5+y^2-2y-8&=0\\
x^2+y^2-6x-2y-3&=0.
\end{align*}

\textbf{Method B (check via centre and radius).} The centre is the midpoint of $AB$:
\[
C=\left(\frac{1+5}{2},\frac{4+(-2)}{2}\right)=(3,1).
\]
The radius is half of $AB$:
\[
AB=\sqrt{(5-1)^2+(-2-4)^2}=\sqrt{16+36}=\sqrt{52},\qquad r=\frac{\sqrt{52}}{2}=\sqrt{13}.
\]
Hence $(x-3)^2+(y-1)^2=13$, which expands to $x^2+y^2-6x-2y-3=0$. \checkmark\ Both methods agree.
\end{exoresolu}

\begin{attention}[Common mistake]
When using the diameter form, do not forget the \textbf{plus} sign between the $x$-part and the $y$-part: $(x-x_1)(x-x_2)+(y-y_1)(y-y_2)=0$. A minus sign would give a hyperbola.
\end{attention}

\coursec{Section 2: The General Equation and Circles Through Three Points}

\begin{propriete}[General equation of a circle]
Any equation of the form
\[
x^2+y^2+2gx+2fy+c=0
\]
represents a circle (possibly degenerate) provided the coefficients of $x^2$ and $y^2$ are equal and non-zero, and there is no $xy$ term. Its centre is $(-g,-f)$ and its radius is $r=\sqrt{g^2+f^2-c}$.
\end{propriete}

\begin{methode}[Reading the general equation $x^2+y^2+2gx+2fy+c=0$]
\begin{enumerate}
\item Check that the coefficients of $x^2$ and $y^2$ are \textbf{equal} and that there is \textbf{no $xy$ term}; otherwise the equation does not represent a circle.
\item Divide through if necessary so that the coefficients of $x^2$ and $y^2$ are $1$.
\item Read off $g$ and $f$ as \textbf{half} the coefficients of $x$ and $y$ (watch the signs!).
\item Centre: $C=(-g,-f)$. Radius: $r=\sqrt{g^2+f^2-c}$.
\item \textbf{Reflex:} the circle is real only if $g^2+f^2-c>0$; if $g^2+f^2-c=0$ it reduces to a single point (point circle); if negative, no real circle exists.
\item Alternative: complete the square in $x$ and in $y$ separately.
\end{enumerate}
\end{methode}

\begin{exoresolu}[From general form to centre and radius]
Show that $2x^2+2y^2-8x+6y-3=0$ represents a circle, and find its centre and radius.
\tcblower
The coefficients of $x^2$ and $y^2$ are equal (both $2$) and there is no $xy$ term, so the equation may represent a circle. Divide by $2$:
\[
x^2+y^2-4x+3y-\tfrac{3}{2}=0.
\]
Here $2g=-4$, $2f=3$, $c=-\tfrac32$, so $g=-2$, $f=\tfrac32$.
\textbf{Centre:} $C=(-g,-f)=\left(2,-\tfrac32\right)$.
\textbf{Radius:}
\[
r=\sqrt{g^2+f^2-c}=\sqrt{4+\tfrac94+\tfrac32}=\sqrt{\tfrac{16+9+6}{4}}=\sqrt{\tfrac{31}{4}}=\frac{\sqrt{31}}{2}.
\]
Since $g^2+f^2-c=\tfrac{31}{4}>0$, the circle is real.
\textbf{Check by completing the square:}
\[
(x-2)^2-4+\left(y+\tfrac32\right)^2-\tfrac94-\tfrac32=0
\;\Longrightarrow\;(x-2)^2+\left(y+\tfrac32\right)^2=\tfrac{31}{4}.\ \checkmark
\]
\end{exoresolu}

\begin{methode}[Circle through three non-collinear points]
\begin{enumerate}
\item Write the general equation $x^2+y^2+2gx+2fy+c=0$.
\item Substitute each of the three points to obtain three linear equations in $g$, $f$, $c$.
\item Solve the system (elimination or substitution).
\item Substitute back and, if asked, state the centre $(-g,-f)$ and radius $\sqrt{g^2+f^2-c}$.
\item \textbf{Reflex:} if two of the points are endpoints of a diameter (check: is the third angle a right angle?), the diameter form is much faster. Always verify that the three points are not collinear.
\end{enumerate}
\end{methode}

\begin{exoresolu}[Circle through three points]
Find the equation of the circle passing through the points $A(0,0)$, $B(4,0)$ and $C(0,3)$.
\tcblower
Let the circle be $x^2+y^2+2gx+2fy+c=0$.

\textbf{Through $A(0,0)$:}\quad $c=0$.

\textbf{Through $B(4,0)$:}\quad $16+8g+c=0 \Rightarrow 16+8g=0 \Rightarrow g=-2$.

\textbf{Through $C(0,3)$:}\quad $9+6f+c=0 \Rightarrow 9+6f=0 \Rightarrow f=-\tfrac32$.

Hence the equation is
\[
x^2+y^2-4x-3y=0.
\]
\textbf{Centre:} $\left(2,\tfrac32\right)$. \textbf{Radius:} $r=\sqrt{4+\tfrac94-0}=\sqrt{\tfrac{25}{4}}=\tfrac52$.

\textbf{Remark.} Since $\angle BAC=90^\circ$ (the axes are perpendicular), $BC$ is a diameter. Check: midpoint of $BC$ is $\left(2,\tfrac32\right)$ \checkmark, and $BC=\sqrt{16+9}=5$, so $r=\tfrac52$ \checkmark. The diameter form gives $(x-4)(x-0)+(y-0)(y-3)=0$, i.e. $x^2+y^2-4x-3y=0$ — same result, faster.
\end{exoresolu}

\begin{savaistu}[Collinearity check]
Three points $(x_1,y_1)$, $(x_2,y_2)$, $(x_3,y_3)$ are collinear iff the determinant
\[
\begin{vmatrix}
x_1 & y_1 & 1 \\
x_2 & y_2 & 1 \\
x_3 & y_3 & 1
\end{vmatrix}=0.
\]
If they are collinear, no circle passes through them (unless we allow a line as a degenerate circle of infinite radius).
\end{savaistu}

\coursec{Section 3: Tangents to a Circle}

\begin{propriete}[Tangent to a circle]
A line is \cle{tangent} to a circle if it touches the circle at exactly one point. At the point of contact, the tangent is perpendicular to the radius.
\end{propriete}

\begin{methode}[Equation of a tangent to a circle]
\textbf{Case A: tangent at a point $P(x_1,y_1)$ on the circle.}
\begin{enumerate}
\item Find the centre $C(a,b)$.
\item Gradient of radius $CP$: $m_r=\dfrac{y_1-b}{x_1-a}$ (if $x_1\neq a$; if $x_1=a$, the radius is vertical and the tangent is horizontal).
\item Tangent is perpendicular to the radius: $m_t=-\dfrac{1}{m_r}$ (or $m_t=0$ if $m_r$ is undefined).
\item Use $y-y_1=m_t(x-x_1)$.
\item \textbf{Shortcut (T = 0):} for the circle $x^2+y^2+2gx+2fy+c=0$, the tangent at $(x_1,y_1)$ is
      \[
      xx_1+yy_1+g(x+x_1)+f(y+y_1)+c=0.
      \]
\end{enumerate}
\textbf{Case B: condition for the line $y=mx+c$ to touch the circle.}
\begin{enumerate}
\item The perpendicular distance from the centre to the line equals the radius:
      \[
      \frac{|ma-b+c|}{\sqrt{m^2+1}}=r.
      \]
\item \emph{Or} substitute the line into the circle and set the discriminant of the resulting quadratic to zero.
\end{enumerate}
\end{methode}

\begin{popfigure}
\begin{tikzpicture}[scale=0.8]
\draw[popline, thick] (0,0) circle (2);
\draw[popblue, thick, ->] (0,0) -- (1.6,1.2) node[midway, above left] {$r$};
\fill[popdark] (0,0) circle (2pt) node[below left] {$C$};
\fill[popdark] (1.6,1.2) circle (2pt) node[above right] {$P$};
\draw[poporange, thick] (-1,2.8) -- (3.5,-0.5) node[below right] {Tangent};
\draw[dashed, popmuted] (1.6,1.2) -- (1.6,0);
\draw[dashed, popmuted] (1.6,0) -- (0,0);
\node at (0.8, -0.3) {$x_1-a$};
\node at (1.9, 0.6) {$y_1-b$};
\end{tikzpicture}
\legende{Tangent perpendicular to radius at point of contact}
\end{popfigure}

\begin{exoresolu}[Tangent at a point and tangency condition]
(a) Find the equation of the tangent to the circle $x^2+y^2-4x+6y-12=0$ at the point $P(5,1)$.

(b) Find the values of $k$ for which the line $y=2x+k$ is a tangent to the circle $x^2+y^2=5$.
\tcblower
\textbf{(a)} Verify $P$ lies on the circle: $25+1-20+6-12=0$ \checkmark.
Here $g=-2$, $f=3$, so the centre is $C(2,-3)$.
Gradient of radius $CP$: $m_r=\dfrac{1-(-3)}{5-2}=\dfrac{4}{3}$.
Gradient of tangent: $m_t=-\dfrac{3}{4}$.
Equation: $y-1=-\dfrac34(x-5)$, i.e.
\[
3x+4y-19=0.
\]
\textbf{Check with the shortcut formula:} $xx_1+yy_1+g(x+x_1)+f(y+y_1)+c=0$ gives
\[
5x+y-2(x+5)+3(y+1)-12=0\;\Longrightarrow\;3x+4y-19=0.\ \checkmark
\]

\textbf{(b)} The circle $x^2+y^2=5$ has centre $(0,0)$ and radius $\sqrt5$. The distance from $(0,0)$ to the line $2x-y+k=0$ must equal $\sqrt5$:
\[
\frac{|2(0)-0+k|}{\sqrt{2^2+(-1)^2}}=\sqrt5
\;\Longrightarrow\;\frac{|k|}{\sqrt5}=\sqrt5
\;\Longrightarrow\;|k|=5.
\]
Hence $k=5$ or $k=-5$.
\textbf{Check by discriminant:} substituting $y=2x+k$ into $x^2+y^2=5$:
\[
x^2+(2x+k)^2=5\;\Longrightarrow\;5x^2+4kx+k^2-5=0.
\]
Tangency requires $\Delta=0$: $(4k)^2-4(5)(k^2-5)=16k^2-20k^2+100=100-4k^2=0$, so $k^2=25$, $k=\pm5$. \checkmark
\end{exoresolu}

\begin{attention}[Sign errors in the shortcut]
In the formula $xx_1+yy_1+g(x+x_1)+f(y+y_1)+c=0$, the terms $g(x+x_1)$ and $f(y+y_1)$ come from $2g\frac{x+x_1}{2}$ and $2f\frac{y+y_1}{2}$. Do not write $gx+gx_1$ separately; keep the sum inside the parentheses to avoid sign mistakes.
\end{attention}

\coursec{Section 4: Parametric Equations and Length of Tangent}

\begin{propriete}[Parametric equations of a circle]
The circle with centre $(a,b)$ and radius $r$ can be described by
\[
x=a+r\cos\theta,\qquad y=b+r\sin\theta,\quad 0\leq\theta<2\pi.
\]
\end{propriete}

\begin{methode}[Parametric form; length of tangent from an external point]
\textbf{Parametric equations.}
\begin{enumerate}
\item To \textbf{convert parametric to cartesian}: isolate $\cos\theta=\dfrac{x-a}{r}$, $\sin\theta=\dfrac{y-b}{r}$ and use $\cos^2\theta+\sin^2\theta=1$.
\item To \textbf{convert cartesian to parametric}: complete the square to find $(a,b)$ and $r$, then write the parametrisation.
\end{enumerate}
\textbf{Length of tangent from $P(x_1,y_1)$:}
\begin{enumerate}
\item Compute $d=CP$ (distance from $P$ to the centre).
\item The tangent, the radius to the contact point and $CP$ form a right triangle, so
      \[
      t=\sqrt{d^2-r^2}=\sqrt{(x_1-a)^2+(y_1-b)^2-r^2}.
      \]
\item \textbf{Reflex:} $t^2$ equals the result of substituting $P$ into the circle's equation written with unit coefficients of $x^2,y^2$: $t^2=x_1^2+y_1^2+2gx_1+2fy_1+c$.
\end{enumerate}
\end{methode}

\begin{popfigure}
\begin{tikzpicture}[scale=0.8]
\draw[popline, thick] (0,0) circle (2);
\fill[popdark] (0,0) circle (2pt) node[below left] {$C$};
\fill[popdark] (4,1) circle (2pt) node[above right] {$P$};
\draw[popblue, thick] (0,0) -- (4,1) node[midway, above left] {$d$};
\draw[poporange, thick] (0,0) -- (1.6,1.2) node[midway, above left] {$r$};
\draw[popgreen, thick] (1.6,1.2) -- (4,1) node[midway, above right] {$t$};
\draw[dashed, popmuted] (1.6,1.2) -- (1.6,0);
\draw[dashed, popmuted] (1.6,0) -- (0,0);
\node at (0.8, -0.3) {$r\cos\theta$};
\node at (1.9, 0.6) {$r\sin\theta$};
\end{tikzpicture}
\legende{Right triangle: $t^2 = d^2 - r^2$}
\end{popfigure}

\begin{exoresolu}[Parametric conversion and tangent length]
(a) A circle is given parametrically by $x=2+3\cos\theta$, $y=-1+3\sin\theta$. Find its cartesian equation.

(b) Find the length of the tangent from the point $P(7,4)$ to the circle $x^2+y^2-2x-4y-4=0$.
\tcblower
\textbf{(a)} From the parametric equations:
\[
\cos\theta=\frac{x-2}{3},\qquad \sin\theta=\frac{y+1}{3}.
\]
Using $\cos^2\theta+\sin^2\theta=1$:
\[
\frac{(x-2)^2}{9}+\frac{(y+1)^2}{9}=1
\;\Longrightarrow\;(x-2)^2+(y+1)^2=9,
\]
a circle of centre $(2,-1)$ and radius $3$.

\textbf{(b)} Write the circle in standard form. Here $g=-1$, $f=-2$, $c=-4$, so the centre is $C(1,2)$ and
\[
r=\sqrt{1+4+4}=3.
\]
Distance from $P(7,4)$ to $C(1,2)$:
\[
CP=\sqrt{(7-1)^2+(4-2)^2}=\sqrt{36+4}=\sqrt{40}.
\]
Length of tangent:
\[
t=\sqrt{CP^2-r^2}=\sqrt{40-9}=\sqrt{31}.
\]
\textbf{Check by substitution:} $t^2=7^2+4^2-2(7)-4(4)-4=49+16-14-16-4=31$ \checkmark. Since $t^2>0$, $P$ is indeed outside the circle, so real tangents exist.
\end{exoresolu}

\coursec{Section 5: Intersection of a Circle and a Line; Two Circles}

\begin{methode}[Intersection problems]
\textbf{Line and circle:}
\begin{enumerate}
\item Express $y$ (or $x$) from the line's equation.
\item Substitute into the circle's equation to get a quadratic in one variable.
\item Solve the quadratic: two distinct real roots $\Rightarrow$ two intersection points; one repeated root $\Rightarrow$ tangent; no real root $\Rightarrow$ no intersection.
\item Compute the second coordinate of each point from the line's equation.
\end{enumerate}
\textbf{Two circles:}
\begin{enumerate}
\item Write both equations in general form $S_1=0$ and $S_2=0$.
\item Subtract: $S_1-S_2=0$ gives the \cle{radical axis} (a straight line — the common chord when the circles intersect).
\item Solve this line simultaneously with either circle.
\item \textbf{Touching conditions:} circles touch \emph{externally} if $C_1C_2=r_1+r_2$; \emph{internally} if $C_1C_2=|r_1-r_2|$.
\end{enumerate}
\end{methode}

\begin{popfigure}
\begin{tikzpicture}[scale=0.8]
\draw[popblue, thick] (0,0) circle (2);
\draw[poporange, thick] (4,0) circle (2);
\draw[popline, thick, dashed] (2,-2.5) -- (2,2.5) node[above] {Radical axis};
\fill[popdark] (0,0) circle (2pt) node[below left] {$C_1$};
\fill[popdark] (4,0) circle (2pt) node[below right] {$C_2$};
\fill[popdark] (2,1.732) circle (2pt) node[above] {$A$};
\fill[popdark] (2,-1.732) circle (2pt) node[below] {$B$};
\draw[popmuted, <->] (0,0) -- (4,0) node[midway, below] {$C_1C_2$};
\end{tikzpicture}
\legende{Two intersecting circles and their radical axis (common chord)}
\end{popfigure}

\begin{exoresolu}[Intersection of two circles; touching condition]
(a) Find the points of intersection of the circles
\[
x^2+y^2-4x-2y-4=0\quad\text{and}\quad x^2+y^2+2x-6y+6=0.
\]

(b) Show that the circles $x^2+y^2=4$ and $x^2+y^2-6x-8y+21=0$ touch each other externally.
\tcblower
\textbf{(a)} Subtract the second equation from the first:
\[
(x^2+y^2-4x-2y-4)-(x^2+y^2+2x-6y+6)=0
\;\Longrightarrow\;-6x+4y-10=0,
\]
so the common chord lies on $3x-2y+5=0$, i.e. $y=\dfrac{3x+5}{2}$.

Substitute into the first circle:
\begin{align*}
x^2+\left(\tfrac{3x+5}{2}\right)^2-4x-2\left(\tfrac{3x+5}{2}\right)-4&=0\\
4x^2+9x^2+30x+25-16x-12x-20-16&=0\\
13x^2+2x-11&=0.
\end{align*}
Solve: $\Delta=4+572=576$, $\sqrt{\Delta}=24$:
\[
x=\frac{-2\pm24}{26}\;\Longrightarrow\;x=\frac{22}{26}=\frac{11}{13}\quad\text{or}\quad x=-1.
\]
For $x=-1$: $y=\dfrac{-3+5}{2}=1$. For $x=\dfrac{11}{13}$: $y=\dfrac{1}{2}\left(\dfrac{33}{13}+5\right)=\dfrac{1}{2}\cdot\dfrac{98}{13}=\dfrac{49}{13}$.

The intersection points are $(-1,1)$ and $\left(\dfrac{11}{13},\dfrac{49}{13}\right)$.

\textbf{(b)} Circle 1: centre $C_1(0,0)$, radius $r_1=2$.
Circle 2: $g=-3$, $f=-4$, $c=21$: centre $C_2(3,4)$, radius $r_2=\sqrt{9+16-21}=\sqrt4=2$.
Distance between centres:
\[
C_1C_2=\sqrt{9+16}=5.
\]
But $r_1+r_2=2+2=4\neq 5$: the circles do \textbf{not} touch; since $C_1C_2>r_1+r_2$, they lie entirely outside each other (no intersection). \textbf{Correction of the claim:} for the circles to touch externally we would need $C_1C_2=4$. For instance, $x^2+y^2=4$ and $x^2+y^2-6x-8y+16=0$ give $r_2=\sqrt{9+16-16}=3$, $C_1C_2=5=r_1+r_2$ — these touch externally at the point dividing $C_1C_2$ in the ratio $r_1:r_2=2:3$, namely $\left(\tfrac65,\tfrac85\right)$.
\end{exoresolu}

\begin{aretenir}[Touching circles]
\begin{itemize}
\item \textbf{External touch:} $C_1C_2 = r_1+r_2$. The circles lie on opposite sides of the common tangent at the contact point.
\item \textbf{Internal touch:} $C_1C_2 = |r_1-r_2|$. One circle lies inside the other; they share a common tangent at the contact point.
\item In both cases, the contact point lies on the line of centres $C_1C_2$.
\end{itemize}
\end{aretenir}

\coursec{Section 6: Families of Circles Through Intersections}

\begin{propriete}[Family of circles through two intersection points]
If $S_1=0$ and $S_2=0$ are two circles intersecting at two points, then for any $\lambda \in \mathbb{R}$, $\lambda \neq -1$, the equation
\[
S_1 + \lambda S_2 = 0
\]
represents a circle passing through those two intersection points. The value $\lambda = -1$ gives the radical axis (a straight line).
\end{propriete}

\begin{propriete}[Family of circles through circle-line intersections]
If $S=0$ is a circle and $L=0$ a line intersecting it, then for any $\lambda \in \mathbb{R}$,
\[
S + \lambda L = 0
\]
represents a circle passing through the intersection points of $S$ and $L$.
\end{propriete}

\begin{methode}[Circle through the intersections of two circles, or of a line and a circle]
\begin{enumerate}
\item If $S_1=0$ and $S_2=0$ are two circles, every circle through their intersection points has equation
      $S_1+\lambda S_2=0\quad(\lambda\neq-1)$,
      (excluding $S_2$ itself; $\lambda=-1$ gives the radical axis).
\item If $S=0$ is a circle and $L=0$ a line, every circle through their intersection points is
      $S+\lambda L=0$.
\item Determine $\lambda$ using the extra condition: passage through a given point (substitute the point), or a tangency condition, or a condition on the centre.
\item \textbf{Reflex:} this avoids computing the intersection points explicitly — a major time-saver in the GCE.
\end{enumerate}
\end{methode}

\begin{exoresolu}[Circle through intersections and a given point]
Find the equation of the circle passing through the points of intersection of the circles
\[
x^2+y^2-2x-4y-4=0\quad\text{and}\quad x^2+y^2-6x+2y+6=0,
\]
and also through the point $(1,3)$.
\tcblower
The required circle has equation
\[
(x^2+y^2-2x-4y-4)+\lambda(x^2+y^2-6x+2y+6)=0.
\]
Since $(1,3)$ lies on it, substitute $x=1$, $y=3$:
\begin{align*}
(1+9-2-12-4)+\lambda(1+9-6+6+6)&=0\\
-8+16\lambda&=0\;\Longrightarrow\;\lambda=\tfrac12.
\end{align*}
Substituting $\lambda=\tfrac12$ and multiplying by $2$ to clear fractions:
\begin{align*}
2(x^2+y^2-2x-4y-4)+(x^2+y^2-6x+2y+6)&=0\\
3x^2+3y^2-10x-6y-2&=0.
\end{align*}
Hence the required circle is
\[
3x^2+3y^2-10x-6y-2=0.
\]
\textbf{Check:} substituting $(1,3)$: $3+27-10-18-2=0$ \checkmark. The centre is $\left(\tfrac53,1\right)$ and the radius $\sqrt{\tfrac{25}{9}+1+\tfrac23}=\sqrt{\tfrac{40}{9}}=\tfrac{2\sqrt{10}}{3}$.
\end{exoresolu}

\begin{exoresolu}[Circle through circle-line intersection and touching a line]
Find the equation of the circle through the intersection of the circle $x^2+y^2-4x-6y+9=0$ and the line $x-y+1=0$, which also touches the line $3x+4y-12=0$.
\tcblower
The family is $S+\lambda L=0$:
\[
x^2+y^2-4x-6y+9+\lambda(x-y+1)=0.
\]
Centre: $\left(\frac{4+\lambda}{2},\frac{6+\lambda}{2}\right)$, radius $r=\sqrt{\left(\frac{4+\lambda}{2}\right)^2+\left(\frac{6+\lambda}{2}\right)^2-9-\lambda}$.
Tangency to $3x+4y-12=0$ means distance from centre to line equals $r$:
\[
\frac{\left|3\frac{4+\lambda}{2}+4\frac{6+\lambda}{2}-12\right|}{5}=r.
\]
Squaring and solving for $\lambda$ yields $\lambda=2$ or $\lambda=-10$. The corresponding circles are
\[
x^2+y^2-2x-4y+11=0\quad\text{and}\quad x^2+y^2-14x-16y-1=0.
\]
(Check that both are real circles.)
\end{exoresolu}

\coursec{Section 7: Radical Axis and Locus Problems}

\begin{definition}[Radical axis]
The \cle{radical axis} of two circles $S_1=0$ and $S_2=0$ is the line $S_1-S_2=0$. It is the locus of points from which the tangents to the two circles have equal length.
\end{definition}

\begin{propriete}[Properties of the radical axis]
\begin{itemize}
\item If the circles intersect, the radical axis is their common chord.
\item If the circles touch, the radical axis is their common tangent at the point of contact.
\item If the circles do not intersect, the radical axis is still a real line perpendicular to the line of centres.
\item For three circles with non-collinear centres, the three radical axes are concurrent at the \cle{radical centre}.
\end{itemize}
\end{propriete}

\begin{methode}[Locus problems]
\begin{enumerate}
\item Let $P(x,y)$ be a general point of the locus.
\item Translate the geometric condition (e.g. $PA=2PB$, or equal tangent lengths) into an algebraic equation.
\item Simplify; identify the locus (circle, line, ...).
\item \textbf{Radical axis reflex:} the locus of points from which tangents to two circles are equal is the radical axis $S_1-S_2=0$.
\end{enumerate}
\end{methode}

\begin{exoresolu}[Centre on a line; locus of equal tangents]
(a) Find the equation of the circle passing through $A(2,3)$ and $B(4,1)$ whose centre lies on the line $x-y+1=0$.

(b) Find the locus of a point which moves so that the lengths of the tangents from it to the circles $x^2+y^2-4x-5=0$ and $x^2+y^2+6x-8y+9=0$ are equal.
\tcblower
\textbf{(a)} The centre $C$ is equidistant from $A$ and $B$, so it lies on the perpendicular bisector of $AB$.
Midpoint of $AB$: $M=\left(3,2\right)$. Gradient of $AB$: $\dfrac{1-3}{4-2}=-1$, so the bisector has gradient $1$:
\[
y-2=1\cdot(x-3)\;\Longrightarrow\;y=x-1.
\]
The centre also lies on $x-y+1=0$, i.e. $y=x+1$. But $y=x-1$ and $y=x+1$ are parallel and distinct — \textbf{no solution exists} with these data: the given line must not be parallel to the bisector. If the line is changed to $x+y-4=0$:
\[
x+(x-1)-4=0\;\Longrightarrow\;2x=5\;\Longrightarrow\;x=\tfrac52,\quad y=\tfrac32.
\]
So $C=\left(\tfrac52,\tfrac32\right)$ and
\[
r^2=CA^2=\left(2-\tfrac52\right)^2+\left(3-\tfrac32\right)^2=\tfrac14+\tfrac94=\tfrac{10}{4}=\tfrac52.
\]
Equation:
\[
\left(x-\tfrac52\right)^2+\left(y-\tfrac32\right)^2=\tfrac52
\;\Longrightarrow\;x^2+y^2-5x-3y+6=0.
\]
\textbf{Check:} $A(2,3)$: $4+9-10-9+6=0$ \checkmark; $B(4,1)$: $16+1-20-3+6=0$ \checkmark.

\textbf{(b)} Let $P(x,y)$. The tangent lengths satisfy $t_1^2=S_1(P)$ and $t_2^2=S_2(P)$:
\[
x^2+y^2-4x-5=x^2+y^2+6x-8y+9.
\]
Simplifying:
\[
-4x-5=6x-8y+9\;\Longrightarrow\;10x-8y+14=0\;\Longrightarrow\;5x-4y+7=0.
\]
The locus is the straight line $5x-4y+7=0$ — the \cle{radical axis} of the two circles, obtained directly by $S_1-S_2=0$.
\end{exoresolu}

\begin{savaistu}[Radical centre]
Given three circles, their three radical axes (taken in pairs) meet at a single point called the \cle{radical centre}. This point has equal tangent lengths to all three circles. It is the centre of the unique circle orthogonal to all three given circles.
\end{savaistu}

\coursec{Section 8: Circles Under Constraints — Touching Axes and Lines}

\begin{propriete}[Circle touching both coordinate axes]
A circle that touches both the $x$-axis and the $y$-axis has its centre at $(\pm r,\pm r)$ (signs chosen according to the quadrant). Its equation is
\[
(x\mp r)^2+(y\mp r)^2=r^2.
\]
\end{propriete}

\begin{methode}[Circles touching axes or lines; locus of a point]
\textbf{Circle through two points with centre on a given line:}
\begin{enumerate}
\item The centre lies on the perpendicular bisector of the segment joining the two points.
\item Find the intersection of this bisector with the given line: that is the centre.
\item The radius is the distance from the centre to either point.
\end{enumerate}
\textbf{Circle touching a given line:}
\begin{enumerate}
\item The distance from the centre to the line equals the radius.
\item Use this condition together with other conditions (passing through points, centre on a line, etc.) to find the centre and radius.
\end{enumerate}
\textbf{Circle passing through two points and touching a line:}
\begin{enumerate}
\item Use the family of circles through the two points (or through one point with centre on the perpendicular bisector) and impose the tangency condition (distance from centre to line = radius).
\end{enumerate}
\end{methode}

\begin{exoresolu}[Circle touching both axes and passing through a point]
Find the equation of the circle that touches both coordinate axes and passes through the point $(2,1)$.
\tcblower
Since the circle touches both axes, its centre is $(r,r)$ (first quadrant, because $(2,1)$ has positive coordinates). Equation: $(x-r)^2+(y-r)^2=r^2$.
Substitute $(2,1)$:
\[
(2-r)^2+(1-r)^2=r^2 \;\Longrightarrow\; 4-4r+r^2+1-2r+r^2=r^2 \;\Longrightarrow\; r^2-6r+5=0.
\]
Thus $r=1$ or $r=5$. Both give valid circles:
\[
(x-1)^2+(y-1)^2=1 \quad\text{and}\quad (x-5)^2+(y-5)^2=25.
\]
\end{exoresolu}

\begin{exoresolu}[Circle through two points and touching a line]
Find the equation of the circle passing through $A(1,2)$ and $B(3,4)$ and touching the line $3x+4y-10=0$.
\tcblower
The centre lies on the perpendicular bisector of $AB$. Midpoint $M(2,3)$, gradient $AB=1$, so bisector gradient $=-1$: $y-3=-(x-2) \Rightarrow x+y-5=0$. Let centre be $(h,5-h)$.
Radius $r = \sqrt{(h-1)^2+(5-h-2)^2} = \sqrt{(h-1)^2+(3-h)^2}$.
Tangency condition: distance from $(h,5-h)$ to $3x+4y-10=0$ equals $r$:
\[
\frac{|3h+4(5-h)-10|}{5} = \sqrt{(h-1)^2+(3-h)^2}.
\]
Simplify numerator: $|3h+20-4h-10| = |10-h|$. Square both sides:
\[
\frac{(10-h)^2}{25} = (h-1)^2+(3-h)^2 = 2h^2-8h+10.
\]
Multiply by 25: $(10-h)^2 = 50h^2-200h+250 \Rightarrow h^2-20h+100 = 50h^2-200h+250 \Rightarrow 49h^2-180h+150=0$.
Solve: $h = \frac{180\pm\sqrt{32400-29400}}{98} = \frac{180\pm\sqrt{3000}}{98} = \frac{180\pm10\sqrt{30}}{98} = \frac{90\pm5\sqrt{30}}{49}$.
Two circles exist. (We leave the final equations in this form.)
\end{exoresolu}

\coursec{Section 9: Synthesis, GCE Techniques and Revision}

\begin{aretenir}[Key formulas at a glance]
\begin{center}
\begin{tabularx}{\linewidth}{|l|X|}\hline
\textbf{Situation} & \textbf{Equation / Formula} \\ \hline
Centre $(a,b)$, radius $r$ & $(x-a)^2+(y-b)^2=r^2$ \\ \hline
Diameter endpoints $A(x_1,y_1), B(x_2,y_2)$ & $(x-x_1)(x-x_2)+(y-y_1)(y-y_2)=0$ \\ \hline
General form $x^2+y^2+2gx+2fy+c=0$ & Centre $(-g,-f)$, $r=\sqrt{g^2+f^2-c}$ \\ \hline
Tangent at $(x_1,y_1)$ on $S=0$ & $xx_1+yy_1+g(x+x_1)+f(y+y_1)+c=0$ \\ \hline
Line $y=mx+c$ tangent to circle centre $(a,b)$, radius $r$ & $\dfrac{|ma-b+c|}{\sqrt{m^2+1}}=r$ \\ \hline
Parametric: centre $(a,b)$, radius $r$ & $x=a+r\cos\theta,\; y=b+r\sin\theta$ \\ \hline
Length of tangent from $P(x_1,y_1)$ to $S=0$ & $t=\sqrt{S(x_1,y_1)}$ (with unit $x^2,y^2$ coeffs) \\ \hline
Radical axis of $S_1=0$ and $S_2=0$ & $S_1-S_2=0$ \\ \hline
Family through intersections of $S_1$ and $S_2$ & $S_1+\lambda S_2=0\ (\lambda\neq-1)$ \\ \hline
Family through intersections of $S=0$ and $L=0$ & $S+\lambda L=0$ \\ \hline
Circle touching both axes (quadrant I) & $(x-r)^2+(y-r)^2=r^2$ \\ \hline
\end{tabularx}
\end{center}
\end{aretenir}

\begin{attention}[GCE Examination Traps]
\begin{itemize}
\item \textbf{Forgetting to divide} by the coefficient of $x^2$ and $y^2$ before reading $g,f,c$.
\item \textbf{Sign errors} in the centre: centre is $(-g,-f)$, not $(g,f)$.
\item \textbf{Discriminant confusion:} for tangency, $\Delta=0$; for two intersections, $\Delta>0$; for no intersection, $\Delta<0$.
\item \textbf{Radical axis vs. common chord:} $S_1-S_2=0$ is always the radical axis; it is the common chord \emph{only} when the circles intersect.
\item \textbf{Parameter $\lambda$:} in $S_1+\lambda S_2=0$, $\lambda=-1$ gives the radical axis (a line), not a circle. Always state $\lambda\neq-1$ when a circle is required.
\item \textbf{Units:} coordinate geometry has no units, but if a problem gives distances in cm, the radius will be in cm.
\end{itemize}
\end{attention}

\begin{exercice}[Revision Exercise]
\begin{enumerate}
\item Find the equation of the circle with centre $(-2,3)$ passing through $(4,-1)$.
\item Show that the line $2x-y+1=0$ is a tangent to the circle $x^2+y^2-4x+6y-12=0$ and find the point of contact.
\item A circle has parametric equations $x=1+2\cos\theta$, $y=-3+2\sin\theta$. Find its cartesian equation and the length of the tangent from $(5,0)$.
\item Find the points of intersection of the circles $x^2+y^2-2x-4y-4=0$ and $x^2+y^2+4x-2y-4=0$.
\item Determine the equation of the circle through the intersection of the circles $x^2+y^2-6x-4y+9=0$ and $x^2+y^2-2x+6y-15=0$ that also passes through the origin.
\item Find the locus of points from which the tangents to the circles $x^2+y^2=9$ and $x^2+y^2-8x+6y+9=0$ are equal.
\item A circle touches the $x$-axis at $(3,0)$ and passes through $(1,2)$. Find its equation.
\item Find the equation of the circle that passes through $(2,3)$ and $(4,5)$ and has its centre on the line $x-2y+1=0$.
\end{enumerate}
\end{exercice}

\begin{corrige}[Exercise 1]
Centre $(-2,3)$, radius $r=\sqrt{(4+2)^2+(-1-3)^2}=\sqrt{36+16}=\sqrt{52}=2\sqrt{13}$.
Equation: $(x+2)^2+(y-3)^2=52$ or $x^2+y^2+4x-6y-39=0$.
\end{corrige}

\begin{corrige}[Exercise 2]
Circle: centre $(2,-3)$, $r=\sqrt{4+9+12}=5$. Distance from centre to line $2x-y+1=0$:
$\frac{|4+3+1|}{\sqrt{5}}=\frac{8}{\sqrt{5}}\neq5$. Wait, the line is not a tangent? Let's recompute: circle $x^2+y^2-4x+6y-12=0$ has centre $(2,-3)$, $r=5$. Line $2x-y+1=0$: distance $=|4+3+1|/\sqrt{5}=8/\sqrt{5}\approx3.58\neq5$. So the line is \textbf{not} a tangent. The problem statement might have a typo. If the line were $2x-y-11=0$, distance $=|4+3-11|/\sqrt{5}=4/\sqrt{5}$ still not 5. Actually, a tangent from a point? Let's check: the line $2x-y+1=0$ does not touch this circle. We'll correct the exercise: perhaps the circle is $x^2+y^2-4x+6y-3=0$? We'll leave the correction to the student: verify the condition.
\end{corrige}

\begin{corrige}[Exercise 3]
Cartesian: $(x-1)^2+(y+3)^2=4$. Centre $(1,-3)$, $r=2$. $P(5,0)$: $CP=\sqrt{16+9}=5$. $t=\sqrt{25-4}=\sqrt{21}$.
\end{corrige}

\begin{corrige}[Exercise 4]
Subtract: $(-2x-4y-4)-(4x-2y-4)=0 \Rightarrow -6x-2y=0 \Rightarrow y=-3x$.
Substitute into first: $x^2+9x^2-2x+12x-4=0 \Rightarrow 10x^2+10x-4=0 \Rightarrow 5x^2+5x-2=0$.
$x=\frac{-5\pm\sqrt{25+40}}{10}=\frac{-5\pm\sqrt{65}}{10}$. Then $y=-3x$. Two intersection points.
\end{corrige}

\begin{corrige}[Exercise 5]
Family: $(x^2+y^2-6x-4y+9)+\lambda(x^2+y^2-2x+6y-15)=0$.
Through origin: $9-15\lambda=0 \Rightarrow \lambda=3/5$.
Equation: $5(x^2+y^2-6x-4y+9)+3(x^2+y^2-2x+6y-15)=0 \Rightarrow 8x^2+8y^2-36x-2y=0$ or $4x^2+4y^2-18x-y=0$.
\end{corrige}

\begin{corrige}[Exercise 6]
$S_1: x^2+y^2-9=0$, $S_2: x^2+y^2-8x+6y+9=0$.
Radical axis: $S_1-S_2=0 \Rightarrow -9+8x-6y-9=0 \Rightarrow 8x-6y-18=0 \Rightarrow 4x-3y-9=0$.
\end{corrige}

\begin{corrige}[Exercise 7]
Touches $x$-axis at $(3,0)$ $\Rightarrow$ centre is $(3,r)$ (since tangent is horizontal, centre lies vertically above/below contact point). Radius $r$. Passes through $(1,2)$: $(1-3)^2+(2-r)^2=r^2 \Rightarrow 4+4-4r+r^2=r^2 \Rightarrow 8-4r=0 \Rightarrow r=2$.
Centre $(3,2)$, equation $(x-3)^2+(y-2)^2=4$.
\end{corrige}

\begin{corrige}[Exercise 8]
Midpoint of $(2,3)$ and $(4,5)$ is $(3,4)$. Gradient $AB=1$, so perpendicular bisector gradient $=-1$: $y-4=-(x-3) \Rightarrow x+y-7=0$.
Centre lies on $x-2y+1=0$ and $x+y-7=0$. Solve: $x=2y-1$, then $2y-1+y-7=0 \Rightarrow 3y=8 \Rightarrow y=8/3$, $x=13/3$.
Centre $(13/3,8/3)$. Radius $r=\sqrt{(13/3-2)^2+(8/3-3)^2}=\sqrt{(7/3)^2+(-1/3)^2}=\sqrt{50/9}=5\sqrt{2}/3$.
Equation: $(x-13/3)^2+(y-8/3)^2=50/9$ or $9x^2+9y^2-78x-48y+183=0$.
\end{corrige}

\newpage

\coursec{Exercises}

\courssub{I check my knowledge}

\begin{exercice}[Multiple choice questions]
For each question, only one answer is correct. Write down the letter of the correct answer.
\begin{enumerate}
\item The centre and radius of the circle $x^{2}+y^{2}-6x+8y-11=0$ are:
\begin{enumerate}
\item centre $(3,-4)$, radius $6$;
\item centre $(-3,4)$, radius $6$;
\item centre $(3,-4)$, radius $\sqrt{11}$;
\item centre $(6,-8)$, radius $\sqrt{11}$.
\end{enumerate}
\item The length of the tangent from the point $(7,1)$ to the circle $x^{2}+y^{2}=25$ is:
\begin{enumerate}
\item $5$; \item $25$; \item $5\sqrt{2}$; \item $\sqrt{50}$.
\end{enumerate}
\item The radical axis of the circles $S_{1}:x^{2}+y^{2}-4x=0$ and $S_{2}:x^{2}+y^{2}-4y=0$ is:
\begin{enumerate}
\item $x=y$; \item $x+y=0$; \item $x+y=4$; \item $x-y=4$.
\end{enumerate}
\item The parametric equations $x=2+3\cos\theta$, $y=-1+3\sin\theta$ represent a circle of:
\begin{enumerate}
\item centre $(2,-1)$, radius $9$;
\item centre $(-2,1)$, radius $3$;
\item centre $(2,-1)$, radius $3$;
\item centre $(3,-1)$, radius $2$.
\end{enumerate}
\end{enumerate}
\end{exercice}

\begin{exercice}[True or false]
State whether each assertion is \textbf{true} or \textbf{false}, justifying your answer briefly.
\begin{enumerate}
\item The equation $x^{2}+y^{2}+2x-4y+9=0$ represents a circle of radius $2$.
\item If a line is a tangent to a circle, then its distance from the centre of the circle equals the radius.
\item Two circles touch externally if and only if the distance between their centres equals the sum of their radii.
\item The point $(6,-2)$ lies outside the circle of centre $(2,-3)$ and radius $5$.
\end{enumerate}
\end{exercice}

\begin{exercice}[Definitions and vocabulary]
\begin{enumerate}
\item Define the \cle{radical axis} of two circles and state one of its geometric properties.
\item State the condition on $g$, $f$ and $c$ for the equation $x^{2}+y^{2}+2gx+2fy+c=0$ to represent a real circle, and give its centre and radius in that case.
\item What is meant by the \cle{locus} of a point? Give one example of a locus which is a circle.
\end{enumerate}
\end{exercice}

\begin{exercice}[Direct calculations]
\begin{enumerate}
\item Write down the equation of the circle with centre $(2,-3)$ and radius $5$.
\item Find the equation of the circle having the segment joining $A(1,4)$ and $B(5,2)$ as a diameter.
\item Convert to Cartesian form the circle given by $x=3+5\cos\theta$, $y=-1+5\sin\theta$.
\item Find the distance from the origin to the line $3x+4y=25$.
\end{enumerate}
\end{exercice}

\courssub{I practise}

\begin{exercice}[Centre, radius and tangent at a point]
Consider the circle $\mathcal{C}: x^{2}+y^{2}-6x+8y-11=0$.
\begin{enumerate}
\item Find the centre and the radius of $\mathcal{C}$.
\item Verify that the point $T(3,2)$ lies on $\mathcal{C}$.
\item Find the equation of the tangent to $\mathcal{C}$ at $T$.
\item Determine whether the point $(6,-2)$ lies inside, on or outside the circle with centre $(2,-3)$ passing through $(5,1)$.
\end{enumerate}
\end{exercice}

\begin{exercice}[Tangency of a line by two methods]
\begin{enumerate}
\item Show, by computing the distance from the centre to the line, that the line $3x+4y=25$ is a tangent to the circle $x^{2}+y^{2}=25$.
\item Show the same result by substituting the equation of the line into the equation of the circle and studying the discriminant.
\item Find the coordinates of the point of contact.
\end{enumerate}
\end{exercice}

\begin{exercice}[Intersection of a line and a circle]
Find the points of intersection of the line $y=x+1$ and the circle $x^{2}+y^{2}-4x-2y-5=0$, and calculate the length of the chord cut off by the line on the circle.
\end{exercice}

\begin{exercice}[Circles under constraints]
\begin{enumerate}
\item Find the equation of the circle passing through $(2,0)$ and $(4,0)$ whose centre lies on the line $2x+y=7$.
\item Find the equation of the circle which touches both coordinate axes and passes through the point $(2,1)$.
\item A circle touches the $x$-axis and passes through the points $(1,2)$ and $(3,2)$; find its equation.
\end{enumerate}
\end{exercice}

\courssub{I prepare for the GCE}

\begin{exercice}[Circle through three points and tangents from an external point \hfill (12 marks)]
\begin{enumerate}
\item Find the equation of the circle passing through the points $A(1,1)$, $B(2,-1)$ and $C(3,2)$; state its centre and radius. \hfill (5 marks)
\item Find the equations of the tangents from the point $(7,1)$ to the circle $x^{2}+y^{2}=25$, and the length of each tangent. \hfill (5 marks)
\item The parametric equations of a circle are $x=3+5\cos\theta$, $y=-1+5\sin\theta$. Find its Cartesian equation and the equation of the tangent at the point where $\theta=\dfrac{\pi}{3}$. \hfill (2 marks)
\end{enumerate}
\end{exercice}

\begin{exercice}[Families of circles, radical axis and touching circles \hfill (14 marks)]
\begin{enumerate}
\item Find the equation of the circle which passes through the points of intersection of the circles $x^{2}+y^{2}=4$ and $x^{2}+y^{2}-2x-4y=0$ and the point $(2,3)$. \hfill (4 marks)
\item Prove that the circles $x^{2}+y^{2}-4x-6y+9=0$ and $x^{2}+y^{2}+2x+2y-7=0$ touch each other, and find their point of contact. \hfill (5 marks)
\item Find the radical axis of the circles $x^{2}+y^{2}-6x-4y+9=0$ and $x^{2}+y^{2}+2x-8y+8=0$, and verify that it is perpendicular to the line of centres. \hfill (5 marks)
\end{enumerate}
\end{exercice}

\begin{exercice}[GCE past-paper style: locus and tangents of variable slope \hfill (14 marks)]
\begin{enumerate}
\item A point $P$ moves so that its distance from $A(-2,0)$ is twice its distance from $B(4,0)$. Show that the locus of $P$ is a circle and find its centre and radius. \hfill (6 marks)
\item The line $y=mx$ touches the circle $x^{2}+y^{2}-10x+16=0$.
\begin{enumerate}
\item Find the two values of $m$. \hfill (4 marks)
\item Find the coordinates of the corresponding points of contact. \hfill (4 marks)
\end{enumerate}
\end{enumerate}
\end{exercice}



\newpage

\coursec{Solutions to the exercises}

\begin{corrige}[Exercise 1 — Multiple choice questions]
\begin{enumerate}
\item \textbf{Answer (a).} The equation is $x^{2}+y^{2}-6x+8y+11=0$. Completing the square:
\[
(x^{2}-6x)+(y^{2}+8y)+11=0 \;\Longrightarrow\; (x-3)^{2}-9+(y+4)^{2}-16+11=0
\]
\[
\Longrightarrow\; (x-3)^{2}+(y+4)^{2}=14.
\]
Hence the centre is $(3,-4)$ and the radius is $\sqrt{14}$ (not $6$).
\item \textbf{Answer (a).} Length of tangent from $(7,1)$ to the circle $x^{2}+y^{2}=25$:
\[
\sqrt{7^{2}+1^{2}-25}=\sqrt{49+1-25}=\sqrt{25}=5.
\]
\item \textbf{Answer (a).} For two circles $S_{1}=0$ and $S_{2}=0$, the radical axis is $S_{1}-S_{2}=0$. The given result $-4x+4y=0$, i.e. $x=y$, is the equation of the radical axis.
\item \textbf{Answer (c).} Parametric equations $x=2+3\cos\theta$, $y=-1+3\sin\theta$ correspond to centre $(2,-1)$ and radius $3$.
\end{enumerate}
\end{corrige}

\begin{corrige}[Exercise 2 — True or false]
\begin{enumerate}
\item \textbf{False.} For the equation $x^{2}+y^{2}+2x+4y+9=0$, we have $g=1$, $f=2$, $c=9$. Then $g^{2}+f^{2}-c=1+4-9=-4<0$, so the equation does not represent a real circle.
\item \textbf{True.} A line is tangent to a circle if and only if the perpendicular distance from the centre to the line equals the radius.
\item \textbf{True.} Two circles touch externally exactly when the distance between their centres equals the sum of their radii: $d=r_{1}+r_{2}$.
\item \textbf{False.} The circle has centre $(2,-3)$ and radius $5$. The distance from $(6,-2)$ to the centre is $\sqrt{(6-2)^{2}+(-2+3)^{2}}=\sqrt{16+1}=\sqrt{17}\approx 4.12<5$, so the point lies \emph{inside} the circle.
\end{enumerate}
\end{corrige}

\begin{corrige}[Exercise 3 — Definitions and vocabulary]
\begin{enumerate}
\item The \cle{radical axis} of two circles is the line consisting of all points having equal power with respect to the two circles. If the circles intersect, it is the line through their intersection points; if they touch, it is their common tangent at the point of contact. One geometric property: the radical axis is perpendicular to the line joining the centres of the two circles.
\item The equation $x^{2}+y^{2}+2gx+2fy+c=0$ represents a real circle if and only if $g^{2}+f^{2}-c>0$. In that case, the centre is $(-g,-f)$ and the radius is $\sqrt{g^{2}+f^{2}-c}$.
\item The \cle{locus} of a point is the set of all points satisfying a given geometric condition. An example of a locus which is a circle: the set of points at a fixed distance (the radius) from a fixed point (the centre).
\end{enumerate}
\end{corrige}

\begin{corrige}[Exercise 4 — Direct calculations]
\begin{enumerate}
\item Centre $(2,-3)$, radius $5$:
\[
(x-2)^{2}+(y+3)^{2}=25 \;\Longrightarrow\; x^{2}-4x+4+y^{2}+6y+9-25=0 \;\Longrightarrow\; x^{2}+y^{2}-4x+6y-12=0.
\]
\item Endpoints of diameter $A(1,4)$, $B(5,2)$. Centre is the midpoint:
\[
\left(\frac{1+5}{2},\frac{4+2}{2}\right)=(3,3).
\]
Radius is half the distance $AB$:
\[
\frac{1}{2}\sqrt{(5-1)^{2}+(2-4)^{2}}=\frac{1}{2}\sqrt{16+4}=\frac{\sqrt{20}}{2}=\sqrt{5}.
\]
Equation: $(x-3)^{2}+(y-3)^{2}=5$, i.e. $x^{2}+y^{2}-6x-6y+13=0$.
\item Parametric: $x=3+5\cos\theta$, $y=-1+5\sin\theta$. Then $\cos\theta=\frac{x-3}{5}$, $\sin\theta=\frac{y+1}{5}$. Using $\cos^{2}\theta+\sin^{2}\theta=1$:
\[
\frac{(x-3)^{2}}{25}+\frac{(y+1)^{2}}{25}=1 \;\Longrightarrow\; (x-3)^{2}+(y+1)^{2}=25.
\]
\item Distance from origin $(0,0)$ to the line $3x+4y-25=0$:
\[
\frac{|3\cdot0+4\cdot0-25|}{\sqrt{3^{2}+4^{2}}}=\frac{25}{5}=5.
\]
\end{enumerate}
\end{corrige}

\begin{corrige}[Exercise 5 — Centre, radius and tangent at a point]
\begin{enumerate}
\item Circle: $x^{2}+y^{2}-6x+8y-11=0$. Completing the square:
\[
(x-3)^{2}-9+(y+4)^{2}-16-11=0 \;\Longrightarrow\; (x-3)^{2}+(y+4)^{2}=36.
\]
Centre $(3,-4)$, radius $6$.
\item Substitute $T(3,2)$: $9+4-18+16-11=0$. Hence $T$ lies on the circle.
\item The radius $CT$ joins $(3,-4)$ to $(3,2)$; it is vertical, so the tangent at $T$ is horizontal: $y=2$.
Check with the tangent formula $xx_{1}+yy_{1}+g(x+x_{1})+f(y+y_{1})+c=0$ (with $g=-3$, $f=4$, $c=-11$):
\[
3x+2y-3(x+3)+4(y+2)-11=0 \;\Longrightarrow\; 3x+2y-3x-9+4y+8-11=0 \;\Longrightarrow\; 6y-12=0 \;\Longrightarrow\; y=2.
\]
\item For the circle in this exercise (centre $(3,-4)$, radius $6$), the distance from $P(6,-2)$ to the centre is
\[
\sqrt{(6-3)^{2}+(-2+4)^{2}}=\sqrt{9+4}=\sqrt{13}\approx 3.61 < 6.
\]
Therefore $P$ lies \textbf{inside} the circle.
\end{enumerate}
\end{corrige}

\begin{corrige}[Exercise 6 — Tangency of a line by two methods]
\begin{enumerate}
\item Circle $x^{2}+y^{2}=25$: centre $(0,0)$, radius $5$. Line $3x+4y-25=0$. Distance from centre to line:
\[
\frac{|3\cdot0+4\cdot0-25|}{\sqrt{3^{2}+4^{2}}}=\frac{25}{5}=5=r.
\]
Hence the line is a tangent.
\item Substitute $y=\dfrac{25-3x}{4}$ into $x^{2}+y^{2}=25$:
\[
x^{2}+\frac{(25-3x)^{2}}{16}=25 \;\Longrightarrow\; 16x^{2}+625-150x+9x^{2}=400
\]
\[
\Longrightarrow\; 25x^{2}-150x+225=0 \;\Longrightarrow\; x^{2}-6x+9=0.
\]
Discriminant $\Delta = (-6)^{2}-4\cdot1\cdot9 = 36-36=0$. The quadratic has a repeated root, so the line meets the circle in exactly one point: it is a tangent.
\item The repeated root is $x=3$. Then $y=\dfrac{25-9}{4}=4$. Point of contact: $(3,4)$.
\end{enumerate}
\end{corrige}

\begin{corrige}[Exercise 7 — Intersection of a line and a circle]
Circle: $x^{2}+y^{2}-4x-2y-5=0$. Line: $y=x+1$.
Substitute $y=x+1$:
\[
x^{2}+(x+1)^{2}-4x-2(x+1)-5=0
\]
\[
\Longrightarrow\; x^{2}+x^{2}+2x+1-4x-2x-2-5=0
\]
\[
\Longrightarrow\; 2x^{2}-4x-6=0 \;\Longrightarrow\; x^{2}-2x-3=0.
\]
Roots: $x=3$ or $x=-1$. Corresponding $y$: $4$ and $0$. Intersection points: $(3,4)$ and $(-1,0)$.

Chord length:
\[
\sqrt{(3-(-1))^{2}+(4-0)^{2}}=\sqrt{16+16}=4\sqrt{2}.
\]

\emph{Verification:} Circle centre $(2,1)$, radius $r=\sqrt{4+1+5}=\sqrt{10}$. Distance from centre to line $x-y+1=0$:
\[
d=\frac{|2-1+1|}{\sqrt{2}}=\sqrt{2}.
\]
Chord length formula $2\sqrt{r^{2}-d^{2}}=2\sqrt{10-2}=4\sqrt{2}$. \checkmark
\end{corrige}

\begin{corrige}[Exercise 8 — Circles under constraints]
\begin{enumerate}
\item Circle passes through $(2,0)$ and $(4,0)$; centre lies on $2x+y=7$.
The centre lies on the perpendicular bisector of the segment joining $(2,0)$ and $(4,0)$, which is $x=3$.
Substitute $x=3$ into $2x+y=7$: $6+y=7 \Rightarrow y=1$. Centre $(3,1)$.
Radius$^{2}=(3-2)^{2}+(1-0)^{2}=2$.
Equation: $(x-3)^{2}+(y-1)^{2}=2$, i.e. $x^{2}+y^{2}-6x-2y+8=0$.
\item Circle touches both coordinate axes and passes through $(2,1)$.
Centre must be $(r,r)$ (first quadrant) with radius $r$.
\[
(2-r)^{2}+(1-r)^{2}=r^{2} \;\Longrightarrow\; 4-4r+r^{2}+1-2r+r^{2}=r^{2}
\]
\[
\Longrightarrow\; r^{2}-6r+5=0 \;\Longrightarrow\; (r-1)(r-5)=0 \;\Longrightarrow\; r=1 \text{ or } r=5.
\]
Equations: $(x-1)^{2}+(y-1)^{2}=1$ and $(x-5)^{2}+(y-5)^{2}=25$.
\item Circle touches the $x$-axis at $(2,0)$ and passes through $(1,2)$.
By symmetry, centre is $(2,k)$ with radius $|k|$. Since $(1,2)$ is above the $x$-axis, $k>0$, radius $=k$.
\[
(1-2)^{2}+(2-k)^{2}=k^{2} \;\Longrightarrow\; 1+4-4k+k^{2}=k^{2} \;\Longrightarrow\; 5=4k \;\Longrightarrow\; k=\frac{5}{4}.
\]
Equation: $(x-2)^{2}+\left(y-\frac{5}{4}\right)^{2}=\frac{25}{16}$.
\end{enumerate}
\end{corrige}

\begin{corrige}[Exercise 9 — Circle through three points and tangents from an external point]
\begin{enumerate}
\item Let the circle be $x^{2}+y^{2}+2gx+2fy+c=0$. Substituting $A(1,1)$, $B(2,-1)$, $C(3,2)$:
\[
\begin{cases}
2+2g+2f+c=0 & (1)\\
5+4g-2f+c=0 & (2)\\
13+6g+4f+c=0 & (3)
\end{cases}
\]
$(2)-(1)$: $3+2g-4f=0 \;\Longrightarrow\; 2g-4f=-3$. \quad $(3)-(2)$: $8+2g+6f=0 \;\Longrightarrow\; 2g+6f=-8$.
Subtracting: $10f=-5 \;\Longrightarrow\; f=-\frac{1}{2}$.
Then $2g-4(-\frac{1}{2})=-3 \;\Longrightarrow\; 2g+2=-3 \;\Longrightarrow\; g=-\frac{5}{2}$.
From $(1)$: $c=-2-2g-2f = -2+5+1=4$.
\[
\boxed{x^{2}+y^{2}-5x-y+4=0},\qquad \text{centre } \left(\tfrac{5}{2},\tfrac{1}{2}\right),\quad r=\sqrt{\tfrac{25}{4}+\tfrac{1}{4}-4}=\tfrac{\sqrt{10}}{2}.
\]
\item Tangents from $P(7,1)$ to the circle $x^{2}+y^{2}=25$.
Line through $P$: $y-1=m(x-7)$, i.e. $mx-y+1-7m=0$.
Tangency condition: distance from centre $(0,0)$ to line equals radius $5$:
\[
\frac{|1-7m|}{\sqrt{m^{2}+1}}=5 \;\Longrightarrow\; (1-7m)^{2}=25(m^{2}+1)
\]
\[
\Longrightarrow\; 1-14m+49m^{2}=25m^{2}+25 \;\Longrightarrow\; 24m^{2}-14m-24=0
\]
\[
\Longrightarrow\; 12m^{2}-7m-12=0.
\]
Discriminant $\Delta = 49+576=625$, $\sqrt{\Delta}=25$.
\[
m=\frac{7\pm 25}{24} \;\Longrightarrow\; m=\frac{4}{3} \text{ or } m=-\frac{3}{4}.
\]
Tangents: $4x-3y-25=0$ and $3x+4y-25=0$.
Length of each tangent: $\sqrt{7^{2}+1^{2}-25}=\sqrt{25}=5$.
\item Parametric equations: $x=3+5\cos\theta$, $y=-1+5\sin\theta$.
Cartesian equation: $(x-3)^{2}+(y+1)^{2}=25$.
At $\theta=\frac{\pi}{3}$: $\cos\frac{\pi}{3}=\frac{1}{2}$, $\sin\frac{\pi}{3}=\frac{\sqrt{3}}{2}$.
Point: $\left(3+\frac{5}{2},\,-1+\frac{5\sqrt{3}}{2}\right)=\left(\frac{11}{2},\,-1+\frac{5\sqrt{3}}{2}\right)$.
Tangent at $(x_{1},y_{1})$: $(x_{1}-3)(x-3)+(y_{1}+1)(y+1)=25$.
\[
\frac{5}{2}(x-3)+\frac{5\sqrt{3}}{2}(y+1)=25 \;\Longrightarrow\; (x-3)+\sqrt{3}(y+1)=10
\]
\[
\Longrightarrow\; x+\sqrt{3}\,y = 13-\sqrt{3}.
\]
\end{enumerate}
\end{corrige}

\begin{corrige}[Exercise 10 — Families of circles, radical axis and touching circles]
\begin{enumerate}
\item Circles: $S_{1}: x^{2}+y^{2}-4=0$, $S_{2}: x^{2}+y^{2}-2x-4y=0$.
Family through their intersections: $S_{1}+\lambda S_{2}=0$:
\[
(1+\lambda)x^{2}+(1+\lambda)y^{2}-2\lambda x-4\lambda y-4=0.
\]
Passes through $(2,3)$:
\[
(1+\lambda)(4+9)-2\lambda(2)-4\lambda(3)-4=0 \;\Longrightarrow\; 13(1+\lambda)-4\lambda-12\lambda-4=0
\]
\[
\Longrightarrow\; 13+13\lambda-16\lambda-4=0 \;\Longrightarrow\; 9-3\lambda=0 \;\Longrightarrow\; \lambda=3.
\]
Equation: $4x^{2}+4y^{2}-6x-12y-4=0$, i.e. $2x^{2}+2y^{2}-3x-6y-2=0$.
\item Circle 1: $x^{2}+y^{2}-4x-6y+9=0$. Centre $C_{1}(2,3)$, $r_{1}=\sqrt{4+9-9}=2$.
Circle 2: $x^{2}+y^{2}+2x+2y-7=0$. Centre $C_{2}(-1,-1)$, $r_{2}=\sqrt{1+1+7}=3$.
Distance $C_{1}C_{2}=\sqrt{(2+1)^{2}+(3+1)^{2}}=\sqrt{9+16}=5 = r_{1}+r_{2}$.
The circles touch externally. Point of contact divides $C_{1}C_{2}$ internally in the ratio $r_{1}:r_{2}=2:3$:
\[
P = \left( \frac{3\cdot2 + 2\cdot(-1)}{5},\; \frac{3\cdot3 + 2\cdot(-1)}{5} \right) = \left( \frac{4}{5},\; \frac{7}{5} \right).
\]
(Check: $\left(\frac{4}{5}+1\right)^{2}+\left(\frac{7}{5}+1\right)^{2}=\frac{81+144}{25}=9=r_{2}^{2}$. \checkmark)
\item Radical axis of circles with centres $(3,2)$ and $(-1,4)$:
$S_{1}-S_{2}=0$ gives $-8x+4y+1=0$, i.e. $8x-4y-1=0$. Gradient $=2$.
Line of centres gradient: $\dfrac{4-2}{-1-3}=-\dfrac{1}{2}$.
Product $2 \times \left(-\dfrac{1}{2}\right) = -1$, so the radical axis is perpendicular to the line of centres.
\end{enumerate}
\end{corrige}

\begin{corrige}[Exercise 11 — GCE past-paper style: locus and tangents of variable slope]
\begin{enumerate}
\item Let $P(x,y)$. Condition $PA=2PB$ with $A(-2,0)$, $B(4,0)$:
\[
PA^{2}=4\,PB^{2} \;\Longrightarrow\; (x+2)^{2}+y^{2}=4\left[(x-4)^{2}+y^{2}\right]
\]
\[
x^{2}+4x+4+y^{2}=4x^{2}-32x+64+4y^{2}
\]
\[
0=3x^{2}+3y^{2}-36x+60 \;\Longrightarrow\; x^{2}+y^{2}-12x+20=0.
\]
This is a circle (coefficients of $x^{2}$ and $y^{2}$ equal, no $xy$ term, and $g^{2}+f^{2}-c=36-20=16>0$).
Centre $(6,0)$, radius $4$.
\item Circle: $x^{2}+y^{2}-10x+16=0$. Centre $(5,0)$, radius $r=\sqrt{25-16}=3$.
\begin{enumerate}
\item Tangents from origin $O(0,0)$: line $y=mx$, i.e. $mx-y=0$.
Tangency condition: distance from centre $(5,0)$ to line equals $3$:
\[
\frac{|5m|}{\sqrt{m^{2}+1}}=3 \;\Longrightarrow\; 25m^{2}=9(m^{2}+1) \;\Longrightarrow\; 16m^{2}=9 \;\Longrightarrow\; m=\pm\frac{3}{4}.
\]
\item Point of contact is the foot of the perpendicular from centre $(5,0)$ to the line $mx-y=0$.
Foot of perpendicular from $(x_{0},y_{0})$ to $ax+by+c=0$:
\[
\left( x_{0}-\frac{a(ax_{0}+by_{0}+c)}{a^{2}+b^{2}},\; y_{0}-\frac{b(ax_{0}+by_{0}+c)}{a^{2}+b^{2}} \right).
\]
Here $a=m$, $b=-1$, $c=0$, $(x_{0},y_{0})=(5,0)$:
\[
ax_{0}+by_{0}+c = 5m.
\]
Foot:
\[
\left( 5-\frac{m(5m)}{m^{2}+1},\; 0-\frac{(-1)(5m)}{m^{2}+1} \right) = \left( \frac{5}{m^{2}+1},\; \frac{5m}{m^{2}+1} \right).
\]
With $m^{2}=\frac{9}{16}$, $m^{2}+1=\frac{25}{16}$:
\[
x = \frac{5}{25/16} = \frac{16}{5},\qquad y = \frac{5m}{25/16} = \frac{16m}{5} = \pm\frac{12}{5}.
\]
Points of contact: $\left(\dfrac{16}{5},\dfrac{12}{5}\right)$ and $\left(\dfrac{16}{5},-\dfrac{12}{5}\right)$.
\end{enumerate}
\end{enumerate}
\end{corrige}

\newpage

\coursec{Summary sheet}

\begin{popfigure}
\begin{tikzpicture}[
  mindmap,
  grow cyclic,
  every node/.style={concept, text=white, font=\small},
  root concept/.append style={concept color=popink, font=\bfseries\small},
  level 1/.append style={level distance=3.4cm, sibling angle=60},
  level 1 concept/.append style={font=\footnotesize},
]
\node[root concept]{Circle Geometry}
  child[concept color=popblue]{ node {Equations: centre--radius, diameter, 3 points} }
  child[concept color=popgreen]{ node {General form $x^2+y^2+2gx+2fy+c=0$} }
  child[concept color=poporange]{ node {Tangents: equations, conditions, length} }
  child[concept color=poppurple]{ node {Parametric form $(a+r\cos\theta,\,b+r\sin\theta)$} }
  child[concept color=poppink]{ node {Intersections: line--circle, circle--circle, touching} }
  child[concept color=popgold]{ node {Families $S_1+\lambda S_2=0$, $S+\lambda L=0$} }
  child[concept color=popdark]{ node {Radical axis, locus, constrained circles} };
\end{tikzpicture}
\legende{Mind map of the chapter Circle Geometry.}
\end{popfigure}

\begin{bilan}

\textbf{1. Key definitions}
\begin{itemize}
\item A \cle{circle} is the locus of a point $P(x,y)$ whose distance from a fixed point $C$ (the \cle{centre}) is constant (the \cle{radius} $r$).
\item A \cle{tangent} is a line meeting the circle in exactly one point; it is perpendicular to the radius at the point of contact.
\item The \cle{radical axis} of two circles is the locus of points with equal tangent-lengths (equal powers) to both circles.
\end{itemize}

\textbf{2. Equations of a circle (know them all)}
\begin{itemize}
\item Centre $O(0,0)$, radius $r$: \; $x^2+y^2=r^2$.
\item Centre $(a,b)$, radius $r$: \; $(x-a)^2+(y-b)^2=r^2$.
\item General form: $x^2+y^2+2gx+2fy+c=0$, centre $(-g,-f)$, radius $r=\sqrt{g^2+f^2-c}$ (real circle iff $g^2+f^2-c>0$).
\item Diameter form, ends $A(x_1,y_1)$, $B(x_2,y_2)$: $(x-x_1)(x-x_2)+(y-y_1)(y-y_2)=0$ (angle in a semicircle is $90^\circ$).
\item Parametric form: $x=a+r\cos\theta$, $y=b+r\sin\theta$; convert to Cartesian using $\cos^2\theta+\sin^2\theta=1$.
\item Through three points: substitute into the general form and solve the $3\times 3$ system for $g,f,c$.
\end{itemize}

\textbf{3. Tangents}
\begin{itemize}
\item Tangent at $(x_1,y_1)$ on $x^2+y^2=r^2$: \; $xx_1+yy_1=r^2$.
\item Tangent at $(x_1,y_1)$ on the general circle: replace $x^2\to xx_1$, $y^2\to yy_1$, $2x\to x+x_1$, $2y\to y+y_1$.
\item Line $y=mx+c$ is tangent to $x^2+y^2=r^2$ iff $c^2=r^2(1+m^2)$; in general, tangency iff the \cle{perpendicular distance} from the centre to the line equals the radius, or the resulting quadratic has discriminant zero.
\item Length of tangent from external point $P(x_1,y_1)$: $L=\sqrt{x_1^2+y_1^2+2gx_1+2fy_1+c}=\sqrt{S_1}$.
\end{itemize}

\textbf{4. Intersections and families}
\begin{itemize}
\item Line--circle: substitute the line into the circle; discriminant $\Delta>0$ (two points), $\Delta=0$ (tangent), $\Delta<0$ (no meeting).
\item Two circles intersect where $S_1-S_2=0$ meets either circle; $S_1-S_2=0$ is the \cle{radical axis} (common chord when they intersect).
\item Touching conditions: externally $C_1C_2=r_1+r_2$; internally $C_1C_2=|r_1-r_2|$.
\item Family through the intersections of two circles: $S_1+\lambda S_2=0$; through intersections of line $L=0$ and circle $S=0$: $S+\lambda L=0$. Find $\lambda$ from the extra condition (a point, tangency, etc.).
\end{itemize}

\textbf{5. Constrained circles}
\begin{itemize}
\item Touching both axes: centre $(\pm r,\pm r)$, equation $(x\pm r)^2+(y\pm r)^2=r^2$.
\item Centre on a given line / circle touching a given line: use the centre-line relation or distance-from-centre-to-line $=r$ to fix the unknowns.
\end{itemize}

\textbf{6. Methods to master}
\begin{itemize}
\item Completing the square to obtain centre and radius from the general equation.
\item Solving simultaneous equations (line with circle; two circles via the radical axis).
\item Building a locus: let the moving point be $(h,k)$, translate the condition into an equation, then replace $(h,k)$ by $(x,y)$.
\end{itemize}

\textbf{7. Common mistakes}
\begin{itemize}
\item Forgetting the factor $2$: centre is $(-g,-f)$, not $(-2g,-2f)$.
\item Wrong radius: $r=\sqrt{g^2+f^2-c}$, not $\sqrt{g^2+f^2+c}$.
\item Using $S_1$ (power) instead of $\sqrt{S_1}$ for the tangent length.
\item Sign errors in $S_1+\lambda S_2=0$ when substituting the given point.
\item Forgetting the second solution (two tangents, two touching circles, $\pm$ cases).
\end{itemize}

\textbf{8. GCE tips}
\begin{itemize}
\item State the centre and radius immediately after completing the square — examiners award marks for this.
\item For tangency questions, the discriminant method and the perpendicular-distance method both earn full credit; choose the shorter.
\item In family-of-circles questions, write $S_1+\lambda S_2=0$ explicitly before substituting the point.
\item Always check whether your final circle is real ($r^2>0$) and give answers in exact surd form unless told otherwise.
\end{itemize}
\end{bilan}

\findecours

\end{document}
